\documentclass[AMA,twocolumn]{USG}

\usepackage{bm}
\usepackage{booktabs}
\usepackage{algorithm}
\usepackage{algpseudocode}
\usepackage{rotating}
\usepackage{url}
\usepackage{placeins}
\def\UrlBreaks{\do\/\do-\do.\do:\do_\do\~}

\graphicspath{{./images/}{./figures/review/}}
\newif\ifggkmsubmission
\ifdefined\GGKMSubmissionLayout
  \ggkmsubmissiontrue
\else
  \ggkmsubmissionfalse
\fi
\ifggkmsubmission
  \newcounter{submissionfigure}
\fi

\newcommand{\GGKMTable}[1]{%
  \ifggkmsubmission\else\input{#1}\fi}
\newcommand{\GGKMWideFigure}[4]{%
  \ifggkmsubmission
    \begin{figure*}[tbp]
      \noindent\textit{[Figure~\ref{#3} near here]}\par
      \caption{#2}\label{#3}
    \end{figure*}%
  \else
    \begin{figure*}[tbp]
      \centerline{\includegraphics[width=\linewidth]{#1}}%
      \caption{#2}\label{#3}
    \end{figure*}%
  \fi}

\articletype{RESEARCH ARTICLE}
\journal{Statistics in Medicine}
\volume{}
\startpage{1}
\articledoi{}

\begin{document}
\title{A generalized gamma kernel-machine promotion-time cure model}
\titlemark{Kernel-machine promotion-time cure model}
\shrttitle{Kernel-machine promotion-time cure model}
\author[1]{Gabriel de Jesus Pereira}[https://orcid.org/0009-0003-3254-4854]
\author[1]{Vera L. D. Tomazella}[https://orcid.org/0000-0002-6780-2089]
\author[2]{Pedro R. D. Marinho}[https://orcid.org/0000-0003-1591-8300]
\author[2]{Marcelo Rodrigo Portela Ferreira}[https://orcid.org/0000-0003-0242-0255]
\author[2]{Agatha S. Rodrigues}[https://orcid.org/0000-0003-1112-8264]
\authormark{PEREIRA \textsc{et al.}}
\address[1]{\orgdiv{Department of Statistics, }\orgname{Universidade Federal de S\~ao Carlos, }\orgaddress{\state{S\~ao Carlos, }\country{Brazil}}}
\address[2]{\orgdiv{Department of Statistics, }\orgname{Universidade Federal da Para\'iba, }\orgaddress{\state{Jo\~ao Pessoa, }\country{Brazil}}}
\corres{Vera L. D. Tomazella, Department of Statistics, Universidade Federal de S\~ao Carlos, S\~ao Carlos, Brazil. \textbf{AUTHOR CONFIRMATION REQUIRED:} corresponding-author email.}
\keywords{cure model | generalized gamma | kernel machine | promotion time | survival analysis}
  \abstract[ABSTRACT]{Cure models make the long-term event-free fraction an explicit quantity, but
promotion-time cure models usually take the mean number of latent causes to be
log-linear in the covariates, which cannot represent nonlinear, threshold, or
interaction effects and can bias both the covariate effects and the estimated
cured fraction. We propose a generalized gamma kernel-machine promotion-time
cure model (GG-KM) that places the log-intensity of the latent-cause process in
a reproducing-kernel Hilbert space, retains a flexible parametric generalized
gamma latency distribution, and allows the latent number of causes to be Poisson
or negative binomial, thereby accommodating both equidispersion and
overdispersion. Estimation uses penalized expectation-maximization for the
Poisson model and expectation/conditional maximization either for the
negative-binomial model, followed by bounded observed-likelihood refinement and
explicit convergence diagnostics. We evaluate the framework through simulation
experiments with linear, quadratic, and periodic covariate effects and a
four-covariate mechanism design, and through applications to melanoma and
breast-cancer cohorts, using event-stratified training and test splits,
training-only tuning, and inverse-probability-of-censoring weighted integrated
Brier scores together with calibration and cure-fraction accuracy. No single
predictor is uniformly best: a correctly specified simpler predictor performs
best when the covariate structure matches its form, whereas GG-KM is competitive
across mechanisms, is the most accurate predictor for the radial four-covariate
mechanism, and improves on the log-linear and additive-spline predictors for the
interaction mechanism, although the correctly aligned quadratic model remains
best in that setting. In the breast-cancer cohort, the spline models attained
the lowest mean integrated Brier score, whereas the mixture-cure model attained
the highest mean cure AUC. These findings support GG-KM as a flexible and robust
alternative when the functional form of the latent-cause mechanism is unknown,
and the melanoma and breast-cancer applications illustrate its use for
reproducible out-of-sample prediction.}

%{Cure models estimate the fraction of individuals who never experience the event of interest, but linear promotion-time predictors can misrepresent nonlinear covariate effects on the latent-cause mechanism. The generalized gamma kernel-machine promotion-time cure model places the log latent-cause intensity in a reproducing-kernel Hilbert space with an explicit unpenalized intercept, retains a common Stacy generalized gamma latency, and accommodates Poisson and negative-binomial latent counts. Estimation uses safeguarded penalized expectation-maximization or expectation/conditional maximization either, followed by bounded observed-likelihood refinement and explicit convergence diagnostics. Prediction for unseen covariates uses training-only preprocessing and the fitted kernel expansion. Evaluation uses common event-stratified 70/30 splits, inner tuning confined to training data, inverse-probability-of-censoring weighted Brier scores, paired uncertainty, retained failures, and versioned simulation and melanoma outputs. The comparison includes log-linear, additive spline, quadratic, and radial-basis kernel predictors under both count laws, together with separate primary and mechanism designs. Quantitative conclusions are restricted to reproduced outputs from the declared implementation and do not assert universal superiority, causal effects, or identification of individual cure status.}

\maketitle

\section{Introduction}\label{sec:introduction}

Standard survival models assign vanishing probability to remaining event
free as follow-up grows arbitrarily long. That limit is inappropriate
when a clinically meaningful subgroup will never experience the event of
interest, a situation indicated by a sustained plateau in the survival
curve under mature follow-up. Cure models describe this setting by
separating the long-term event-free probability from the timing of
events among individuals who remain susceptible, which makes the cured
proportion an interpretable quantity in oncology and other long-term-survivor
studies \cite{berkson1952,tsodikov2003}.

{
Accommodating this plateau requires a model that allows a positive
limiting survival probability, and three main constructions are available
in the literature for this purpose, differing in the mechanism by which
they generate the cure mass. The first, the mixture cure model,
introduces a latent binary indicator that separates long-term event-free
individuals from those who remain susceptible, so that the cured
proportion and the survival of the susceptible subpopulation are
specified as separate components, typically a regression model for
incidence together with a proper distribution for the latency
\cite{boag1949,berkson1952,farewell1982,sy2000}. The second, the
defective-distribution model, does not represent the cured fraction
through a separate component; instead, the cured fraction arises
intrinsically from the parameter space of a single parametric survival
distribution. For most parameter values the distribution is proper and
the survival function decays to zero, as in conventional survival models;
over a distinct region of the same parameter space the distribution
becomes improper and the survival function instead converges to a
strictly positive limit, which is the cured fraction. A single family
then governs both the failure times of the susceptible individuals and
the size of the cured group, with no separate incidence model and with
covariates entering through its parameters
\cite{balka2009,rocha2017,scudilio2019}. The third, based on latent
competing risks, generates the cure mass from an unobserved count of
competing causes, so that the cured fraction is The probability that this
count is zero \cite{yakovlev1996,chen1999,tsodikov2003,rodrigues2009unif}.
} %This construction can be re-expressed as a mixture whose cured proportion equals that zero-count probability and whose susceptible latency is induced by the count law, so it is a structured special case of the mixture model in which the incidence and latency components are linked through shared parameters rather than specified separately \cite{rodrigues2009}.
{
This mechanism admits a two-stage description. An initiation stage sets
the unobserved number $N$ of latent causes carried by an individual, and
a maturation stage assigns each cause an activation time; the earliest
activation produces the event, so $N=0$ describes an individual who never
fails and identifies the cure mass \cite{rodrigues2025}. Encoding the
count law through its probability-generating function links the latency
distribution to both the population survival and the cure fraction, and
it is this representation on which the model developed here is built.
}

{
In the promotion-time construction the mean number of latent causes is
almost always taken to be log-linear in the covariates
\cite{chen1999,rodrigues2009}. This specification is convenient and
keeps the cure fraction in closed form, but it forces every covariate
to act monotonically on the latent-cause intensity, so threshold effects,
interactions, and periodic or otherwise nonlinear associations cannot be
represented. Because the intensity governs the cure fraction, a
misspecified incidence model biases both the estimated covariate effects
and the cured proportion \cite{amico2018}. The activation-time component
carries another issue: fixing the latency to a restrictive family can be
equally misspecified, and in a cure model this error propagates to the
covariate effects and to the cure fraction. A flexible parametric family
mitigates this while preserving interpretability and a tractable
likelihood, and the generalized gamma is a natural choice, since it
nests the exponential, Weibull, gamma, and log-normal distributions and
spans a wide range of hazard shapes
\cite{stacy1962,prentice1974,cox2007}. Together, these concerns motivate
treating both components flexibly, through a flexible latent-intensity
predictor for the incidence and a sufficiently general parametric latency
for the activation times.
}

{
Flexible predictors have been proposed in cure modelling: smoothers,
support-vector machines, neural networks, and deep architectures have all
been used to relax the usual linear specification. In the promotion-time
setting, Pal and Aselisewine\cite{pal2023} introduced a support-vector
formulation for the incidence component, and Xie and Yu\cite{xie2021}
used a neural network for an unstructured promotion-time component,
while deep promotion-time models and structured nonlinear cure models
further broaden the available high-dimensional architectures\cite{medina2024,feng2026}.
These methods make the covariate effect on the latent-cause intensity
nonlinear, but they differ substantially in how they treat that
intensity, the latency distribution, and the count structure. To our
knowledge, no existing promotion-time approach combines all three within
a single likelihood-based estimation and prediction framework: a
flexible nonlinear predictor for the latent-cause intensity, which shapes
both the cured fraction and the timing of events; a transparent and
sufficiently general parametric latency; and an explicit count-family
treatment of overdispersion.
}

To address this gap, we develop a generalized gamma kernel-machine
promotion-time cure model (GG-KM) that places the log-intensity of
the latent-cause process in a reproducing-kernel Hilbert space (RKHS)
\cite{kimeldorf1971}, while retaining a flexible parametric
generalized gamma latency distribution \cite{stacy1962} and allowing
either Poisson or negative-binomial latent counts. This formulation
separates two sources of modelling flexibility: the kernel predictor
captures potentially complex nonlinear and interaction effects of
covariates on the latent-cause mechanism, whereas the generalized
gamma distribution provides a flexible parametric representation of
the activation-time process.


The remainder of this article is organized as follows. Section~\ref{sec:methods}
presents the proposed generalized gamma kernel-machine promotion-time cure
model, covering the general promotion-time construction, the generalized gamma
latency distribution, the kernel-based latent-cause predictor, and the penalized
EM and ECME estimation procedures with their observed-likelihood refinement and
convergence assessment. Section~\ref{sec:simulation} describes the simulation
study and evaluates the predictive performance of the proposed model against
log-linear, additive-spline, and quadratic alternatives under different
covariate structures and latent-count distributions. Section~\ref{sec:melanoma}
applies the model to the melanoma data and reports its predictive behaviour.
Section~\ref{sec:discussion} discusses the main findings, limitations, and
directions for future research.

\section{Methods}\label{sec:methods}

This section first recalls, as background, the general promotion-time cure construction that the model builds on, and then develops the proposed GG-KM model: a generalized gamma latency, a regularized reproducing-kernel Hilbert space predictor for the latent-cause intensity, and Poisson or negative-binomial latent counts. The remaining subsections give the penalized estimation algorithms and the prediction and evaluation protocol. Observed likelihoods and conditional objective functions are stated in the main text, while lengthy derivations and analytical gradients are deferred to the Supporting Information.

%{ Figure~\ref{fig:ggkm-workflow} provides an overview of the GG-KM model and its estimation and evaluation framework. Panel A illustrates the kernel-based latent-cause intensity $\theta(\bm{x})=\exp\{f(\bm{x})\}$ and its link to the Poisson and negative-binomial latent-count distributions. Panel B depicts the promotion-time mechanism, in which individuals with $N_i=0$ are cured, whereas for $N_i\geq1$ the event time is the minimum of the latent activation times. Panel C shows the resulting population survival curves and the cure plateau $p_c(\bm{x})$. Panel D summarizes the estimation and prediction pipeline, including the event-stratified data split, training-only preprocessing, penalized EM or ECME estimation, out-of-sample prediction, and IPCW integrated Brier evaluation.}

 %Figure~\ref{fig:ggkm-workflow} lays out the GG-KM model and its estimation and evaluation path. Panel A shows the latent-cause intensity $\theta(\bm{x})=\exp\{f(\bm{x})\}$ built by a radial-basis kernel expansion over $(x_1,x_2)\in[-1,1]^2$, with the six kernel centres $\bm{x}_j$ marked and the count law named beneath the surface. Panel B plots four parallel time axes for two cured individuals ($N_i=0$) and two susceptible ones ($N_i\ge1$). Its key maps every glyph to a quantity: an open circle is an activation time $Z_{ij}$ from the generalized gamma latency, a filled circle is the event time $T_i=\min_jZ_{ij}$, and a vertical bar is the censoring time $C_i$; the resulting pair $(Y_i,\delta_i)$ is printed at the right of each axis. Panel C draws the matched Poisson and negative-binomial population survival curves, the cure plateau $p_c(\bm{x})$, and the shaded interval at or below the scoring horizon $\tau$. Panel D traces the leakage-controlled pipeline: event-stratified 70/30 split, Gram matrix construction, penalized EM or ECME with bounded refinement, kernel-coefficient prediction, and IPCW Brier evaluation. The dashed arc in Panel D is the only feedback path and marks the single iterative update that revises the kernel coefficients, the latency parameters, and the negative-binomial dispersion when applicable.

%\GGKMWideFigure{ggkm-method-diagram.pdf}{Structure of the GG-KM model and of its estimation and evaluation. Panel A shows the generative mechanism: a kernel predictor $f(\bm{x})=\beta_0+\sum_j\alpha_jK(\bm{x},\bm{x}_j)$ over the marked kernel centres $\bm{x}_j$ determines the positive latent-cause intensity $\theta(\bm{x})=\exp\{f(\bm{x})\}$, which indexes the Poisson or negative-binomial count $N_i$ named beneath the surface. Panel B shows the observation mechanism. Given $N_i\ge1$ the promotion times $Z_{ij}$ are independent generalized gamma variables and $T_i=\min_jZ_{ij}$, whereas $N_i=0$ gives $T_i=\infty$, so only $Y_i=\min(T_i,C_i)$, $\delta_i=I(T_i\leq C_i)$ and $\bm{x}_i$ are observed; the key maps the open circle, the filled circle and the vertical bar to $Z_{ij}$, $T_i$ and $C_i$, and the realized pair $(Y_i,\delta_i)$ is printed at the right of each time axis. Panel C shows the population survival curves, where P denotes Poisson and NB negative-binomial counts. They use the simulation latency $(a,d,p)=(1.6,1.3,1.1)$, a matched cure fraction of $0.45$ and dispersion $\phi=1$; they are illustrative rather than fitted, and line style with markers distinguishes them in addition to colour. Shading marks times at or below the frozen horizon $\tau$ within which predictions are scored, so the approach of $S_{\mathrm{pop}}(t\mid\bm{x})$ to $p_c(\bm{x})$ is extrapolated rather than evaluated. Panel D shows the estimation and evaluation path: one event-stratified 70/30 split with training-only preprocessing, the validated Gram matrix, penalized EM or ECME with an ascent safeguard, out-of-sample prediction, and IPCW integrated Brier scoring with paired differences. The dashed arc in Panel D is the only feedback path and marks the iterative update of $(\beta_0,\bm{\alpha})$, $(a,d,p)$ and $\phi$.}{fig:ggkm-workflow}{Four-panel schematic. Panel A: a kernel predictor with marked centres gives the latent-cause intensity that indexes a Poisson or negative-binomial latent count. Panel B: four time axes for two susceptible and two cured individuals, with a key mapping open circles to activation times, filled circles to event times and vertical bars to censoring times, and the observed time and event indicator printed at the right of each axis. Panel C: population survival curves for both count laws approach a common cure plateau beyond the shaded scored horizon. Panel D: event-stratified split with training-only preprocessing, validated Gram matrix, penalized EM or ECME with ascent safeguard, out-of-sample prediction, and IPCW integrated Brier evaluation, with one dashed feedback path.}

\subsection{General promotion-time cure construction}\label{sec:promotion-time}

{
	%Let $\mathcal{D}_n=\{(Y_i,\delta_i,\bm{x}_i):i=1,\dots,n\}$ denote an independent sample of size $n$, where $T_i$ is the event time, $C_i$ the right-censoring time, $Y_i=\min(T_i,C_i)$, $\delta_i=I(T_i\leq C_i)$, and $\bm{x}_i$ a covariate vector. Censoring is assumed conditionally independent, $T_i\perp C_i\mid\bm{x}_i$. The latent variable $N_i$ counts the potential causes of the event, and $N_i=0$ identifies a cured individual.

Let $T_i$ denote the event time and $C_i$ the right-censoring time. The observed time is the random variable $Y_i=\min(T_i,C_i)$, whose realized value is written $t_i$, and $\delta_i=I(T_i\leq C_i)$ is the event indicator, with $\bm{x}_i$ a covariate vector, for $i=1,\ldots,n$. The observed data are therefore $\mathcal{D}_{\text{obs}}=\{(t_i,\delta_i,\bm{x}_i)\}^{n}_{i=1}$. Censoring is assumed conditionally independent, $T_i\perp C_i\mid\bm{x}_i$. The latent variable $N_i$ counts the potential causes of the event, and $N_i=0$ identifies a cured individual.

Conditional on the covariates, $N_i\mid\bm{x}_i$ has probability mass function $\pi_k(\bm{x}_i)=P(N_i=k\mid\bm{x}_i)$, $k=0,1,2,\dots$, and probability-generating function $A_i(u)=\mathbb{E}\left(u^{N_i}\mid\bm{x}_i\right)=\sum_{k\ge0}\pi_k(\bm{x}_i)u^k$. Given $N_i=k\ge1$, the activation times $Z_{i1},\dots,Z_{ik}$ are independent and identically distributed with a proper distribution function $F_0$, survival function $S_0=1-F_0$, and density $f_0$; conditional on $\bm{x}_i$ they are independent of $N_i$, and their distribution does not depend on $N_i$. In the formulation considered here the covariates act on the latent-count distribution but not on the latency $F_0$. The event time is the earliest activation, $T_i=\min_{1\le j\le N_i}Z_{ij}$ when $N_i\ge1$, whereas $T_i=\infty$ when $N_i=0$.

Conditional on $N_i=k$, the event exceeds $t$ only if all $k$ activation times do, so $P(T_i>t\mid N_i=k,\bm{x}_i)=S_0(t)^k$. Averaging over the latent-count distribution gives the population survival function and, as $t\to\infty$ with $S_0(t)\to0$, the cure fraction,
\begin{equation}
\begin{aligned}
S_{\mathrm{pop}}(t\mid\bm{x}_i)
&=P(T_i>t\mid\bm{x}_i)
=\sum_{k\ge0}\pi_k(\bm{x}_i)S_0(t)^k
=A_i\{S_0(t)\},\\
p_c(\bm{x}_i)
&=\lim_{t\to\infty}S_{\mathrm{pop}}(t\mid\bm{x}_i)
=A_i(0)=P(N_i=0\mid\bm{x}_i).
\end{aligned}
\label{eq:promotion-population}
\end{equation}
The population survival is the generating function evaluated at the latency survival, and the cure fraction is its value at zero, equal to the probability of no latent cause. Differentiating with respect to time, and using $S_0'=-f_0$, gives the population density and hazard,
\begin{equation}
\begin{aligned}
f_{\mathrm{pop}}(t\mid\bm{x}_i)
&=f_0(t)\,A_i'\{S_0(t)\},\\
h_{\mathrm{pop}}(t\mid\bm{x}_i)
&=\frac{f_{\mathrm{pop}}(t\mid\bm{x}_i)}
        {S_{\mathrm{pop}}(t\mid\bm{x}_i)},
\end{aligned}
\label{eq:promotion-density-hazard}
\end{equation}
where $A_i'$ denotes differentiation with respect to its argument. When $p_c(\bm{x}_i)>0$, $f_{\mathrm{pop}}$ is an improper subdensity that integrates to $1-p_c(\bm{x}_i)$ rather than to one. The construction satisfies the required boundary behaviour: $S_{\mathrm{pop}}(0\mid\bm{x})=1$, $S_{\mathrm{pop}}(t\mid\bm{x})$ is non-increasing in $t$, and $\lim_{t\to\infty}S_{\mathrm{pop}}(t\mid\bm{x})=p_c(\bm{x})$. % A degenerate cure fraction $p_c(\bm{x})=1$ corresponds to the latent-count mass concentrating entirely on $N_i=0$, that is $\pi_0(\bm{x})=1$, which is a limiting boundary case rather than a value attained at interior count parameters.

Equations~\eqref{eq:promotion-population} and \eqref{eq:promotion-density-hazard} give the standard promotion-time construction for a generic activation-time distribution. The proposed GG-KM model, developed in the following subsection, specializes the latency $F_0$ to the Stacy generalized gamma distribution and models the latent-count mean $\theta_i=\theta(\bm{x}_i)$ through a reproducing-kernel Hilbert space predictor, with the count taken to be Poisson or negative binomial.}

\subsection{The proposed GG-KM model}\label{sec:model-specification}

%{The proposed GG-KM model specializes the general construction of Section~\ref{sec:promotion-time} in three linked components: the activation-time distribution, the latent-cause intensity, and the latent-count law. The activation times follow a Stacy generalized gamma latency with parameters $(a,d,p)$ shared across individuals; the latent-cause intensity is $\theta(\bm{x})=\exp\{f(\bm{x})\}$, where $f$ is a regularized reproducing-kernel Hilbert space predictor; and the latent count $N_i\mid\bm{x}_i$ is Poisson or negative binomial with mean $\theta_i=\theta(\bm{x}_i)$. Covariates act on the model only through the intensity $\theta(\bm{x})$, and hence only through $f$. The parameters satisfy $a,d,p>0$ and $\theta_i>0$, and the negative-binomial count adds a dispersion parameter $\phi>0$. The three components are specified in turn below.}

The activation time of each latent cause follows the Stacy generalized gamma distribution with scale $a>0$ and shape parameters $d>0$ and $p>0$.\cite{stacy1962} Its density, distribution function, and survival function are
\begin{equation}
	\begin{aligned}
		f_{GG}\left(t;a,d,p\right)
		&=\frac{p t^{d-1}\exp\{-(t/a)^p\}}
		{a^d\Gamma(d/p)},\\
		F_{GG}\left(t;a,d,p\right)
		&=P\{d/p,(t/a)^p\},\\
		S_{GG}\left(t;a,d,p\right)
		&=Q\{d/p,(t/a)^p\},
	\end{aligned}
	\label{eq:gg-distribution}
\end{equation}
where $P$ and $Q$ are the regularized lower and upper incomplete gamma
functions. This parameterization contains Weibull models when $d=p$,
gamma models when $p=1$, and the exponential model when $d=p=1$. The
lognormal distribution is not a special case or limit of the Stacy
parameterization at fixed $(a,d,p)$; obtaining it requires a separate
reparameterized limit and is therefore not claimed here. The
implementation verifies numerically that the density integrates to one
and that the derivative of the distribution function equals the density,
so $F_{GG}$, $f_{GG}$, and $S_{GG}$ remain mutually consistent. Figure~\ref{fig:gg-latency-flexibility} illustrates how the verified special cases and a general parameter setting produce distinct density and hazard shapes.

\GGKMWideFigure{gg-latency-flexibility.pdf}{Generalized gamma latency flexibility under the Stacy parameterization. Panel A shows the density $f_{GG}\left(t; a,d,p\right)$ and Panel B the corresponding hazard $h_{GG}\left(t;a,d,p\right)=f_{GG}\left(t;a,d,p\right)/S_{GG}\left(t;a,d,p\right)$ for the exponential case $(a,d,p)=(1.6,1,1)$, Weibull case $(1.6,0.7,0.7)$, gamma case $(1.6,2,1)$, and the general generalized gamma setting $(1.6,1.3,1.1)$ used in the simulation study. Line type and markers distinguish the four cases in addition to colour.}{fig:gg-latency-flexibility}{Two panels compare four generalized gamma latency distributions. The left panel shows their densities over time and the right panel shows their hazards. Exponential, Weibull, gamma, and general generalized gamma curves use distinct line types and markers.}


{

	The GG-KM model relaxes the usual log-linear specification of the latent-cause
	intensity by representing
	$
	    \eta(\bm{x})
	    =
	    \log\{\theta(\bm{x})\}
	$
	through a reproducing kernel Hilbert space (RKHS). Here,
	$\theta(\bm{x})>0$ denotes the conditional mean of the latent number of causes,
	\begin{equation}
	    \theta(\bm{x})
	    =
	    \mathbb{E}(N\mid\bm{x}),
	\end{equation}
	under both the Poisson and negative-binomial parameterizations considered in
	this work. Modelling $\eta(\bm{x})$ in an RKHS allows the association between
	the covariates and the latent-cause mechanism to be nonlinear without requiring
	a prespecified finite-dimensional functional form.

	% The covariate effects enter the model through the mean parameter of the
	% latent number of competing causes. Let
	% $\mathcal{X}\subseteq\mathbb{R}^{r}$ denote the covariate space and $\bm{x},\bm{x}'\in\mathcal{X}$ two covariate vectors, and let $K:\mathcal{X}\times\mathcal{X}\to\mathbb{R}$ be a kernel measuring their similarity. The
	% nonlinear component of the predictor is modelled through a reproducing
	% kernel selected from the following family:
	% \begin{equation}
	% 	\begin{aligned}
	% 		K_{\mathrm{lin}}(\bm{x},\bm{x}')
	% 		&=\bm{x}^{\mathsf T}\bm{x}',\\
	% 		K_{\mathrm{poly}}(\bm{x},\bm{x}')
	% 		&=(\gamma\,\bm{x}^{\mathsf T}\bm{x}'+c_0)^{q},\\
	% 		K_{\mathrm{rbf}}(\bm{x},\bm{x}')
	% 		&=\exp\{-\gamma\|\bm{x}-\bm{x}'\|_2^2\},\\
	% 		K_{\mathrm{exp}}(\bm{x},\bm{x}')
	% 		&=\exp\{-\gamma\|\bm{x}-\bm{x}'\|_2\},\\
	% 		K_{\mathrm{lap}}(\bm{x},\bm{x}')
	% 		&=\exp\{-\gamma\|\bm{x}-\bm{x}'\|_1\},\\
	% 		K_{\mathrm{cau}}(\bm{x},\bm{x}')
	% 		&=\{1+\gamma\|\bm{x}-\bm{x}'\|_2^2\}^{-1},\\
	% 		K_{\mathrm{wav}}(\bm{x},\bm{x}')
	% 		&=
	% 		\prod_{j=1}^{r}
	% 		\left[
	% 		\cos\!\left(
	% 		1.75\,\frac{x_j-x'_j}{A}
	% 		\right)
	% 		\exp\!\left\{
	% 		-\frac{1}{2}
	% 		\left(
	% 		\frac{x_j-x'_j}{A}
	% 		\right)^2
	% 		\right\}
	% 		\right].
	% 	\end{aligned}
	% 	\label{eq:kernels}
	% \end{equation}
	% Here $\gamma>0$ is a scale parameter, $q\geq1$ is an integer polynomial
	% degree, $c_0\geq0$ is an offset, and $A>0$ is the scale parameter of the
	% wavelet kernel. The RBF kernel corresponds to the Gaussian kernel. The
	% exponential and Laplacian kernels differ in the distance used in their
	% construction: the former uses the Euclidean ($L_2$) norm, whereas the
	% latter uses the $L_1$ norm.

	% Each kernel in Equation~\eqref{eq:kernels} is positive semidefinite and
	% therefore defines an associated reproducing-kernel Hilbert space
	% $\mathcal H_K$. In particular, the wavelet kernel is constructed as a
	% product across coordinates of cosine-modulated Gaussian kernels;
	% positive semidefiniteness is preserved under products of
	% positive-semidefinite kernels.

	% The semiparametric version of the generalized representer theorem
	% \cite{scholkopf2001generalized} allows the RKHS component to be
	% augmented by a finite-dimensional, unpenalized parametric component.
	% Let
	% $
	% \{\psi_\ell\}_{\ell=1}^{M}
	% $
	% be a collection of real-valued functions on the covariate space
	% $\mathcal X$, and suppose that the $n\times M$ matrix
	% $
	% \bm\Psi
	% =
	% \bigl(\psi_\ell(\bm{x}_i)\bigr)_
	% {\substack{i\in\left[n\right]\\\ell\in\left[M\right]}}
	% $
	% has rank $M$. Consider a predictor of the form
	% \begin{equation*}
	% 	\widetilde f=f+h,
	% 	\qquad
	% 	f\in\mathcal H_K,
	% 	\qquad
	% 	h\in\operatorname{span}\{\psi_1,\ldots,\psi_M\},
	% \end{equation*}
	% obtained by minimizing a regularized criterion of the form
	% \begin{equation*}
	% 	\mathcal C
	% 	\left(
	% 	(\bm{x}_1,y_1,\widetilde f(\bm{x}_1)),
	% 	\ldots,
	% 	(\bm{x}_n,y_n,\widetilde f(\bm{x}_n))
	% 	\right)
	% 	+
	% 	G\!\left(\|f\|_{\mathcal H_K}\right),
	% \end{equation*}
	% where $G$ is strictly increasing. Under these conditions, the
	% semiparametric representer theorem implies that any minimizer admits
	% the representation
	% \begin{equation}
	% 	\widetilde f^{\,\star}(\bm{x})
	% 	=
	% 	\sum_{j=1}^{n}
	% 	\alpha_j K(\bm{x},\bm{x}_j)
	% 	+
	% 	\sum_{\ell=1}^{M}
	% 	\beta_\ell\psi_\ell(\bm{x}).
	% 	\label{eq:semiparametric-representer}
	% \end{equation}
	% Thus, only the component belonging to $\mathcal H_K$ is regularized,
	% whereas the finite-dimensional component spanned by the functions
	% $\psi_\ell$ remains unpenalized.

	A semiparametric extension of the generalized representer theorem
	\cite{scholkopf2001generalized} allows an RKHS component to be combined with
	an unpenalized finite-dimensional parametric component. Let
	$
	    \mathcal H_0
	    =
	    \operatorname{span}\{\psi_1,\ldots,\psi_M\},
	$
	where $\psi_1,\ldots,\psi_M:\mathcal X\to\mathbb R$, and define
	\begin{equation*}
	    \Psi_{i\ell}
	    =
	    \psi_\ell(\bm{x}_i),
	    \qquad
	    i\in \left[n\right],
	    \quad
	    \ell\in\left[M\right].
	\end{equation*}
	Assume that the resulting $n\times M$ matrix $\bm\Psi$ has full column rank
	$M$, so that the coefficients of the finite-dimensional component are
	identifiable on the observed sample.

	Consider predictors of the form
	\begin{equation*}
	    \widetilde f
	    =
	    f+h,
	    \qquad
	    f\in\mathcal H_K,
	    \qquad
	    h\in\mathcal H_0,
	\end{equation*}
	and suppose that estimation is based on minimizing a criterion of the form
	\begin{equation*}
	    \mathcal C
	    \left(
	        \widetilde f(\bm{x}_1),
	        \ldots,
	        \widetilde f(\bm{x}_n)
	    \right)
	    +
	    G\!\left(\|f\|_{\mathcal H_K}\right),
	\end{equation*}
	where the data are regarded as fixed arguments of $\mathcal C$ and
	$G:[0,\infty)\to\mathbb R$ is strictly increasing. Under these conditions,
	the generalized representer theorem implies that the RKHS component of any
	minimizer can be represented using only kernel sections evaluated at the
	training observations. Consequently, a minimizer admits the representation
	\begin{equation}
	    \widetilde f^{\,\star}(\bm{x})
	    =
	    \sum_{j=1}^{n}
	        \alpha_j K(\bm{x},\bm{x}_j)
	    +
	    \sum_{\ell=1}^{M}
	        \beta_\ell\psi_\ell(\bm{x}).
	    \label{eq:semiparametric-representer}
	\end{equation}
	Thus, regularization is applied only to the RKHS component, whereas the
	finite-dimensional component remains unpenalized.

	For the GG-KM model, estimation is formulated as the minimization of a
	penalized negative log-likelihood over the latent-cause predictor
	$\eta$. Let
	$\ell(\eta;\mathcal D_{\text{obs}})$ denote the corresponding observed-data
	log-likelihood. The corresponding penalized estimation problem is
	\begin{equation}
	    \min_{\eta}
	    \left\{
	        -\ell(\eta;\mathcal D_{\text{obs}})
	        +
	        \frac{\lambda}{2}
	        \|f\|_{\mathcal H_K}^{2}
	    \right\},
	    \qquad
	    \lambda>0,
	    \label{eq:ggkm-functional-objective}
	\end{equation}
	where the penalization is applied only to the RKHS component of
	$\eta$.

	To place this optimization problem within the semiparametric
	representer-theorem framework, the unpenalized finite-dimensional
	space is chosen to contain only an intercept. Specifically, setting
	$
	    M=1,
		\text{ }
	    \psi_1(\bm{x})=1,
	$
	so that
	$
	    \mathcal H_0
	    =
	    \operatorname{span}\{1\}.
	$
	Consequently, the class of predictors over which
	Equation~\eqref{eq:ggkm-functional-objective} is minimized is
	\begin{equation}
	    \eta(\bm{x})
	    =
	    \beta_0+f(\bm{x}),
	    \qquad
	    \beta_0\in\mathbb R,
	    \qquad
	    f\in\mathcal H_K.
	    \label{eq:ggkm-semiparametric-predictor}
	\end{equation}
	Thus, the intercept constitutes the unpenalized finite-dimensional
	component, whereas $f$ contains the possibly nonlinear covariate
	effect represented in the RKHS. In this case, the matrix
	$\bm\Psi$ reduces to the $n$-vector of ones and has rank one for
	every nonempty sample.

	Substituting
	Equation~\eqref{eq:ggkm-semiparametric-predictor} into the likelihood,
	the estimation problem can be written explicitly as
	\begin{equation}
	    (\widehat{\beta}_0,\widehat f)
	    =
	    \arg\min_{
	        \beta_0\in\mathbb R,\,
	        f\in\mathcal H_K
	    }
	    \left\{
	        -\ell(\beta_0,f;\mathcal D_{\text{obs}})
	        +
	        \frac{\lambda}{2}
	        \|f\|_{\mathcal H_K}^{2}
	    \right\}.
	    \label{eq:ggkm-penalized-objective}
	\end{equation}
	Here, $\ell(\beta_0,f;\mathcal D_{\text{obs}})$ denotes the same observed-data
	log-likelihood evaluated through the predictor
	$\eta(\bm{x})=\beta_0+f(\bm{x})$.

	The regularization term in
	Equation~\eqref{eq:ggkm-penalized-objective} is of the form required by
	the representer theorem. Indeed, for
	\begin{equation*}
	    G(\|f\|_{\mathcal H_K})
	    =
	    \frac{\lambda}{2}\|f\|_{\mathcal H_K}^{2},
	    \qquad
	    \|f\|_{\mathcal H_K}\geq 0,
	\end{equation*}
	$G$ is strictly increasing when $\lambda>0$. Therefore, the
	semiparametric representer theorem applies to the minimization problem
	in Equation~\eqref{eq:ggkm-penalized-objective}. It follows that the
	RKHS component of any minimizer can be expressed as
	$
	    \widehat f(\bm{x})
	    =
	    \sum_{j=1}^{n}
	        \widehat{\alpha}_j
	        K(\bm{x},\bm{x}_j).
	$
	Hence, although
	Equation~\eqref{eq:ggkm-penalized-objective} is initially posed as an
	optimization problem over the infinite-dimensional space
	$\mathcal H_K$, its solution can be sought within the finite-dimensional
	family
	\begin{equation}
	    \widehat{\eta}(\bm{x})
	    =
	    \widehat{\beta}_0
	    +
	    \sum_{j=1}^{n}
	        \widehat{\alpha}_j
	        K(\bm{x},\bm{x}_j).
	    \label{eq:rkhs-predictor}
	\end{equation}

	Let $\bm K$ denote the $n\times n$ training Gram matrix, with entries
	$
	    K_{ij}
	    =
	    K(\bm{x}_i,\bm{x}_j).
	$
	For
	$
	    f(\cdot)
	    =
	    \sum_{j=1}^{n}
	        \alpha_j K(\cdot,\bm{x}_j),
	$
	the squared RKHS norm satisfies
	$
	    \|f\|_{\mathcal H_K}^{2}
	    =
	    \bm{\alpha}^{\mathsf T}
	    \bm K
	    \bm{\alpha}.
	$
	Accordingly, the penalized estimation problem reduces to the
	finite-dimensional optimization
	\begin{equation}
	    \min_{\beta_0\in\mathbb R,\,
	         \bm{\alpha}\in\mathbb R^n}
	    \left\{
	        -\ell(\beta_0,\bm{\alpha};\mathcal D_{\text{obs}})
	        +
	        \frac{\lambda}{2}
	        \bm{\alpha}^{\mathsf T}
	        \bm K
	        \bm{\alpha}
	    \right\}.
	    \label{eq:ggkm-finite-objective}
	\end{equation}
	Thus, increasing $\lambda$ imposes stronger regularization on the
	nonparametric RKHS component of the predictor, while the intercept
	$\beta_0$ remains unpenalized.

	Finally, the GG-KM predictor is linked to the conditional mean of the
	latent number of causes through
	\begin{equation}
	    \theta(\bm{x})
	    =
	    \exp\{\eta(\bm{x})\}.
	    \label{eq:theta-link}
	\end{equation}
	Therefore, $\eta(\bm{x})$ is unrestricted on $\mathbb R$, whereas the
	exponential link guarantees $\theta(\bm{x})>0$ without requiring an
	explicit positivity constraint.


	% The estimation of $\eta\left(\bm{x}\right) = \log\left(\theta\left(\bm{x}\right)\right)$
	% is formulated through the minimization of a penalized negative log-likelihood function:

	% \begin{equation}
	% 	\min \left\{- \ell\left(\bm{x}; \beta_0, f\right) + \frac{\lambda}{2} \|f\|_{\mathcal H_K}^{2}\right\},
	% \end{equation}

	% For the GG-KM model, the finite-dimensional component consists only of
	% an intercept. Setting
	% $M=1,
	% \text{ }
	% \psi_1(\bm{x})=1,
	% $
	% gives
	% $
	% h(\bm{x})=\beta_0.
	% $
	% The corresponding matrix $\bm\Psi$ is the $n$-vector of ones and has
	% rank one for any nonempty sample, so the rank condition of the
	% semiparametric representer theorem is satisfied.

	% Identifying the complete predictor $\widetilde f$ with the GG-KM
	% predictor $\eta$, therefore
	% \begin{equation*}
	% 	\eta(\bm{x})
	% 	=
	% 	\widetilde f(\bm{x})
	% 	=
	% 	f(\bm{x})+\beta_0,
	% 	\qquad
	% 	f\in\mathcal H_K.
	% \end{equation*}
	% The corresponding penalized objective is
	% \begin{equation*}
	% 	\mathcal J(\beta_0,f)
	% 	=
	% 	\mathcal L\!\left(
	% 	\beta_0+f(\bm{x}_1),\ldots,
	% 	\beta_0+f(\bm{x}_n)
	% 	\right)
	% 	+
	% 	\frac{\lambda}{2}
	% 	\|f\|_{\mathcal H_K}^{2},
	% 	\qquad
	% 	\lambda>0.
	% \end{equation*}
	% Since the penalty depends only on the RKHS component $f$, the intercept
	% $\beta_0$ is not regularized. By
	% Equation~\eqref{eq:semiparametric-representer}, the RKHS component of
	% the minimizer has the finite-dimensional representation
	% \begin{equation*}
	% 	f^\star(\bm{x})
	% 	=
	% 	\sum_{j=1}^{n}
	% 	\alpha_jK(\bm{x},\bm{x}_j),
	% \end{equation*}
	% and consequently the GG-KM predictor takes the form
	% \begin{equation}
	% 	\eta(\bm{x})
	% 	=
	% 	\beta_0
	% 	+
	% 	\sum_{j=1}^{n}
	% 	\alpha_jK(\bm{x},\bm{x}_j).
	% 	\label{eq:rkhs-predictor}
	% \end{equation}

	% If $\bm K$ denotes the training Gram matrix, with
	% $
	% K_{ij}=K(\bm{x}_i,\bm{x}_j),
	% $
	% then
	% \begin{equation*}
	% 	\|f\|_{\mathcal H_K}^{2}
	% 	=
	% 	\bm{\alpha}^{\mathsf T}
	% 	\bm K
	% 	\bm{\alpha},
	% \end{equation*}
	% and the regularization term becomes
	% $
	% \frac{\lambda}{2}
	% \bm{\alpha}^{\mathsf T}
	% \bm K
	% \bm{\alpha}.
	% $
	% Hence, increasing $\lambda$ imposes stronger regularization on the
	% RKHS component of the predictor while leaving the intercept
	% $\beta_0$ unpenalized.

	% For $\lambda>0$, the quadratic penalty is a strictly increasing
	% function of $\|f\|_{\mathcal H_K}$ and therefore satisfies the
	% condition required by the representer theorem.

	% Finally, the mean parameter of the latent number of competing causes
	% is related to the predictor through
	% \begin{equation*}
	% 	\theta(\bm{x})
	% 	=
	% 	\exp\{\eta(\bm{x})\},
	% \end{equation*}
	% or equivalently, $\eta(\bm{x})=\log\theta(\bm{x})$.
	% Thus, $\eta(\bm{x})$ is unrestricted on the real line, while the
	% exponential link guarantees
	% $\theta(\bm{x})>0$ without requiring an additional positivity
	% constraint.
}

{The count law converts the intensity into a population survival function and a cure fraction through its generating function. Writing $\theta_i=\theta(\bm{x}_i)$, Poisson counts have generating function
\begin{equation*}
A_i(u)=\exp\{\theta_i(u-1)\},
\end{equation*}
so that substituting $u=S_{GG}(t)$ in Equation~\eqref{eq:promotion-population} for the population survival and $u=0$ for the cure fraction gives
\begin{equation}
\begin{aligned}
S_{\mathrm{P}}(t\mid\bm{x})
&=\exp\{-\theta(\bm{x})F_{GG}(t)\},\\
p_{c,\mathrm{P}}(\bm{x})
&=\exp\{-\theta(\bm{x})\},
\end{aligned}
\label{eq:poisson-survival-cure}
\end{equation}
with population density and hazard
\begin{equation*}
\begin{aligned}
f_{\mathrm{pop}}(t\mid\bm{x})
&=\theta(\bm{x})f_{GG}(t)\exp\{-\theta(\bm{x})F_{GG}(t)\},\\
h_{\mathrm{pop}}(t\mid\bm{x})
&=\theta(\bm{x})f_{GG}(t),
\end{aligned}
\end{equation*}
so the Poisson population hazard is the latency density scaled by the intensity, and $\mathrm{Var}(N_i\mid\bm{x}_i)=\theta_i$. Negative-binomial counts with dispersion $\phi>0$ have generating function
\begin{equation*}
A_i(u)=\{1+\phi\theta_i(1-u)\}^{-1/\phi},
\end{equation*}
with mean $\theta_i$ and variance $\theta_i(1+\phi\theta_i)$, so $\phi$ controls overdispersion relative to the Poisson case; the same substitutions give
\begin{equation}
\begin{aligned}
S_{\mathrm{NB}}(t\mid\bm{x})
&=\{1+\phi\theta(\bm{x})F_{GG}(t)\}^{-1/\phi},\\
p_{c,\mathrm{NB}}(\bm{x})
&=\{1+\phi\theta(\bm{x})\}^{-1/\phi},
\end{aligned}
\label{eq:nb-survival-cure}
\end{equation}
with population density and hazard
\begin{equation*}
\begin{aligned}
f_{\mathrm{pop}}(t\mid\bm{x})
&=\theta(\bm{x})f_{GG}(t)\{1+\phi\theta(\bm{x})F_{GG}(t)\}^{-1/\phi-1},\\
h_{\mathrm{pop}}(t\mid\bm{x})
&=\frac{\theta(\bm{x})f_{GG}(t)}{1+\phi\theta(\bm{x})F_{GG}(t)}.
\end{aligned}
\end{equation*}
The negative-binomial law reduces to the Poisson law as $\phi\to0$, so that $S_{\mathrm{NB}}\to S_{\mathrm{P}}$ and $p_{c,\mathrm{NB}}\to p_{c,\mathrm{P}}$, giving a direct check on the two count models. The cure fraction is strictly positive at every finite $\theta(\bm{x})$ and approaches one only as $\theta(\bm{x})\to0$.}

\subsection{Estimation Procedures}\label{sec:penalized-estimation}

Estimation of the GG-KM model is based on the EM algorithm, with
regularization incorporated through a penalized EM procedure. The
Poisson specification is estimated using penalized EM, whereas the
negative-binomial specification is estimated using the
Expectation/Conditional Maximization Either (ECME) algorithm
\cite{mclachlan2008algorithm}. The latter is particularly useful for
the negative-binomial model because the additional dispersion parameter
$\phi$ leads to a more complicated complete-data maximization step.

Let
$
    \mathcal D_{\mathrm{comp}}
    =
    \{(t_i,\delta_i,N_i,\bm{x}_i)\}_{i=1}^{n}
$
denote the corresponding complete data obtained by augmenting the
observations with the latent number of causes $N_i$. The parameter
vectors for the Poisson and negative-binomial models are respectively
$
    \bm\Theta_{\mathrm P}
    =
    (\beta_0,\bm{\alpha},a,d,p)
$
and
$
    \bm\Theta_{\mathrm{NB}}
    =
    (\beta_0,\bm{\alpha},a,d,p,\phi).
$
For both specifications,
$
    \eta_i=
    \eta(\bm{x}_i)
    =
    \beta_0+\bm K_i^{\mathsf T}\bm{\alpha},
    \text{ }
    \theta_i=\exp(\eta_i),
$
where
$
    \bm K_i
    =
    \bigl(
        K(\bm{x}_i,\bm{x}_1),
        \ldots,
        K(\bm{x}_i,\bm{x}_n)
    \bigr)^{\mathsf T}.
$

The target criterion is the observed penalized log-likelihood, obtained
by combining the count distribution, the generalized gamma latency
distribution in Equation~\eqref{eq:gg-distribution}, the kernel
predictor in Equation~\eqref{eq:rkhs-predictor}, and the RKHS penalty.
For the Poisson model,
\begin{equation}
    \ell^{\mathrm{pen}}_{\mathrm P}
    (\bm\Theta_{\mathrm P})
    =
    \sum_{i=1}^{n}
    \left[
        \delta_i
        \{\eta_i+\log f_{GG}(t_i)\}
        -
        \theta_iF_{GG}(t_i)
    \right]
    -
    \frac{\lambda}{2}
    \bm{\alpha}^{\mathsf T}
    \bm K
    \bm{\alpha}.
    \label{eq:poisson-observed-likelihood}
\end{equation}
For the negative-binomial model,
\begin{equation}
    \begin{aligned}
        \ell^{\mathrm{pen}}_{\mathrm{NB}}
        (\bm\Theta_{\mathrm{NB}})
        ={}&
        \sum_{i=1}^{n}
        \left[
            \delta_i
            \{\eta_i+\log f_{GG}(t_i)\}
        \right.\\
        &\left.
            \qquad-
            \left(
                \frac{1}{\phi}+\delta_i
            \right)
            \log
            \{1+\phi\theta_iF_{GG}(t_i)\}
        \right]
        -
        \frac{\lambda}{2}
        \bm{\alpha}^{\mathsf T}
        \bm K
        \bm{\alpha},
    \end{aligned}
    \label{eq:nb-observed-likelihood}
\end{equation}
with $\phi>0$. Here and throughout this subsection,
$f_{GG}(t_i)$, $F_{GG}(t_i)$, and $S_{GG}(t_i)$ are evaluated at the
current generalized gamma parameters $(a,d,p)$.

Direct optimization of these observed-data criteria is complicated by
the latent counts $N_i$. The EM construction instead uses the
conditional expectation of the complete-data log-likelihood. Given the
current parameter estimate $\bm\Theta^{(r)}$, define
\begin{equation}
    Q_{\lambda}
    \left(
        \bm\Theta
        \mid
        \bm\Theta^{(r)}
    \right)
    =
    \mathbb E_{\bm\Theta^{(r)}}
    \left[
        \ell_c
        \left(
            \bm\Theta;
            \mathcal D_{\mathrm{comp}}
        \right)
        \,\middle|\,
        \mathcal D_{\mathrm{obs}}
    \right]
    -
    \frac{\lambda}{2}
    \bm{\alpha}^{\mathsf T}
    \bm K
    \bm{\alpha}.
    \label{eq:penalized-Q-function}
\end{equation}
The E-step therefore replaces functions of the unobserved $N_i$ by
their conditional expectations under the current parameter values.

For the Poisson specification, the complete-data log-likelihood,
up to terms that do not depend on the model parameters, is
\begin{equation}
    \begin{aligned}
        \ell_{c,\mathrm P}
        ={}&
        \sum_{i=1}^{n}
        \left\{
            N_i\eta_i-\theta_i
        \right\}\\
        &+
        \sum_{i=1}^{n}
        \left\{
            \delta_i\log f_{GG}(t_i)
            +
            (N_i-\delta_i)
            \log S_{GG}(t_i)
        \right\}.
    \end{aligned}
    \label{eq:poisson-complete-data}
\end{equation}
The omitted terms, such as those involving the factorial of $N_i$,
are constant with respect to $\bm\Theta_{\mathrm P}$ when the complete
data are held fixed.

At iteration $r$, the conditional expectation of the latent count for
the Poisson model is
\begin{equation}
    \widehat N_{i,\mathrm P}^{(r)}
    =
    \mathbb E
    \left[
        N_i
        \mid
        t_i,\delta_i,\bm{x}_i;
        \bm\Theta_{\mathrm P}^{(r)}
    \right]
    =
    \delta_i
    +
    \theta_i^{(r)}
    S_{GG}^{(r)}(t_i).
    \label{eq:poisson-estep}
\end{equation}
For the negative-binomial specification, the corresponding conditional
expectation is
\begin{equation}
    \widehat N_{i,\mathrm{NB}}^{(r)}
    =
    \delta_i
    +
    \frac{
        \bigl(1+\phi^{(r)}\delta_i\bigr)
        \theta_i^{(r)}
        S_{GG}^{(r)}(t_i)
    }{
        1+
        \phi^{(r)}
        \theta_i^{(r)}
        F_{GG}^{(r)}(t_i)
    }.
    \label{eq:nb-estep}
\end{equation}

Once the conditional counts have been computed, the expected
complete-data objective separates naturally into a predictor block and
a generalized gamma latency block. For notational simplicity, the
iteration index on $\widehat N_i$ is suppressed in the following
expressions.

For the Poisson model, the terms depending on
$(\beta_0,\bm{\alpha})$ are
\begin{equation}
    Q_{\eta,\mathrm P}
    =
    \sum_{i=1}^{n}
    \left\{
        \widehat N_i\eta_i
        -
        \exp(\eta_i)
    \right\}
    -
    \frac{\lambda}{2}
    \bm{\alpha}^{\mathsf T}
    \bm K
    \bm{\alpha}.
    \label{eq:poisson-predictor-block}
\end{equation}
For the negative-binomial model, conditional on the current value of
$\phi$, the corresponding predictor block is
\begin{equation}
    Q_{\eta,\mathrm{NB}}
    =
    \sum_{i=1}^{n}
    \left[
        \widehat N_i\eta_i
        -
        \left(
            \widehat N_i+\frac{1}{\phi}
        \right)
        \log\{1+\phi\exp(\eta_i)\}
    \right]
    -
    \frac{\lambda}{2}
    \bm{\alpha}^{\mathsf T}
    \bm K
    \bm{\alpha},
    \label{eq:nb-predictor-block}
\end{equation}
up to terms that do not depend on $(\beta_0,\bm{\alpha})$.

For fixed conditional counts, both count models share the same
generalized gamma latency block,
\begin{equation}
    Q_{GG}
    =
    \sum_{i=1}^{n}
    \left\{
        \delta_i\log f_{GG}(t_i)
        +
        (\widehat N_i-\delta_i)
        \log S_{GG}(t_i)
    \right\}.
    \label{eq:latency-block}
\end{equation}
Consequently, $(\beta_0,\bm{\alpha})$ appears only in the corresponding
predictor block, whereas $(a,d,p)$ appears only in
Equation~\eqref{eq:latency-block} once $\widehat N_i$ is held fixed.
This separation permits the predictor and latency parameters to be
updated in distinct maximization blocks without altering the underlying
penalized EM objective.

For the Poisson specification, these conditional maximizations
constitute the penalized EM update. The negative-binomial specification
contains the additional dispersion parameter $\phi$. Its complete-data
count contribution contains terms involving $\phi$ and gamma functions
of the latent count, which makes direct maximization of the expected
complete-data criterion with respect to $\phi$ considerably less
convenient. The ECME algorithm avoids this additional expectation by
updating $\phi$ through the observed-data log-likelihood rather than
through the corresponding expected complete-data block.

Specifically, after updating
$(\beta_0,\bm{\alpha})$ and $(a,d,p)$ conditionally on the current
dispersion, the ECME update is
\begin{equation}
    \phi^{(r+1)}
    =
    \underset{\phi>0}{\arg\max}
    \;
    \ell^{\mathrm{pen}}_{\mathrm{NB}}
    \left(
        \beta_0^{(r+1)},
        \bm{\alpha}^{(r+1)},
        a^{(r+1)},
        d^{(r+1)},
        p^{(r+1)},
        \phi
    \right).
    \label{eq:nb-ecme-phi-update}
\end{equation}
Because the RKHS penalty does not depend on $\phi$, this update is
equivalent to maximizing the unpenalized observed-data log-likelihood
with respect to $\phi$ while all remaining parameters are held fixed.

The Supporting Information provides the complete-data derivations,
the conditional expectations in
Equations~\eqref{eq:poisson-estep}--\eqref{eq:nb-estep}, and the
analytical gradients used to optimize the predictor and latency
blocks.

\begin{algorithm*}[tbp]
	\caption{Penalized EM algorithm for the Poisson GG-KM model}
	\label{alg:poisson-em}
	\begin{algorithmic}[1]

		\Require Training data $(\bm X,\bm t,\bm\delta)$, kernel configuration,
		$\lambda>0$, tolerance $\varepsilon>0$, and maximum number of
		iterations $k_{\max}$

		\State Construct $\bm K$ and initialize
		$(\beta_0^{(0)},\bm\alpha^{(0)},\bm\xi^{(0)})$, with
		$\bm\xi=(\log a,\log d,\log p)^\top$; evaluate
		$\ell_{\mathrm P}^{\mathrm{pen},(0)}$.

		\For{$k=0,\ldots,k_{\max}-1$}

			\State Set
			$\eta_i^{(k)}
			=\beta_0^{(k)}+\bm K_i^\top\bm\alpha^{(k)}$ and
			$(a^{(k)},d^{(k)},p^{(k)})=\exp\{\bm\xi^{(k)}\}$.

			\State \textbf{E-step:} for $i=1,\ldots,n$, compute
			\[
			\widehat N_i^{(k+1)}
			=
			\delta_i+
			\exp\{\eta_i^{(k)}\}
			S_{GG}\!\left(
				t_i;a^{(k)},d^{(k)},p^{(k)}
			\right).
			\]

			\State \textbf{M-step:} with $\widehat N_i^{(k+1)}$ fixed, update
			\[
			(\beta_0^{(k+1)},\bm\alpha^{(k+1)})
			=
			\arg\max_{\beta_0,\bm\alpha}
			\left\{
			\sum_{i=1}^{n}
			\left[
				\widehat N_i^{(k+1)}\eta_i-e^{\eta_i}
			\right]
			-
			\frac{\lambda}{2}
			\bm\alpha^\top\bm K\bm\alpha
			\right\}.
			% \qquad
			% \eta_i=\beta_0+\bm K_i^\top\bm\alpha.
			\]

			\State Update the latency parameters on the log scale:
			\[
			\bm\xi^{(k+1)}
			=
			\arg\max_{\bm\xi}
			\sum_{i=1}^{n}
			\left[
				\delta_i
				\log f_{GG}\!\left(t_i;e^{\xi_a},e^{\xi_d},e^{\xi_p}\right)
				+
				(\widehat N_i^{(k+1)}-\delta_i)
				\log S_{GG}\!\left(t_i;e^{\xi_a},e^{\xi_d},e^{\xi_p}\right)
			\right].
			\]

			\State Set
			$(a^{(k+1)},d^{(k+1)},p^{(k+1)})
			=\exp\{\bm\xi^{(k+1)}\}$ and evaluate
			$\ell_{\mathrm P}^{\mathrm{pen},(k+1)}$.

			\If{
			$\left|
			\ell_{\mathrm P}^{\mathrm{pen},(k+1)}
			-
			\ell_{\mathrm P}^{\mathrm{pen},(k)}
			\right|
			\leq\varepsilon$
			}
				\State \textbf{break}
			\EndIf

		\EndFor

		\State \Return
		$(\widehat\beta_0,\widehat{\bm\alpha},
		\widehat a,\widehat d,\widehat p)$.

	\end{algorithmic}
\end{algorithm*}

\begin{algorithm*}[tbp]
	\caption{Penalized ECME algorithm for the negative-binomial GG-KM model}
	\label{alg:nb-ecme}
	\begin{algorithmic}[1]

		\Require Training data $(\bm X,\bm t,\bm\delta)$, kernel configuration,
		$\lambda>0$, tolerance $\varepsilon>0$, and maximum number of
		iterations $k_{\max}$

		\State Construct $\bm K$ and initialize
		$(\beta_0^{(0)},\bm\alpha^{(0)},\bm\xi^{(0)},\zeta^{(0)})$, with
		$\bm\xi=(\log a,\log d,\log p)^\top$ and $\zeta=\log\phi$;
		evaluate $\ell_{\mathrm{NB}}^{\mathrm{pen},(0)}$.

		\For{$k=0,\ldots,k_{\max}-1$}

			\State Set
			$\eta_i^{(k)}
			=\beta_0^{(k)}+\bm K_i^\top\bm\alpha^{(k)}$,
			$(a^{(k)},d^{(k)},p^{(k)})=\exp\{\bm\xi^{(k)}\}$, and
			$\phi^{(k)}=e^{\zeta^{(k)}}$.

			\State \textbf{E-step:} for $i=1,\ldots,n$, compute
			\[
			\widehat N_i^{(k+1)}
			=
			\delta_i+
			\frac{
				(1+\phi^{(k)}\delta_i)
				e^{\eta_i^{(k)}}
				S_{GG}(t_i;a^{(k)},d^{(k)},p^{(k)})
			}{
				1+\phi^{(k)}
				e^{\eta_i^{(k)}}
				F_{GG}(t_i;a^{(k)},d^{(k)},p^{(k)})
			}.
			\]

			\State \textbf{M-step:} with $\widehat N_i^{(k+1)}$ and
			$\phi^{(k)}$ fixed, update
			\[
			(\beta_0^{(k+1)},\bm\alpha^{(k+1)})
			=
			\arg\max_{\beta_0,\bm\alpha}
			\left\{
			\sum_{i=1}^{n}
			\left[
				\widehat N_i^{(k+1)}\eta_i
				-
				\left(
					\widehat N_i^{(k+1)}
					+\frac{1}{\phi^{(k)}}
				\right)
				\log\{1+\phi^{(k)}e^{\eta_i}\}
			\right]
			-
			\frac{\lambda}{2}
			\bm\alpha^\top\bm K\bm\alpha
			\right\}.
			% \qquad
			% \eta_i=\beta_0+\bm K_i^\top\bm\alpha.
			\]

			\State Update the latency parameters on the log scale:
			\[
			\bm\xi^{(k+1)}
			=
			\arg\max_{\bm\xi}
			\sum_{i=1}^{n}
			\left[
				\delta_i
				\log f_{GG}(t_i;e^{\xi_a},e^{\xi_d},e^{\xi_p})
				+
				(\widehat N_i^{(k+1)}-\delta_i)
				\log S_{GG}(t_i;e^{\xi_a},e^{\xi_d},e^{\xi_p})
			\right].
			\]

			\State \textbf{ECME-step:} update the dispersion on the log scale
			by maximizing the observed penalized log-likelihood:
			\[
			\zeta^{(k+1)}
			=
			\arg\max_{\zeta\in\mathbb R}
			\ell_{\mathrm{NB}}^{\mathrm{pen}}
			\left(
			    \beta_0^{(k+1)},
			    \bm\alpha^{(k+1)},
			    e^{\bm\xi^{(k+1)}},
			    e^\zeta
			\right).
			\]

			\State Set
			$(a^{(k+1)},d^{(k+1)},p^{(k+1)})
			=\exp\{\bm\xi^{(k+1)}\}$ and
			$\phi^{(k+1)}=e^{\zeta^{(k+1)}}$;
			evaluate $\ell_{\mathrm{NB}}^{\mathrm{pen},(k+1)}$.

			\If{
			$\left|
			\ell_{\mathrm{NB}}^{\mathrm{pen},(k+1)}
			-
			\ell_{\mathrm{NB}}^{\mathrm{pen},(k)}
			\right|
			\leq\varepsilon$
			}
				\State \textbf{break}
			\EndIf

		\EndFor

		\State \Return
		$(\widehat\beta_0,\widehat{\bm\alpha},
		\widehat a,\widehat d,\widehat p,\widehat\phi)$.

	\end{algorithmic}
\end{algorithm*}

Algorithms~\ref{alg:poisson-em} and \ref{alg:nb-ecme} summarize the
penalized EM and ECME procedures for the Poisson and negative-binomial
GG-KM models, respectively.

The simulations and application developed throughout this work use different
families of kernels to represent distinct forms of similarity between
covariate vectors. Let $\bm{x},\bm{x}'\in\mathcal X\subseteq\mathbb R^r$.
The kernel families considered are
\begin{equation}
	\begin{aligned}
		K_{\mathrm{lin}}(\bm{x},\bm{x}')
		&=\bm{x}^{\mathsf T}\bm{x}',\\
		K_{\mathrm{poly}}(\bm{x},\bm{x}')
		&=(\gamma\,\bm{x}^{\mathsf T}\bm{x}'+c_0)^q,\\
		K_{\mathrm{rbf}}(\bm{x},\bm{x}')
		&=\exp\{-\gamma\|\bm{x}-\bm{x}'\|_2^2\},\\
		K_{\mathrm{exp}}(\bm{x},\bm{x}')
		&=\exp\{-\gamma\|\bm{x}-\bm{x}'\|_2\},\\
		K_{\mathrm{lap}}(\bm{x},\bm{x}')
		&=\exp\{-\gamma\|\bm{x}-\bm{x}'\|_1\},\\
		K_{\mathrm{cau}}(\bm{x},\bm{x}')
		&=\{1+\gamma\|\bm{x}-\bm{x}'\|_2^2\}^{-1},\\
		K_{\mathrm{wav}}(\bm{x},\bm{x}')
		&=
		\prod_{j=1}^{r}
		\left[
		\cos\!\left(
		1.75\,\frac{x_j-x_j'}{A}
		\right)
		\exp\!\left\{
		-\frac{1}{2}
		\left(
		\frac{x_j-x_j'}{A}
		\right)^2
		\right\}
		\right].
	\end{aligned}
	\label{eq:kernels}
\end{equation}
Here, $\gamma>0$ is a kernel scale parameter, $q$ is the polynomial degree,
$c_0$ is the polynomial offset, and $A>0$ controls the scale of the
wavelet kernel. The RBF kernel corresponds to the Gaussian kernel and
depends on the squared Euclidean distance, whereas the exponential and
Laplacian kernels use the Euclidean and $L_1$ distances, respectively.
Kernel-specific hyperparameters are selected within the corresponding
tuning procedure.

The GG-KM model is implemented in the open-source Python package \texttt{scikit-cure},
available at \href{https://github.com/cowvin0/scikit-cure}{GitHub}.

\subsection{Prediction, tuning, and evaluation}
\label{sec:prediction-evaluation}

Prediction for an unseen covariate vector uses the training expansion
$
\eta(\bm{x}_{\mathrm{new}})
=
\widehat\beta_0
+
K(\bm{x}_{\mathrm{new}},\bm{X}_{\mathrm{train}})
\widehat{\bm{\alpha}}.
$
Equations~\eqref{eq:poisson-survival-cure} and
\eqref{eq:nb-survival-cure} then produce the individual cure fraction
and population survival curve. Covariate preprocessing is estimated
from the training data and applied unchanged to validation and test
data.

Predictive evaluation uses common event-stratified outer splits with
$70\%$ of observations for training and $30\%$ for testing. Kernel
parameters, the penalty, model variants, and all preprocessing choices
are selected inside the outer training set by a shared inner tuning
protocol and comparable search budget. The same outer test observations
and evaluation grid are used for every paired method comparison.

Let $\widehat G$ be the Kaplan--Meier estimate of the censoring survival
distribution fitted from the outer training data, with censoring treated
as the event. The weighting convention follows the prediction-error
construction of Graf and colleagues.\cite{graf1999} For a test set of
size $m$, the inverse-probability-of-censoring weighted Brier score is
\begin{equation}
	\begin{aligned}
		\widehat{\mathrm{BS}}(t)
		=
		\frac{1}{m}\sum_{i=1}^{m}
		\Biggl\{
		&\frac{I(t_i\leq t,\delta_i=1)}
		{\widehat G(t_i-)}
		\widehat S_i(t)^2\\
		+{}&
		\frac{I(t_i>t)}
		{\widehat G(t)}
		[1-\widehat S_i(t)]^2
		\Biggr\}.
	\end{aligned}
	\label{eq:ipcw-brier}
\end{equation}
The integration horizon is selected from the outer training data and
restricted to times with reliable censoring support. The integrated
Brier score is the time average obtained by trapezoidal integration of
Equation~\eqref{eq:ipcw-brier} over the frozen grid and division by the
horizon length, without renormalizing by the realized sum of weights.

Discrimination with respect to long-term cure status is additionally
assessed using an imputation-based area under the receiver operating
characteristic curve for cure models, denoted by
$\mathrm{AUC}_{\mathrm{cure}}$ \cite{medina2024}. Let
\begin{equation}
	\widehat\pi_i
	=
	\lim_{t\to\infty}
	\widehat S_i(t)
\end{equation}
denote the estimated cure probability for subject $i$. For the Poisson
GG-KM model,
$
\widehat\pi_i
=
\exp\{-\exp(\widehat\eta_i)\},
$
whereas for the negative-binomial GG-KM model,
$
\widehat\pi_i
=
\left(
1+\widehat\phi\exp(\widehat\eta_i)
\right)^{-1/\widehat\phi}
$.
Because cure status is not directly observed for censored individuals,
$\widehat\pi_i$ and $1-\widehat\pi_i$ are used as imputation weights for
the cured and susceptible states, respectively. For a cut-off
$c\in[0,1]$, the estimated true-positive and false-positive rates are
\begin{equation}
	\widehat{\mathrm{TPR}}(c)
	=
	\frac{
		\displaystyle
		\sum_{i=1}^{m}
		I(\widehat\pi_i\leq c)
		(1-\widehat\pi_i)
	}{
		\displaystyle
		\sum_{i=1}^{m}
		(1-\widehat\pi_i)
	},
	\label{eq:auc-cure-tpr}
\end{equation}
\begin{equation}
	\widehat{\mathrm{FPR}}(c)
	=
	\frac{
		\displaystyle
		\sum_{i=1}^{m}
		I(\widehat\pi_i\leq c)
		\widehat\pi_i
	}{
		\displaystyle
		\sum_{i=1}^{m}
		\widehat\pi_i
	}.
	\label{eq:auc-cure-fpr}
\end{equation}
Here, lower values of $\widehat\pi_i$ correspond to stronger evidence
that an individual is susceptible, so susceptibility is treated as the
positive class. The $\mathrm{AUC}_{\mathrm{cure}}$ is obtained by
trapezoidal integration of the receiver operating characteristic curve
formed by Equations~\eqref{eq:auc-cure-tpr} and \eqref{eq:auc-cure-fpr}, that is
		$\{\widehat{\mathrm{FPR}}(c),\widehat{\mathrm{TPR}}(c)\}$
over $c\in[0,1]$. Larger values indicate better discrimination between
individuals with higher and lower probabilities of belonging to the
susceptible subpopulation.

Evaluation summaries retain paired score differences, cure-fraction
error in simulations, $\mathrm{AUC}_{\mathrm{cure}}$, calibration,
convergence, failures, and runtime. A common generalized gamma latency
is used throughout the reported simulation so that the generating
latency matches the fitted latency structure and predictor-shape
contrasts are not confounded with unmodelled latency heterogeneity.
Covariate-dependent latency was not evaluated in the reported scenarios
and remains a future misspecification stress test.

% \subsection{Prediction, tuning, and evaluation}\label{sec:prediction-evaluation}

% Prediction for an unseen covariate vector uses the training expansion $\eta(\bm{x}_{\mathrm{new}})=\widehat\beta_0+K(\bm{x}_{\mathrm{new}},\bm{X}_{\mathrm{train}})\widehat{\bm{\alpha}}$. Equations~\eqref{eq:poisson-survival-cure} and \eqref{eq:nb-survival-cure} then produce the individual cure fraction and population survival curve. Covariate preprocessing is estimated from the training data and applied unchanged to validation and test data.

% Predictive evaluation uses common event-stratified outer splits with $70\%$ of observations for training and $30\%$ for testing. Kernel parameters, the penalty, model variants, and all preprocessing choices are selected inside the outer training set by a shared inner tuning protocol and comparable search budget. The same outer test observations and evaluation grid are used for every paired method comparison.

% Let $\widehat G$ be the Kaplan--Meier estimate of the censoring survival distribution fitted from the outer training data, with censoring treated as the event. The weighting convention follows the prediction-error construction of Graf and colleagues.\cite{graf1999} For a test set of size $m$, the inverse-probability-of-censoring weighted Brier score is
% \begin{equation}
	% \begin{aligned}
		% \widehat{\mathrm{BS}}(t)=\frac{1}{m}\sum_{i=1}^{m}
		% \Biggl\{
		% &\frac{I(t_i\leq t,\delta_i=1)}{\widehat G(t_i-)}
		% \widehat S_i(t)^2\\
		% +{}&\frac{I(t_i>t)}{\widehat G(t)}
		% [1-\widehat S_i(t)]^2
		% \Biggr\}.
		% \end{aligned}
	% \label{eq:ipcw-brier}
	% \end{equation}
% The integration horizon is selected from the outer training data and restricted to times with reliable censoring support. The integrated Brier score is the time average obtained by trapezoidal integration of Equation~\eqref{eq:ipcw-brier} over the frozen grid and division by the horizon length, without renormalizing by the realized sum of weights.

% Evaluation summaries retain paired score differences, cure-fraction error in simulations, calibration, convergence, failures, and runtime. A common generalized gamma latency is used throughout the reported simulation so that the generating latency matches the fitted latency structure and predictor-shape contrasts are not confounded with unmodelled latency heterogeneity. Covariate-dependent latency was not evaluated in the reported scenarios and remains a future misspecification stress test.

\section{Simulation study}\label{sec:simulation}

The simulation study was designed to evaluate the predictive performance
of the proposed GG-KM model under different cure mechanisms, latent-count
distributions, and sample sizes. Specifically, we investigated three
questions. First, we examined whether a simpler predictor performs best
when its functional form is correctly specified. Second, we evaluated
whether GG-KM provides robust predictions when the true cure-probability
function is not adequately represented by the simpler models. Third, we
assessed the extent to which predictive performance depends on the
latent-count distribution and its degree of overdispersion. The
generalized gamma latency distribution was correctly specified and held
fixed across all scenarios. Therefore, differences in predictive
performance primarily reflect the ability of each predictor to recover
the covariate-dependent cure mechanism.

Each data-generating scenario was defined by one cure-probability surface,
one latent-count distribution, and one sample size. We considered three
cure-probability surfaces, linear, quadratic, and periodic. We also
considered three latent-count distributions, Poisson, negative binomial
with dispersion $\phi=0.25$, and negative binomial with dispersion
$\phi=1.00$. The two negative-binomial settings represent modest and
stronger overdispersion, respectively, whereas the Poisson distribution
represents the equidispersed limiting case. The sample sizes were
$n=300$ and $n=600$, resulting in $18$ data-generating scenarios.

In every scenario, two covariates were generated independently as
$x_1,x_2\sim\operatorname{Uniform}(-1,1)$. The bounded uniform domain
deliberately differs from the historical source because it gives all
three geometries a common domain for level-set and recovery comparisons.
The source used standard-normal covariates and a covariate-dependent
Weibull latency; retaining those choices would confound cure-predictor
shape with unmodelled latency heterogeneity.

The true cure probability was generated from one of three fixed surfaces.
Specifically, with $\operatorname{expit}(v)=\{1+\exp(-v)\}^{-1}$, the
designs were
\begin{equation}
	\begin{aligned}
		p_c^{\mathrm{lin}}(\bm{x})
		&=\operatorname{expit}(0.15-1.00x_1+0.75x_2),\\
		p_c^{\mathrm{quad}}(\bm{x})
		&=\operatorname{expit}(0.60-1.45x_1^2-1.05x_2^2+0.55x_1x_2),\\
		p_c^{\mathrm{per}}(\bm{x})
		&=\operatorname{expit}\{0.10+1.15\sin(\pi x_1)+0.80\cos(\pi x_2)\}.
	\end{aligned}
	\label{eq:simulation-cure-surfaces}
\end{equation}
For Poisson generation, the latent-count mean was
$\theta(\bm{x})=-\log p_c(\bm{x})$; for negative-binomial generation, it
was $\theta(\bm{x})=\{p_c(\bm{x})^{-\phi}-1\}/\phi$. This mapping
preserves the same cure-probability surface across count distributions,
so comparisons do not confound dispersion with a different cure target.

Figure~\ref{fig:cure-surfaces} shows the three surfaces on the common
covariate square. The linear surface increases along one tilted gradient
and has parallel straight level sets, so the log-linear predictor is
correctly specified for it. The quadratic surface has closed elliptical
level sets generated by the two squared terms and the interaction, and no
additive predictor reproduces that geometry. The periodic surface reverses
its gradient twice inside the domain and produces alternating regions of
high and low cure probability, a shape an additive spline can follow one
coordinate at a time but a degree-two map cannot. The three designs
therefore separate correct linear specification, smooth nonadditivity and
oscillation while holding the latency distribution fixed. They are not
matched on prevalence: the mean of $p_c(\bm{x})$ over the square is
0.533 for the linear surface, 0.448 for the quadratic surface and 0.520
for the periodic surface, with attainable ranges $[0.168,0.870]$,
$[0.079,0.646]$ and $[0.136,0.886]$. Scores are consequently comparable
between methods within a scenario and not between scenarios.

\GGKMWideFigure{cure-surfaces.pdf}{Frozen cure-probability surfaces used
in the simulation design. Panels show the linear, quadratic, and
periodic functions in Equation~\eqref{eq:simulation-cure-surfaces} over
$(x_1,x_2)\in[-1,1]^2$. All panels use the same $p_c(\bm{x})$ colour
scale; superimposed contour lines provide a colour-independent encoding
of the surface geometry. Contours are labelled at $p_c=0.25$, $0.50$ and
$0.60$, the same levels used in Figure~\ref{fig:cure-surface-recovery},
and each panel title reports the mean of $p_c(\bm{x})$ over the square.
}{fig:cure-surfaces}{Three contour panels show cure probability over two
covariates. The linear panel has diagonal straight contours, the quadratic
panel has curved contours with an interaction, and the periodic panel has
repeated peaks and troughs. Each panel title states the mean cure
probability for that surface.
}

All latent-cause activation times followed the generalized gamma
distribution with $(a,d,p)=(1.6,1.3,1.1)$. Independent censoring times
were exponential with mean 3.5, and observed time was the minimum of the
event and censoring times. The inverse population distribution was used
to generate finite event times directly for susceptible individuals, while
individuals in the limiting cure mass were assigned no event. Every
replicate retained the true individual cure probability and population
survival curve for evaluation.

Algorithm~\ref{alg:simulation-generation} states the implemented
inverse-transform generator. Sampling from the improper population
distribution avoids simulating a potentially large latent count while
preserving its cure mass and finite-event distribution exactly.
\begin{algorithm*}[tbp]
\caption{Right-censored simulation from a promotion-time population distribution}\label{alg:simulation-generation}
\begin{algorithmic}[1]
\Require Scenario, sample size $n$, generalized gamma parameters $(a,d,p)$, censoring scale, and deterministic replication seed
\For{$i=1,\ldots,n$}
  \State Draw $x_{i1},x_{i2}\stackrel{\mathrm{ind}}{\sim}\operatorname{Uniform}(-1,1)$ and evaluate the declared $p_c(\bm{x}_i)$
  \State Set $\theta_i=-\log p_c(\bm{x}_i)$ for Poisson counts, or $\theta_i=\{p_c(\bm{x}_i)^{-\phi}-1\}/\phi$ for negative-binomial counts
  \State Draw $U_i\sim\operatorname{Uniform}(0,1)$ from the population event-time distribution
  \If{$U_i\geq1-p_c(\bm{x}_i)$}
    \State Set $T_i=\infty$ (the draw falls in the cure mass)
  \Else
    \State Set $q_i=-\log(1-U_i)/\theta_i$ for Poisson counts
    \State Set $q_i=\{(1-U_i)^{-\phi}-1\}/(\phi\theta_i)$ for negative-binomial counts
    \State Set $T_i=F_{GG}^{-1}(q_i;a,d,p)$
  \EndIf
  \State Draw $C_i$ from the declared independent censoring distribution
  \State Record the observed time $t_i=\min(T_i,C_i)$, indicator $\delta_i=I(T_i\leq C_i)$, $p_c(\bm{x}_i)$, and the true population survival grid
\EndFor
\State Return the complete observed sample, generating truth, scenario identifier, and replication seed
\end{algorithmic}
\end{algorithm*}

The fitted comparison set shared the same generalized gamma latency and comprised
GG-KM predictors with linear, radial-basis, Laplacian, exponential, Cauchy,
polynomial, and wavelet kernels, each with Poisson and negative-binomial latent
counts, together with additive cubic-spline Poisson and negative-binomial models
and the classical mixture-cure and promotion-time comparators. Linear-kernel
candidates used $\lambda\in\{0.3,3\}$. Radial-basis, Laplacian, exponential, and
Cauchy candidates used $\gamma\in\{0.25,1\}$ with
$\lambda\in\{0.3,3\}$. Polynomial candidates additionally used
degree $q\in\{2,3\}$ with $c_0=1$, while wavelet candidates used
$A\in\{0.5,1,2\}$, in both cases with $\lambda\in\{0.3,3\}$.
Spline candidates used three or five training-fold quantile knots with
$\lambda\in\{0.3,3\}$. All tuned methods used the same two-fold inner tuning
criterion and evaluation-grid size. The promotion-time support-vector method of
Pal and Aselisewine was not included because the reproducibility audit found no
public executable implementation at the article or arXiv records, and a new
reimplementation would not provide a controlled software comparison.
\cite{pal2023}

Each generated sample was split once into a common event-stratified
$70\%$ training set and $30\%$ test set for all fitted methods. Tuning,
feature-map fitting, and preprocessing were confined to the training set,
and the IPCW integration horizon was set from that set at the smaller of
its 90th observed-time percentile and the limit at which estimated
censoring survival remained at least 0.05. Primary test-set estimands were
integrated Brier score, cure AUC, and cure-probability mean squared error.

A target-size diagnostic pilot used five planned replicates per pilot
scenario and was not pooled with the final experiment; at least four
replicates were successful in every pilot cell. The largest observed
between-replicate standard deviations were 0.0418 for integrated Brier
score and 0.0332 for cure-probability mean squared error. Requiring
projected Monte Carlo standard errors no larger than 0.003 and 0.005,
respectively, implied at least 195 replicates for integrated Brier score
and 45 for cure-probability mean squared error. The more restrictive
criterion therefore required 195 replicates, so the final count was fixed
at 250 replicates per scenario before final estimation. At
$R=250$, the corresponding projected Monte Carlo standard errors were
0.00265 for integrated Brier score and 0.00210 for cure-probability mean
squared error. For cure AUC, the largest pilot standard deviation was
0.1211, corresponding to a projected Monte Carlo standard error of
0.00766 at $R=250$. Final summaries report the realized Monte Carlo
standard error for every primary metric and paired differences computed
within each shared replication.

\subsection{Simulation results}\label{sec:simulation-results}

The primary final experiment contained 81,000 method rows from
250 replications in each of 18 scenarios. The largest realized
Monte Carlo standard errors were 0.00175 for the IPCW integrated
Brier score, 0.00434 for cure-probability RMSE, and
0.00694 for cure AUC. Of the 81,000 fits, 80,842 completed
successfully and 158 failed; among the successful fits, 80,832
satisfied the explicit convergence criterion and 10 did not.
Method-specific successful-fit totals ranged from 4,467 to 4,498
of 4,500. Mean per-fit runtime ranged from 1.13 to 40.32 seconds
across methods, with a maximum of 520.48 seconds.

Table~\ref{tab:simulation-best} reports the best observed mean IPCW
integrated Brier score, cure root mean squared error, and cure AUC in each
scenario. For the linear surface, a linear-kernel GG-KM attained the lowest
mean IPCW integrated Brier score and cure root mean squared error in all six
scenarios. For the quadratic surface, a polynomial-kernel GG-KM attained the
lowest cure root mean squared error in all six scenarios and the lowest
integrated Brier score in five of six scenarios, with an RBF-kernel GG-KM
attaining the lowest integrated Brier score in the remaining case. For the
periodic surface, an additive spline model attained the lowest integrated
Brier score and cure root mean squared error in all six scenarios. For cure
AUC, a polynomial-kernel GG-KM attained the highest mean in all six linear
scenarios and all six periodic scenarios. On the quadratic surface,
polynomial-kernel GG-KM models attained the highest mean cure AUC in four
of six scenarios, whereas wavelet-kernel GG-KM models attained the highest
mean in the remaining two.

\GGKMTable{sections/generated/simulation-best-table}

Figure~\ref{fig:simulation-paired-differences} shows paired differences
between the RBF GG-KM and the count-matched linear-kernel GG-KM,
additive-spline, and polynomial-kernel GG-KM comparators across the six
generating-law and sample-size cells. Each contrast is defined as the RBF
GG-KM value minus the comparator value; hence, negative differences favor
the RBF model for IPCW integrated Brier score and cure root mean squared
error, whereas positive differences favor it for cure AUC.

For the linear cure surface, the linear-kernel GG-KM consistently
outperformed the RBF model in integrated Brier score and cure root mean
squared error, while the spline and polynomial comparisons were generally
closer. Cure AUC also favored the linear and polynomial kernels, whereas
the RBF--spline differences remained comparatively small. Under the
quadratic surface, the RBF GG-KM substantially improved on the linear
kernel across all three metrics, but its differences from the spline and
polynomial models were much smaller; the polynomial kernel generally had
lower cure root mean squared error. For the periodic surface, the RBF
model improved markedly on the linear and polynomial kernels in integrated
Brier score and cure root mean squared error, particularly at $n=600$,
but remained inferior to the spline. Cure AUC showed a different pattern,
with the RBF model clearly outperforming the linear kernel and, at
$n=600$, the Poisson polynomial-kernel GG-KM, while its differences from
the spline remained comparatively small. Overall, the paired comparisons
show that relative performance depends strongly on both the geometry of
the cure surface and the evaluation metric.

\GGKMWideFigure
{simulation_rbf_paired_differences_with_auc.pdf}
{Count-matched paired predictive differences of the RBF GG-KM against the linear-kernel GG-KM, additive-spline, and polynomial-kernel GG-KM comparators. Columns correspond to the linear, quadratic, and periodic cure surfaces, while rows show IPCW integrated Brier score, cure root mean squared error, and cure AUC differences. Each contrast is defined as the RBF GG-KM value minus the corresponding comparator value, with Poisson and negative-binomial fits compared separately. Error bars span $\pm1.96$ paired Monte Carlo standard errors, and a horizontal line marks zero. Negative differences favor the RBF GG-KM for integrated Brier score and cure root mean squared error, whereas positive differences favor it for cure AUC.}
{fig:simulation-paired-differences}
{
A three-row by three-column grid displays count-matched paired
differences between the RBF GG-KM and linear-kernel GG-KM,
additive-spline, and polynomial-kernel GG-KM comparators. Columns
represent linear, quadratic, and periodic cure surfaces; rows represent
integrated Brier score, cure root mean squared error, and cure AUC.
Poisson and negative-binomial contrasts are distinguished by marker and
line type, with error bars spanning $\pm1.96$ paired Monte Carlo standard
errors around a horizontal zero line. Negative values favor the RBF model
for integrated Brier score and cure root mean squared error, whereas
positive values favor it for cure AUC.
}

Figure~\ref{fig:cure-surface-recovery} maps where in the
covariate square each predictor loses accuracy. Every panel
fills the $41\times41$ evaluation grid with the pointwise root
mean squared error of the predicted cure probability against
the generating surface, averaged over the valid Poisson
replications at $n=600$, on a scale shared by all twelve panels.
Solid lines repeat the true level sets at $p_c=0.25$, $0.50$
and $0.60$, and the dashed line is the mean estimated
$p_c=0.50$ level, which separates an error of magnitude from a
displacement of shape. Failed fits reduce the pointwise
$n_{\mathrm{valid}}$ rather than being removed from the total
number of replications. The panel means summarize the gridwise
recovery of the cure surface. On the linear surface the
linear-kernel GG-KM attains 0.0490, compared with 0.0671 for
the additive spline, 0.0694 for the polynomial-kernel GG-KM,
and 0.0730 for the RBF GG-KM. On the quadratic surface, the
RBF GG-KM attains the lowest mean error, 0.0802, followed
closely by the additive spline at 0.0816 and the
polynomial-kernel GG-KM at 0.0827, whereas the linear-kernel
GG-KM increases to 0.1282. On the periodic surface, the
additive spline performs best at 0.0824, followed by the RBF
GG-KM at 0.0959, while the linear- and polynomial-kernel GG-KM
predictors increase to 0.1504 and 0.1510, respectively. The
spatial pattern adds what these scalar summaries omit. A
misaligned predictor does not spread its error evenly: the
linear-kernel GG-KM shows concentrated error on the quadratic
and periodic surfaces, and its estimated $p_c=0.50$ level is
absent on the quadratic surface because the fitted cure surface
does not reach that level. The polynomial-kernel GG-KM
similarly deteriorates on the periodic surface. In contrast,
the RBF GG-KM maintains comparatively stable recovery across
the three generating mechanisms relative to the linear and
polynomial kernels, with its estimated $p_c=0.50$ contour
closely following the true level wherever it is present.
Figure~\ref{fig:cure-surface-recovery} therefore complements
the regret summary in Table~\ref{tab:simulation-regret} by
showing where prediction errors occur spatially: no single
predictor dominates across all three surfaces, with the
linear kernel performing best on the linear surface, the RBF
kernel on the quadratic surface, and the additive spline on
the periodic surface.

\GGKMWideFigure{simulation_cure_surface_recovery.pdf}{Pointwise accuracy of the three primary cure surfaces at $n=600$ under Poisson generation. Rows show the linear, quadratic and periodic generating surfaces; columns show the linear-kernel GG-KM, additive-spline, polynomial-kernel GG-KM and radial-basis-function (RBF) GG-KM predictors. The filled field is the root mean squared error of the predicted cure probability at each of the $41\times41$ grid points, taken across all valid replications, on one scale shared by every panel and truncated above 0.30. Solid labelled lines are the true cure levels $0.25$, $0.50$ and $0.60$; the dashed line is the mean estimated $0.50$ level and is annotated as absent when the fitted range does not contain it. The value printed in each panel is the mean of the displayed field. Failed fits remain in $n_{\mathrm{total}}$ and reduce the pointwise $n_{\mathrm{valid}}$ rather than disappearing.}{fig:cure-surface-recovery}{A three-row by four-column grid of shaded panels shows the pointwise root mean squared error of the predicted cure probability for the linear, quadratic and periodic generating surfaces. Columns correspond to the linear-kernel GG-KM, additive-spline, polynomial-kernel GG-KM and RBF GG-KM predictors. Darker shading marks larger error, and the printed value in each panel is the mean of the displayed field. Solid labelled lines mark the true cure levels, while a dashed line marks the mean estimated 0.50 level and is explicitly reported as absent wherever that level is outside the fitted range.}

The complete mean IPCW integrated Brier scores for all eight methods,
with their Monte Carlo standard errors, appear in the Supporting
Information. They confirm the geometry-specific ordering in
Table~\ref{tab:simulation-best} and retain both fitted count laws for
every predictor.

To complement the scenario-specific comparisons, we
summarized performance using regret relative to the
best-performing method within the retained comparison set.
Table~\ref{tab:simulation-regret} restricts this analysis to
the Poisson and negative-binomial spline models and RBF GG-KM
models across the 18 scenarios. For IPCW integrated Brier score,
the Poisson spline had the smallest maximum paired regret,
0.00065 with a 95\% Monte Carlo half-width of 0.00056, whereas
the negative-binomial spline had the smallest equally weighted
mean regret, 0.00012 (half-width 0.00010). For cure root mean
squared error, the Poisson spline had the smallest maximum regret,
0.00464 (half-width 0.00210), whereas the negative-binomial spline
had the smallest mean regret, 0.00073 (half-width 0.00025). In
contrast, the negative-binomial RBF GG-KM was most competitive
for cure AUC, attaining the smallest maximum regret, 0.00687
(half-width 0.00464), the smallest mean regret, 0.00144
(half-width 0.00060), and the best descriptive mean rank of
2.06. Thus, the regret summaries favor the spline models for
prediction error but show greater competitiveness of the RBF
GG-KM for cure discrimination.

\GGKMTable{sections/generated/simulation-regret-table}

\subsection{Higher-dimensional interaction surfaces}\label{sec:highdim-simulation}

The four-covariate experiment is a secondary mechanism study developed
after the primary analysis. It probes interaction structures that
an additive spline cannot represent, but it is neither
independent confirmation nor evidence of a general kernel
advantage. Covariates satisfy
$x_1,\dots,x_4\sim\operatorname{Uniform}(-1,1)$, and the first
target is the \emph{interaction} surface
\begin{equation}
p_c(\bm{x})=\operatorname{expit}\bigl(C_{\mathrm{int}}+1.60\,x_1x_2-1.30\,x_3x_4+0.70\,x_1-0.50\,x_4\bigr),
\end{equation}
and the \emph{radial} surface is
\begin{equation}
p_c(\bm{x})=\operatorname{expit}\bigl(C_{\mathrm{rad}}+3.00\exp\{-0.80(x_1^2+x_2^2+x_3^2+x_4^2)\}\bigr).
\end{equation}
Intercepts $C_{\mathrm{int}}$ and $C_{\mathrm{rad}}$ were
calibrated by bisection on a fixed sample of $200\,000$
observations (seed 20260901) so that each surface yields a mean
cure fraction of $0.45$. The latency remained the common
generalized gamma with $(a,d,p)=(1.6,1.3,1.1)$ and no
covariate-dependent scale; censoring was exponential with mean
3.5. The same fitted methods as in the primary experiment were
compared under Poisson, negative-binomial ($\phi=0.25$), and
negative-binomial ($\phi=1.00$) generating laws at $n=300$ and
$n=600$, giving 12 scenarios. The final experiment used 250
replicates per scenario, and every declared cell remains in the
analysis regardless of its ordering.

Table~\ref{tab:highdim-best} and
Figure~\ref{fig:highdim-paired-differences} report the directed
comparisons with paired Monte Carlo uncertainty. Their
interpretation is limited to the evaluated interaction and
radial mechanisms, and it retains every declared cell,
including comparisons that do not favor GG-KM.

    \GGKMTable{sections/generated/highdim-best-table}
\GGKMWideFigure
{highdim_rbf_vs_linear_spline_polynomial_paired_differences_with_auc.pdf}
{Paired predictive differences of the radial-basis GG-KM predictor against the count-matched GG-KM linear-kernel, spline, and GG-KM polynomial-kernel comparators on the interaction and radial cure surfaces. The top row shows paired mean IPCW integrated Brier score differences and the bottom row paired mean cure root mean squared error differences. Each plotted contrast is defined as the RBF GG-KM value minus the corresponding comparator value, so negative differences favor the RBF GG-KM and positive differences favor the comparator. Solid lines with filled markers denote Poisson fits, whereas dashed lines with open markers denote negative-binomial fits. Each panel displays the six generating-law and sample-size cells, with error bars spanning $\pm1.96$ paired Monte Carlo standard errors.}
{fig:highdim-paired-differences}
{Two-row by two-column grid of paired-difference plots for the interaction and radial cure surfaces, comparing RBF GG-KM with GG-KM linear-kernel, spline, and GG-KM polynomial-kernel models under Poisson and negative-binomial latent-count specifications.}

The complete per-scenario means with Monte Carlo standard errors, together with the numerical diagnostics, are
reported in the Supporting Information.

\section{Melanoma application}\label{sec:melanoma}

The empirical application used the malignant melanoma cohort distributed
as the \texttt{melanoma} data set in version 1.3-32 of the \texttt{boot}
R package and documented in the original counting-process source.
\cite{boot2025,andersen1993} The cohort contains 205 patients operated
on in Odense, Denmark between 1962 and 1977, with follow-up recorded through
1977. There were 57 deaths from melanoma, 14 deaths from other causes,
and 134 patients alive at the end of follow-up. The analysis treated
death from melanoma as the event of interest and censored both administrative
survivors and deaths from other causes, so the fitted model is cause
specific in its event definition and is not a competing-risks analysis;
$p_c(\bm{x})$ is the probability of no melanoma death from the modeled
process, not an all-cause survival probability.
The predictor set comprised sex, age at operation, calendar year of
operation, the natural logarithm of tumour thickness in millimetres, and
ulceration status. Follow-up time was converted from days to years, and
no observations or covariate values were missing. Continuous and binary
predictors passed through the same training-fold standardization step so
that no test-set information entered preprocessing. The versioned
analysis copy was confirmed row for row against \texttt{boot::melanoma}
and audited against its recorded SHA-256 checksum before estimation.

Predictive performance was assessed in 50 event-stratified outer holdouts,
each assigning 143 patients to training and 62 to testing. The same split
indices were used for every fitted method, and all tuning choices were made
within three inner training folds using the IPCW integrated Brier score
and the horizon rule defined in Section~\ref{sec:prediction-evaluation}.
The comparison comprised GG-KM predictors with linear, radial-basis,
Laplacian, exponential, Cauchy, polynomial, and wavelet kernels, each
with Poisson and negative-binomial latent-count laws, together with
additive cubic-spline Poisson and negative-binomial models and the
classical mixture-cure and promotion-time comparators. Kernel and spline
hyperparameters were selected from prespecified grids.

\GGKMTable{sections/generated/melanoma-performance-table}

Table~\ref{tab:melanoma-performance} reports split-level means and
variation. The Poisson and negative-binomial GG-KM models with the RBF
kernel attained the lowest mean IPCW integrated Brier score, both at
0.1418 with split standard deviation 0.0206, followed closely by the
Cauchy kernel under both count laws, with mean IBS 0.1419. The
mixture-cure model had the highest mean cure AUC overall, 0.8492, while
among the GG-KM models the negative-binomial polynomial kernel had the
highest mean cure AUC, 0.8117. Across the 900 planned split-model fits,
897 completed successfully and 893 satisfied the explicit convergence
criteria.

Figure~\ref{fig:melanoma-performance} places the paired score
distributions beside descriptive survival calibration. The comparison of
split-averaged predicted survival with the Kaplan--Meier curve is a
calibration check on the population survival function; it does not
identify individual cure status. Repeated holdouts are dependent because
the same patients recur across splits, so empirical quantiles describe
variation across training sets rather than independent-sample confidence
intervals. Similar mean scores do not establish an approximately linear
latent-cause structure, and a smaller split standard deviation alone does
not prove greater predictive stability.

\GGKMWideFigure{melanoma_predictive_assessment_ibs_auc.pdf}{Predictive assessment in 50 paired repeated holdouts of the melanoma cohort. Panel A compares the full-cohort Kaplan--Meier curve with split-averaged out-of-sample survival predictions over the common supported horizon. Panel B shows split-level IPCW integrated Brier scores, with each point representing one outer split and lower values indicating better prediction. Line type, marker or box encoding, and labels distinguish all eight methods without relying on color alone. Panel A shows mean curves without an uncertainty band; the dependence caveat for repeated holdouts is stated in the text.}{fig:melanoma-performance}{Two-panel melanoma prediction figure. The left panel shows a black Kaplan-Meier step curve and eight mean predicted survival curves from log-linear, spline, quadratic, and GG-KM models with Poisson or negative-binomial count laws. The right panel shows split-level IPCW integrated Brier scores as box-and-strip distributions for every method, with each point representing one outer split and lower values indicating better prediction.}
\GGKMWideFigure{melanoma_ggkm_kernel_comparison_ibs_auc.pdf}{Kernel-specific predictive performance of the GG-KM models in 50 paired repeated holdouts of the melanoma cohort. Panel A shows the split-level IPCW integrated Brier score and Panel B the cure-discrimination area under the curve ($\mathrm{AUC}_{\mathrm{cure}}$) across the seven candidate kernels. Within each kernel, Poisson and negative-binomial GG-KM results are shown side by side, with color and hatching distinguishing the two latent-count laws. Each point represents one outer split, and the boxplots summarize the corresponding split-level distributions.}{fig:melanoma-performance-ggkm-kernel}{Two-panel comparison of GG-KM kernels in the melanoma application. Panel A displays split-level IPCW integrated Brier scores and Panel B cure AUC values for the linear, RBF, Laplacian, exponential, Cauchy, polynomial, and wavelet kernels. For each kernel, Poisson and negative-binomial GG-KM boxplots are shown side by side and distinguished by color and hatching. Boxplots summarize the repeated-holdout distributions, with individual outer-split values overlaid as points.}

Mean predicted cure fractions ranged from 0.572 for the Poisson
GG-KM with the wavelet kernel to 0.639 for the Poisson GG-KM with the
polynomial kernel. These values are model-based extrapolations because
individual cure status is unobserved and follow-up is finite, and the
Kaplan--Meier comparison cannot identify a cure fraction. Across-split
prediction summaries were retained for every individual to expose
variation caused by repeated training samples rather than to create
surrogate cure labels. The application consequently evaluates
reproducible out-of-sample prediction and numerical behavior while
leaving biological interpretation of individual cure probabilities
outside its evidential scope.

Figure~\ref{fig:melanoma-performance-ggkm-kernel} first indicates a marked
similarity between the Poisson and negative-binomial GG-KM specifications.
For most kernels, the corresponding boxplots have strongly overlapping
distributions for both the IPCW integrated Brier score and
$\mathrm{AUC}_{\mathrm{cure}}$, and their mean performances are generally
close.
% This pattern provides empirical evidence, in this application, that
% the predictive behavior of the GG-KM model is relatively insensitive to the
% specific latent-count distribution used. It should not, however, be
% interpreted as establishing distributional independence of the model, since
% the comparison is limited to the Poisson and negative-binomial
% specifications and to this data set. Rather, the results suggest a degree of
% robustness to the latent-count specification, with the kernel choice having
% a more visible influence on predictive performance.

For the IPCW integrated Brier score (Panel A), the distributions overlap
substantially across kernels, and the differences between the Poisson and
negative-binomial versions of the same kernel are generally small. Mean IBS
values were relatively close across specifications, with both the Poisson
and negative-binomial RBF models attaining the lowest mean value, 0.1418.
The Poisson polynomial model had the largest mean IBS, 0.1515, and also
showed comparatively large between-split variability. Although several
kernels exhibited isolated splits with substantially higher Brier scores,
the central parts of the distributions remained similar. Overall, the IBS
results do not indicate a clearly dominant kernel.

Panel B shows greater sensitivity to the kernel choice for
$\mathrm{AUC}_{\mathrm{cure}}$, while still retaining substantial agreement
between the two latent-count laws within most kernels. The linear kernel
provided consistently strong cure discrimination under both count laws,
with mean cure AUCs of approximately 0.809. The negative-binomial
polynomial model achieved the highest mean cure AUC among the GG-KM
specifications, 0.8117, and showed a relatively concentrated distribution,
whereas the corresponding Poisson polynomial model was substantially more
variable across splits. The RBF, Laplacian, exponential, and Cauchy kernels
generally produced lower cure AUCs, concentrated around 0.68--0.70, while
the wavelet kernel showed the weakest discrimination, with mean cure AUCs
of approximately 0.62 under both count laws. Thus, the results suggest that
kernel specification has a stronger practical influence than the choice
between Poisson and negative-binomial latent-count laws in this application.
No single kernel dominated simultaneously in prediction error and cure
discrimination: the RBF kernel was most favorable in terms of mean IBS
under both count laws, whereas the polynomial and linear kernels provided
the strongest cure discrimination.

\section{Breast Cancer Data Application}\label{sec:cancer-data-application}

The second empirical application considered a hospital-based cohort of
patients with breast cancer from Paraíba State, Brazil. Follow-up time was
represented by the recorded survival time, and death from breast cancer was
treated as the event of interest; observations without a recorded
breast-cancer death were treated as censored. In contrast with the melanoma application, the breast-cancer data
contained a substantially richer set of demographic, socioeconomic,
clinical, health-care access, lifestyle, and treatment-related information,
allowing the cure component to depend on a broader patient profile.

The predictor set included race or skin-colour classification, educational
level, income category, age at first consultation, age group, health
macroregion, health-system origin and referral pathway, and the hospital
source of the record. Clinical characteristics included primary tumour
site, histological type, and clinical stage, with stage grouped into
stages I--II, stages III--IV, and a separate category for unavailable
staging information. Lifestyle and history variables comprised tobacco
use, alcohol use, and family history of cancer. Information on the first
treatment received at the hospital was represented through separate
indicators for surgery, radiotherapy, chemotherapy, hormone therapy,
immunotherapy, no treatment, other treatments, and unavailable treatment
information. Health-care timing was represented by the intervals from
first consultation to diagnosis, diagnosis to first treatment, and first
consultation to first treatment. In addition to their continuous versions,
these intervals were represented by indicators distinguishing delays of up
to 90 days from delays exceeding 90 days.

Preprocessing was estimated exclusively from each training sample and then
applied unchanged to the corresponding validation or test observations.
Continuous age and delay variables were transformed using the Yeo--Johnson
power transformation\cite{yeo2000new} and standardized. Explicit missingness indicators were
retained for treatment-delay variables, while several categorical predictors
included dedicated levels for unavailable information. Because primary
tumour site and histological type contained more complex categorical
structures, they were represented using cross-validated survival-target
encodings based on restricted mean survival time, with smoothing applied
within the training data. The resulting representations were carried
unchanged to held-out observations, thereby preventing outcome information
from the validation or test samples from entering the preprocessing stage.

Several prespecified interaction terms were also included to allow clinically
relevant associations to vary across patient subgroups. These comprised
interactions between clinical stage and first treatment, between stage and
hospital source, and between stage and educational level. The continuous
consultation-to-diagnosis delay was additionally interacted with stage and
with hormone-therapy and immunotherapy indicators. These interaction terms,
together with the original transformed continuous variables, survival-target
encodings, and categorical indicators, formed the predictor matrix used by
all competing models.

Predictive performance was assessed using 50 event-stratified paired outer
holdouts, with approximately 70\% of the cohort assigned to training and
30\% to testing in each split. The same outer split indices were used for
all fitted methods. Hyperparameter selection was performed exclusively
within the outer training sample using four inner folds and the IPCW
integrated Brier score, together with the prediction-horizon rule defined
in Section~\ref{sec:prediction-evaluation}. The comparison comprised GG-KM
models under Poisson and negative-binomial latent-count laws, with linear,
radial-basis, Laplacian, exponential, Cauchy, polynomial, and
wavelet kernels, as well as spline Poisson and negative-binomial
models and the classical mixture-cure and promotion-time comparators.
All kernel and spline hyperparameters were selected from prespecified
grids.

\GGKMTable{sections/generated/breast-cancer-performance-table}

Table~\ref{tab:breast-cancer-performance} reports split-level means and
variation. The spline models under both Poisson and negative-binomial
latent-count laws achieved the lowest mean IPCW integrated Brier score,
0.0877 with split standard deviation 0.0019, followed by the mixture-cure
and promotion-time models, both with mean IBS 0.0883. The mixture-cure
model attained the highest mean cure AUC overall, 0.7851, followed by the
promotion-time model at 0.7775 and the spline models at 0.7707. Among the
GG-KM specifications, the negative-binomial linear kernel had the lowest
mean IBS, 0.0897, whereas the exponential kernel produced the highest mean
cure AUC, 0.7124, under both latent-count laws. Across the 900 planned
split-model fits, 873 completed successfully and 871 satisfied the explicit
convergence criteria.

Figure~\ref{fig:breast-cancer-performance} summarizes out-of-sample survival
calibration, prediction error, and cure discrimination. The split-averaged
predicted survival curves closely followed the Kaplan--Meier estimate over
the common prediction horizon, with relatively small differences among the
fitted methods. The spline Poisson and negative-binomial models showed the
lowest IPCW integrated Brier scores, whereas the mixture-cure model achieved
the highest cure-discrimination performance. The GG-KM models with the
exponential kernel were less competitive in overall prediction error than
the spline and classical cure models, but retained substantially better cure
discrimination than several alternative kernel specifications.

\GGKMWideFigure{breast_cancer_predictive_assessment_ibs_auc.pdf}
{Predictive assessment in 50 paired repeated holdouts of the breast cancer
cohort. Panel A compares the full-cohort Kaplan--Meier curve with
split-averaged out-of-sample survival predictions over the common supported
horizon. Panel B shows split-level IPCW integrated Brier scores and Panel C
shows the corresponding cure-discrimination area under the curve
($\mathrm{AUC}_{\mathrm{cure}}$). Each point represents one outer split, and
lower IBS values indicate better prediction whereas larger cure AUC values
indicate better discrimination.}
{fig:breast-cancer-performance}
{Three-panel predictive assessment for the breast cancer application. Panel A
compares the Kaplan--Meier survival curve with mean out-of-sample survival
predictions from the mixture-cure, promotion-time, spline, and selected GG-KM
models. Panels B and C show the corresponding repeated-holdout distributions
of IPCW integrated Brier scores and cure AUC values.}

Figure~\ref{fig:breast-cancer-performance-ggkm-kernel} shows that kernel
choice had a marked influence on GG-KM performance. For prediction error,
the exponential and Cauchy kernels yielded the lowest mean IBS values among
the nonlinear GG-KM specifications, while the wavelet and RBF kernels were
less favorable. Cure discrimination varied more substantially across
kernels. The exponential kernel produced the highest mean cure AUC among the
GG-KM models, 0.7124 under both count laws, whereas the wavelet kernel showed
considerably weaker discrimination.

% The linear Poisson model was particularly
% poor for cure discrimination, with mean AUC close to 0.5, while its
% negative-binomial counterpart performed substantially better. For most of
% the RBF, Laplacian, exponential, Cauchy, and wavelet specifications, however,
% the Poisson and negative-binomial results were very similar.

\GGKMWideFigure{breast_cancer_ggkm_kernel_comparison_ibs_auc.pdf}
{Kernel-specific predictive performance of the GG-KM models in 50 paired
repeated holdouts of the breast cancer cohort. Panel A shows the split-level
IPCW integrated Brier score and Panel B the cure-discrimination area under
the curve ($\mathrm{AUC}_{\mathrm{cure}}$) across the seven candidate
kernels. Within each kernel, Poisson and negative-binomial GG-KM results are
shown side by side, with color and hatching distinguishing the two
latent-count laws. Individual points represent outer splits and boxplots
summarize their distributions.}
{fig:breast-cancer-performance-ggkm-kernel}
{Two-panel comparison of GG-KM kernels in the breast cancer application.
Panel A displays split-level IPCW integrated Brier scores and Panel B cure
AUC values for the linear, RBF, Laplacian, exponential, Cauchy, polynomial,
and wavelet kernels. Poisson and negative-binomial specifications are shown
side by side for each kernel.}

% Discussion, limitations, and proportionate conclusions.
\section{Discussion}\label{sec:discussion}

The generalized gamma kernel-machine promotion-time cure model extends the promotion-time
framework by placing the latent-cause log-intensity in a reproducing-kernel Hilbert space
while retaining a flexible generalized gamma latency and both Poisson and negative-binomial
latent counts. The primary and secondary simulations and the melanoma and breast-cancer
applications evaluate this construction against simpler predictors under a shared protocol.
Their interpretation is limited to the evaluated designs. Kernel families were compared
within prespecified grids, although robustness to broader kernel and hyperparameter choices
remains for future work.

The principal simulation pattern was consistent across the two latent-count laws. On the
quadratic surface, the RBF GG-KM lowered the paired IPCW integrated Brier score against the
linear-kernel GG-KM by about 0.007 to 0.012 and the paired cure root mean squared error by
about 0.04 to 0.07, while its differences from the spline were much smaller and
scenario-dependent. On the periodic surface, the RBF GG-KM improved substantially on the
linear-kernel GG-KM yet trailed the spline by a small margin, and on the linear surface it
paid a small resolved penalty against the correctly specified linear-kernel GG-KM. Within
the original three-surface design, the count law was a secondary factor, since Poisson and
negative-binomial variants within a predictor family differed by at most 0.02 in integrated
Brier score and 0.01 in cure root mean squared error. The melanoma cohort produced overlapping
score distributions in which the linear-kernel Poisson model had the lowest mean integrated
Brier score and GG-KM was competitive with the smallest split-to-split variation. The
breast-cancer cohort showed a complementary pattern: the spline models achieved the lowest
mean integrated Brier score, 0.0877, whereas the mixture-cure model attained the
highest mean cure AUC, 0.7851. Among the GG-KM specifications, the negative-binomial
linear kernel had the lowest mean integrated Brier score, 0.0897, while the exponential
kernel attained the highest mean cure AUC, 0.7124, under both latent-count laws. Together,
the two empirical applications reinforce the simulation message that no
single predictor is uniformly best and that the preferred predictor and
kernel depend on the predictive target and data structure.

Several features of the cure-model setting limit how far these results can be interpreted. Cure fractions are
only weakly identified from right-censored data with finite follow-up, because a long plateau of censored
survivors is consistent with a range of cure probabilities. The melanoma follow-up ends in 1977 and the
estimated cure fractions between 0.577 and 0.648 are model-based extrapolations, not quantities read directly
from the data. The comparison of split-averaged predicted survival with the Kaplan-Meier curve is a
calibration check on the population survival function and does not identify which individuals are cured. The
analysis therefore makes no claim to identify individual cure status, and simulation calibration slopes near
identity do not change that limitation.

Kernel choice and its hyperparameters shape the estimated intensity because they define which subjects are similar
and how rapidly the fitted surface may change. In the primary three-surface comparisons, multiple GG-KM kernel
families were evaluated, with their hyperparameters selected from prespecified grids by leakage-free inner tuning;
the high-dimensional analysis likewise compared multiple kernel families. Bandwidths that are too small can track
idiosyncratic training variation, whereas overly large bandwidths approach a much smoother response; both effects
depend on the training-fold scaling used to measure distance. The high-dimensional results further illustrate that
kernel family can materially affect performance, with polynomial GG-KM performing particularly well on the
interaction surface. A different kernel family or hyperparameter range could therefore shift the balance between
GG-KM and the spline, particularly on surfaces whose geometry does not align well with a given kernel, as the
periodic surface illustrates. Kernel and penalty selection are model-selection steps, and uncertainty from that
selection is not represented by the fitted curve alone. The comparisons fixed a common tuning budget and evaluation
protocol for all methods, so the reported differences describe performance at a shared and modest search intensity
rather than each method's best attainable configuration.

The kernel-machine estimator also carries a computational cost that grows with sample size. An
$n_{\mathrm{train}}\times n_{\mathrm{train}}$ Gram matrix requires quadratic storage, and dense
factorization can require cubic time; nested tuning repeats those operations across candidate kernels and
inner folds. Reusing a Gram matrix within a candidate and retained eigenspace avoids unnecessary work, but
matrices cannot be shared across folds whose preprocessing and training subjects differ. In the primary
simulation experiment, these costs remained manageable at the sample sizes studied, with mean per-fit
runtimes across methods ranging from 1.13 to 40.32 seconds. The richer breast-cancer application was
substantially more computationally demanding, particularly for several nonlinear GG-KM kernels and the
spline models, underscoring that computational burden can increase sharply with a larger and more complex
predictor representation; these timings do not establish scalability to substantially larger cohorts.
Low-rank, Nystr\"om, random-feature, or inducing-point approximations could reduce the burden, although each
changes the effective penalty and out-of-sample predictor and would require separate numerical validation.
Sample size affected the strength of the conclusions as well as runtime: increasing the sample from 300 to
600 reduced Monte Carlo uncertainty and resolved the quadratic and periodic contrasts more clearly.

Censoring and the event definition further bound the interpretation. The evaluation used
inverse-probability-of-censoring weighting with a horizon confined to the region of reliable censoring
support, which controls bias from dependent censoring only to the extent that the censoring model is
adequate. In the melanoma cohort, deaths from other causes were treated as censored. In the breast-cancer
cohort, death from breast cancer was treated as the event, while observations without a recorded
breast-cancer death were censored. Accordingly, both empirical analyses are cause specific in their event
definition and are not competing-risks analyses. Within these latent cause-specific processes,
$p_c(\bm{x})$ describes the probability of no disease-specific death from the modelled process; it is
neither an all-cause survival probability nor a competing-risks cumulative incidence. The interpretation
relies on the corresponding independent-censoring assumption. The negative-binomial dispersion was only
weakly identified from cure and latency information: under Poisson generation the estimated dispersion
correctly stayed near zero, but under true dispersion of 1.00 it was systematically underestimated, ranging
from 0.053 to 0.239. This weak identifiability explains the frequent contact of the dispersion estimate with
its lower bound and cautions against interpreting the fitted dispersion as a precise measure of
overdispersion.

Both empirical applications are observational, and no causal interpretation is attached to the covariate effects
or to the predicted cure fractions. The reported associations describe out-of-sample predictive behaviour under
a fixed analysis protocol, not the effect of intervening on any covariate.

Future work could extend the estimator in directions that were deliberately left outside the present scope.
A controlled comparison against reproducible support-vector or deep cure models would sharpen the
positioning of kernel machines among nonlinear approaches, but a fair comparison requires public
implementations of the same promotion-time model with common generalized gamma latency, which were not
available, so such comparators are noted as open rather than completed. Covariate-dependent latency,
low-rank Gram approximations for larger cohorts, adaptive or multiple-kernel methods, broader kernel
hyperparameter searches, and joint modelling of the censoring mechanism are natural extensions. The current
contribution is the validated combination of an RKHS latent-cause intensity, generalized gamma latency,
Poisson and negative-binomial latent counts, penalized EM and ECME estimation, and reproducible
leakage-free evaluation with paired Monte Carlo uncertainty.

\section{Declarations}\label{sec:declarations}

\textbf{Data and code availability.} The malignant melanoma data are distributed as \texttt{boot::melanoma} in version 1.3-32 of the \texttt{boot} R package, whose official CRAN metadata report an Unlimited license, and are cited to the package documentation and original source. The reproducible article workspace contains the row-for-row verified analysis copy, its SHA-256 checksum, executable Python code, frozen design records, session information, checkpoints, and machine-readable outputs for every table and figure. \textbf{Author confirmation required:} deposit the final data and code package in a public archive, insert its persistent identifier, and verify that the archived files reproduce the submitted version.

\textbf{Ethics and consent.} This study reanalyzes a publicly distributed historical cohort and does not report newly collected participant data. The available thesis and package documentation do not establish whether a new institutional determination is required for this secondary analysis. \textbf{Author confirmation required:} provide the journal-appropriate ethics approval, exemption, or not-applicable statement and confirm whether a consent statement is required.

\textbf{Funding, competing interests, and acknowledgments.} This study was supported by the Espírito Santo Research and Innovation Support Foundation – FAPES (grant number FAPES nº 1432/2025 P: 2025-R8KXL).
{Research carried out using the computational resources of the Center for
Mathematical Sciences Applied to Industry (CeMEAI) funded by FAPESP
(grant 2013/07375-0).} The authors declare no competing interests.
\textbf{Author confirmation required:} supply any additional funding
agencies and grant numbers, if applicable, and approve the final
acknowledgment and competing-interest statements.

\textbf{Author contributions.} The authorship metadata, namely contributor identities, order, affiliations, and ORCID identifiers, were supplied by the corresponding author, Vera L. D. Tomazella. No contributor has been added, removed, or reordered during preparation of the article workspace.

\textbf{Use of generative artificial intelligence.} Generative artificial-intelligence tools were used during preparation and revision of this article for language editing, code drafting, and consistency checks. Scientific assumptions, analysis choices, citations, executable outputs, and final wording remain subject to author verification, and the confirmed authors retain full responsibility for the submitted work. \textbf{Author confirmation required:} confirm the specific tool identities and dates of use, verify this disclosure against current Wiley policy, add any other qualifying tools or uses, and approve the final wording.


\bibliography{references}

\ifggkmsubmission
\clearpage
\begin{table*}[tbp]
\caption{Predictive performance in 50 paired repeated holdouts of the melanoma cohort.
Values for the IPCW integrated Brier score (IBS), cure-discrimination area under the curve
($\mathrm{AUC}_{\mathrm{cure}}$), and mean predicted cure fraction are means across successful
splits, with split standard deviations in parentheses. Convergence is reported as the number
of converged fits among successful fits. Runtime is the mean with the maximum in parentheses.
}
\label{tab:melanoma-performance}
\centering
\begin{tabular}{llccccc}
\hline
Predictor & Count law & Converged & IPCW IBS &
$\mathrm{AUC}_{\mathrm{cure}}$ & Mean cure prediction & Runtime (s) \\
\hline

Mixture cure
& Mixture cure
& 49/50
& 0.1430 (0.0185)
& 0.8492 (0.0491)
& 0.607 (0.079)
& 2.28 (12.39) \\

Promotion time
& Poisson
& 49/49
& 0.1456 (0.0239)
& 0.8240 (0.0283)
& 0.633 (0.047)
& 2.21 (7.29) \\

\hline

Spline
& Poisson
& 50/50
& 0.1432 (0.0217)
& 0.7912 (0.0290)
& 0.613 (0.049)
& 42.56 (99.65) \\

Spline
& Negative binomial
& 50/50
& 0.1434 (0.0216)
& 0.7911 (0.0291)
& 0.614 (0.049)
& 46.37 (112.34) \\

\hline

GG-KM (linear)
& Poisson
& 50/50
& 0.1426 (0.0224)
& 0.8091 (0.0274)
& 0.627 (0.046)
& 15.48 (37.21) \\

GG-KM (linear)
& Negative binomial
& 50/50
& 0.1426 (0.0224)
& 0.8091 (0.0274)
& 0.627 (0.046)
& 16.29 (38.53) \\

GG-KM (RBF)
& Poisson
& 50/50
& 0.1418 (0.0206)
& 0.6937 (0.0340)
& 0.586 (0.047)
& 45.85 (97.33) \\

GG-KM (RBF)
& Negative binomial
& 50/50
& 0.1418 (0.0206)
& 0.6937 (0.0341)
& 0.586 (0.047)
& 54.00 (184.25) \\

GG-KM (Laplacian)
& Poisson
& 48/49
& 0.1426 (0.0253)
& 0.6987 (0.0490)
& 0.578 (0.050)
& 47.26 (97.01) \\

GG-KM (Laplacian)
& Negative binomial
& 50/50
& 0.1438 (0.0257)
& 0.7021 (0.0526)
& 0.575 (0.062)
& 55.12 (140.77) \\

GG-KM (exponential)
& Poisson
& 49/50
& 0.1450 (0.0242)
& 0.6984 (0.0800)
& 0.580 (0.048)
& 43.01 (89.43) \\

GG-KM (exponential)
& Negative binomial
& 50/50
& 0.1448 (0.0241)
& 0.6949 (0.0768)
& 0.580 (0.049)
& 49.45 (110.53) \\

GG-KM (Cauchy)
& Poisson
& 50/50
& 0.1419 (0.0246)
& 0.6817 (0.0390)
& 0.585 (0.047)
& 46.38 (92.49) \\

GG-KM (Cauchy)
& Negative binomial
& 50/50
& 0.1419 (0.0246)
& 0.6823 (0.0390)
& 0.586 (0.046)
& 51.95 (173.17) \\

GG-KM (polynomial)
& Poisson
& 50/50
& 0.1515 (0.0274)
& 0.7457 (0.1338)
& 0.639 (0.063)
& 113.78 (232.80) \\

GG-KM (polynomial)
& Negative binomial
& 50/50
& 0.1477 (0.0280)
& 0.8117 (0.0261)
& 0.613 (0.054)
& 368.69 (806.30) \\

GG-KM (wavelet)
& Poisson
& 48/49
& 0.1488 (0.0137)
& 0.6176 (0.0337)
& 0.572 (0.061)
& 64.57 (146.75) \\

GG-KM (wavelet)
& Negative binomial
& 50/50
& 0.1482 (0.0131)
& 0.6193 (0.0320)
& 0.573 (0.060)
& 77.38 (219.31) \\

\hline
\end{tabular}
\end{table*}
%
\begin{table*}[tbp]
\caption{Predictive performance in 50 paired repeated holdouts of the breast cancer cohort
from Paraíba State, Brazil. Values for the IPCW integrated Brier score (IBS),
cure-discrimination area under the curve ($\mathrm{AUC}_{\mathrm{cure}}$), and mean
predicted cure fraction are means across successful splits, with split standard deviations
in parentheses. Convergence is reported as the number of converged fits among successful
fits. Runtime is the mean with the maximum in parentheses.
}
\label{tab:breast-cancer-performance}
\centering
\begin{tabular}{llccccc}
\hline
Predictor & Count law & Converged & IPCW IBS &
$\mathrm{AUC}_{\mathrm{cure}}$ & Mean cure prediction & Runtime (s) \\
\hline

Mixture cure
& Mixture cure
& 48/50
& 0.0883 (0.0019)
& 0.7851 (0.0156)
& 0.814 (0.029)
& 34.04 (132.92) \\

Promotion time
& Poisson
& 50/50
& 0.0883 (0.0022)
& 0.7775 (0.0072)
& 0.816 (0.009)
& 38.06 (71.25) \\

\hline

Spline
& Poisson
& 50/50
& 0.0877 (0.0019)
& 0.7707 (0.0067)
& 0.816 (0.009)
& 1402.30 (2462.32) \\

Spline
& Negative binomial
& 50/50
& 0.0877 (0.0019)
& 0.7707 (0.0067)
& 0.816 (0.009)
& 1593.56 (4530.27) \\

\hline

GG-KM (linear)
& Poisson
& 50/50
& 0.0956 (0.0008)
& 0.5007 (0.0015)
& 0.884 (0.000)
& 137.77 (169.54) \\

GG-KM (linear)
& Negative binomial
& 50/50
& 0.0897 (0.0017)
& 0.7085 (0.0103)
& 0.818 (0.008)
& 1934.49 (7967.88) \\

GG-KM (RBF)
& Poisson
& 50/50
& 0.0927 (0.0020)
& 0.6262 (0.0699)
& 0.807 (0.011)
& 6648.94 (11951.68) \\

GG-KM (RBF)
& Negative binomial
& 50/50
& 0.0927 (0.0020)
& 0.6262 (0.0699)
& 0.807 (0.011)
& 5364.04 (11483.93) \\

GG-KM (Laplacian)
& Poisson
& 50/50
& 0.0921 (0.0021)
& 0.6745 (0.0557)
& 0.801 (0.012)
& 5410.98 (9729.17) \\

GG-KM (Laplacian)
& Negative binomial
& 50/50
& 0.0921 (0.0021)
& 0.6745 (0.0557)
& 0.801 (0.012)
& 5701.58 (12885.37) \\

GG-KM (exponential)
& Poisson
& 50/50
& 0.0901 (0.0022)
& 0.7124 (0.0579)
& 0.804 (0.013)
& 11443.53 (23233.80) \\

GG-KM (exponential)
& Negative binomial
& 50/50
& 0.0901 (0.0022)
& 0.7124 (0.0579)
& 0.804 (0.013)
& 8428.04 (17262.88) \\

GG-KM (Cauchy)
& Poisson
& 50/50
& 0.0899 (0.0019)
& 0.6927 (0.0240)
& 0.798 (0.011)
& 12069.98 (25675.73) \\

GG-KM (Cauchy)
& Negative binomial
& 50/50
& 0.0899 (0.0019)
& 0.6927 (0.0240)
& 0.798 (0.011)
& 8726.52 (19034.95) \\

% GG-KM (sigmoid)
% & Poisson
% & 49/49
% & 0.0945 (0.0008)
% & 0.5000 (0.0000)
% & 0.819 (0.008)
% & 1974.30 (4980.47) \\

% GG-KM (sigmoid)
% & Negative binomial
% & 40/40
% & 0.0947 (0.0008)
% & 0.5000 (0.0000)
% & 0.818 (0.007)
% & 1680.01 (3205.18) \\

GG-KM (polynomial)
& Poisson
& 37/37
& 0.0927 (0.0013)
& 0.6721 (0.0375)
& 0.885 (0.004)
& 630.28 (717.68) \\

GG-KM (polynomial)
& Negative binomial
& 36/36
& 0.0902 (0.0015)
& 0.7019 (0.0090)
& 0.820 (0.007)
& 2973.05 (6276.62) \\

GG-KM (wavelet)
& Poisson
& 50/50
& 0.0944 (0.0011)
& 0.5456 (0.0274)
& 0.807 (0.012)
& 5598.89 (10148.10) \\

GG-KM (wavelet)
& Negative binomial
& 50/50
& 0.0944 (0.0011)
& 0.5456 (0.0274)
& 0.807 (0.012)
& 5417.58 (13914.95) \\

\hline
\end{tabular}
\end{table*}
%
\begin{table*}[tbp]
\caption{
Best observed mean predictive performance in each simulation scenario
among Poisson and negative-binomial models.
For IPCW IBS and cure RMSE, smaller values are better, whereas for cure AUC,
larger values are better. Parenthesized values are Monte Carlo standard errors.
The displayed optimum is a descriptive comparison and is not a claim of
superiority when paired uncertainty includes zero.
GG, Spl, P, and NB denote generalized-gamma kernel machine, spline,
Poisson, and negative binomial, respectively.
For GG-KM fits, the selected kernel is shown in parentheses.
}
\label{tab:simulation-best}
\centering
\small
\setlength{\tabcolsep}{4pt}
\begin{tabular}{lllccc}
\hline
True surface
& Generating law
& $n$
& Lowest mean IPCW IBS
& Lowest mean cure RMSE
& Highest mean cure AUC \\
\hline

Linear
& P
& 300
& GG-NB (linear): 0.1929 (0.0013)
& GG-P (linear): 0.0630 (0.0019)
& GG-NB (poly): 0.6845 (0.0025) \\

Linear
& NB(0.25)
& 300
& GG-NB (linear): 0.1965 (0.0015)
& GG-NB (linear): 0.0602 (0.0017)
& GG-NB (poly): 0.6930 (0.0025) \\

Linear
& NB(1.00)
& 300
& GG-NB (linear): 0.1975 (0.0014)
& GG-P (linear): 0.0606 (0.0017)
& GG-NB (poly): 0.7032 (0.0025) \\

Linear
& P
& 600
& GG-NB (linear): 0.1937 (0.0010)
& GG-P (linear): 0.0458 (0.0011)
& GG-NB (poly): 0.6875 (0.0020) \\

Linear
& NB(0.25)
& 600
& GG-NB (linear): 0.1944 (0.0010)
& GG-P (linear): 0.0448 (0.0011)
& GG-NB (poly): 0.6910 (0.0019) \\

Linear
& NB(1.00)
& 600
& GG-NB (linear): 0.1968 (0.0011)
& GG-NB (linear): 0.0453 (0.0012)
& GG-NB (poly): 0.7045 (0.0018) \\

\hline

Quadratic
& P
& 300
& GG-NB (poly): 0.2097 (0.0014)
& GG-NB (poly): 0.0990 (0.0018)
& GG-NB (poly): 0.6411 (0.0029) \\

Quadratic
& NB(0.25)
& 300
& GG-P (rbf): 0.2128 (0.0014)
& GG-NB (poly): 0.0962 (0.0016)
& GG-P (poly): 0.6443 (0.0031) \\

Quadratic
& NB(1.00)
& 300
& GG-P (poly): 0.2145 (0.0016)
& GG-NB (poly): 0.0933 (0.0018)
& GG-P (poly): 0.6586 (0.0031) \\

Quadratic
& P
& 600
& GG-NB (poly): 0.2079 (0.0010)
& GG-NB (poly): 0.0720 (0.0013)
& GG-P (wavelet): 0.6422 (0.0018) \\

Quadratic
& NB(0.25)
& 600
& GG-NB (poly): 0.2096 (0.0009)
& GG-NB (poly): 0.0687 (0.0013)
& GG-P (wavelet): 0.6452 (0.0018) \\

Quadratic
& NB(1.00)
& 600
& GG-NB (poly): 0.2128 (0.0010)
& GG-NB (poly): 0.0646 (0.0012)
& GG-NB (poly): 0.6555 (0.0019) \\

\hline

Periodic
& P
& 300
& Spl-NB: 0.1913 (0.0017)
& Spl-NB: 0.1091 (0.0017)
& GG-NB (poly): 0.7237 (0.0034) \\

Periodic
& NB(0.25)
& 300
& Spl-NB: 0.1906 (0.0015)
& Spl-NB: 0.1050 (0.0017)
& GG-NB (poly): 0.7216 (0.0033) \\

Periodic
& NB(1.00)
& 300
& Spl-NB: 0.1905 (0.0016)
& Spl-NB: 0.1005 (0.0016)
& GG-NB (poly): 0.7345 (0.0033) \\

Periodic
& P
& 600
& Spl-NB: 0.1851 (0.0010)
& Spl-P: 0.0794 (0.0012)
& GG-NB (poly): 0.7301 (0.0023) \\

Periodic
& NB(0.25)
& 600
& Spl-NB: 0.1858 (0.0011)
& Spl-NB: 0.0778 (0.0012)
& GG-NB (poly): 0.7361 (0.0022) \\

Periodic
& NB(1.00)
& 600
& Spl-P: 0.1878 (0.0012)
& Spl-P: 0.0728 (0.0013)
& GG-NB (poly): 0.7501 (0.0020) \\

\hline
\end{tabular}
\end{table*}
%
% Generated from the Poisson/NB regret results,
% restricted to spline models and GG-KM with the RBF kernel.
\begin{table*}[tbp]
\caption{
Paired regret and descriptive rank summaries over the 18 primary simulation
scenarios for the Poisson and negative-binomial models.
The comparison is restricted to spline models and GG-KM models using the RBF
kernel. Within each scenario and metric, the reference is the retained method
with the best finite mean. For IPCW integrated Brier score and cure RMSE,
regret is the method value minus the reference value; for cure AUC, regret is
the reference value minus the method value, so that smaller regret is better
for all metrics. Regret uses only replication pairs for which both methods are
finite. The worst regret identifies its scenario and reports a 95\% Monte Carlo
half-width; mean regret weights scenarios equally and propagates their Monte
Carlo standard errors. Mean rank is based on scenario means and is descriptive.
}
\label{tab:simulation-regret}
\centering
\resizebox{\textwidth}{!}{%
\begin{tabular}{llccc}
\hline
Metric
& Method
& Worst paired regret ($\pm$95\%; scenario)
& Mean paired regret ($\pm$95\%)
& Mean rank \\
\hline

IPCW integrated Brier score
& Spline Negative binomial
& 0.00067 $\pm$ 0.00056 (Quadratic, P, $n=600$)
& 0.00012 $\pm$ 0.00010
& 2.06 \\

IPCW integrated Brier score
& Spline Poisson
& 0.00065 $\pm$ 0.00056 (Quadratic, P, $n=600$)
& 0.00013 $\pm$ 0.00010
& 2.06 \\

IPCW integrated Brier score
& GG-KM Poisson (RBF)
& 0.00348 $\pm$ 0.00119 (Periodic, NB(1.00), $n=300$)
& 0.00108 $\pm$ 0.00017
& 3.22 \\

IPCW integrated Brier score
& GG-KM Negative binomial (RBF)
& 0.00359 $\pm$ 0.00120 (Periodic, NB(1.00), $n=300$)
& 0.00108 $\pm$ 0.00017
& 2.67 \\

\hline

Cure RMSE
& Spline Negative binomial
& 0.00464 $\pm$ 0.00210 (Quadratic, NB(1.00), $n=600$)
& 0.00073 $\pm$ 0.00025
& 1.94 \\

Cure RMSE
& Spline Poisson
& 0.00464 $\pm$ 0.00210 (Quadratic, NB(1.00), $n=600$)
& 0.00076 $\pm$ 0.00025
& 1.94 \\

Cure RMSE
& GG-KM Poisson (RBF)
& 0.01651 $\pm$ 0.00268 (Periodic, NB(1.00), $n=300$)
& 0.00631 $\pm$ 0.00054
& 3.17 \\

Cure RMSE
& GG-KM Negative binomial (RBF)
& 0.01627 $\pm$ 0.00272 (Periodic, NB(1.00), $n=300$)
& 0.00632 $\pm$ 0.00053
& 2.94 \\

\hline

Cure AUC
& Spline Negative binomial
& 0.00768 $\pm$ 0.00529 (Quadratic, P, $n=300$)
& 0.00208 $\pm$ 0.00069
& 2.61 \\

Cure AUC
& Spline Poisson
& 0.00693 $\pm$ 0.00538 (Quadratic, P, $n=300$)
& 0.00224 $\pm$ 0.00070
& 2.94 \\

Cure AUC
& GG-KM Poisson (RBF)
& 0.00795 $\pm$ 0.00464 (Periodic, P, $n=300$)
& 0.00157 $\pm$ 0.00061
& 2.39 \\

Cure AUC
& GG-KM Negative binomial (RBF)
& 0.00687 $\pm$ 0.00464 (Periodic, P, $n=300$)
& 0.00144 $\pm$ 0.00060
& 2.06 \\

\hline
\end{tabular}%
}
\end{table*}
%
% Generated from highdim_best_poisson_nb_parametric_updated.csv
\begin{table*}[tbp]
\caption{
Best observed mean predictive performance in each high-dimensional simulation
scenario. For IPCW IBS and cure RMSE, smaller values are better, whereas for
cure AUC, larger values are better. Parenthesized values are Monte Carlo
standard errors. The comparison includes Poisson and negative-binomial GG-KM
models, Poisson and negative-binomial spline models, the mixture-cure model,
and the promotion-time model. Int and Rad denote Interaction and Radial
surfaces, respectively. GG, Spl, P, and NB denote generalized-gamma kernel
machine, spline, Poisson, and negative binomial, respectively. For GG-KM fits,
the selected kernel is shown in parentheses.
}
\label{tab:highdim-best}
\centering
\resizebox{\textwidth}{!}{%
\begin{tabular}{lllccc}
\hline
True surface
& Generating law
& $n$
& Lowest mean IPCW IBS
& Lowest mean cure RMSE
& Highest mean cure AUC \\
\hline

Interaction
& P
& 300
& GG-NB (polynomial): 0.2045 (0.0017)
& GG-NB (polynomial): 0.1250 (0.0017)
& GG-NB (polynomial): 0.7052 (0.0024) \\

Interaction
& NB(0.25)
& 300
& GG-NB (polynomial): 0.2075 (0.0015)
& GG-NB (polynomial): 0.1254 (0.0017)
& GG-NB (polynomial): 0.7179 (0.0024) \\

Interaction
& NB(1.00)
& 300
& GG-NB (polynomial): 0.2093 (0.0017)
& GG-NB (polynomial): 0.1179 (0.0018)
& GG-NB (polynomial): 0.7289 (0.0023) \\

Interaction
& P
& 600
& GG-NB (polynomial): 0.2018 (0.0011)
& GG-NB (polynomial): 0.0961 (0.0014)
& GG-NB (polynomial): 0.7043 (0.0017) \\

Interaction
& NB(0.25)
& 600
& GG-NB (polynomial): 0.2027 (0.0011)
& GG-NB (polynomial): 0.0912 (0.0013)
& GG-NB (polynomial): 0.7091 (0.0017) \\

Interaction
& NB(1.00)
& 600
& GG-NB (polynomial): 0.2027 (0.0011)
& GG-NB (polynomial): 0.0869 (0.0012)
& GG-P (polynomial): 0.7250 (0.0017) \\

\hline

Radial
& P
& 300
& Spl-P: 0.2138 (0.0014)
& Spl-P: 0.1168 (0.0017)
& GG-NB (polynomial): 0.6521 (0.0023) \\

Radial
& NB(0.25)
& 300
& Spl-P: 0.2167 (0.0014)
& Spl-P: 0.1150 (0.0016)
& GG-NB (polynomial): 0.6602 (0.0023) \\

Radial
& NB(1.00)
& 300
& Spl-NB: 0.2226 (0.0015)
& Spl-NB: 0.1117 (0.0015)
& GG-NB (polynomial): 0.6633 (0.0024) \\

Radial
& P
& 600
& Spl-P: 0.2138 (0.0009)
& Spl-NB: 0.0926 (0.0011)
& GG-NB (polynomial): 0.6426 (0.0019) \\

Radial
& NB(0.25)
& 600
& Spl-P: 0.2136 (0.0009)
& Spl-NB: 0.0897 (0.0010)
& GG-NB (polynomial): 0.6480 (0.0018) \\

Radial
& NB(1.00)
& 600
& Spl-P: 0.2174 (0.0010)
& Spl-P: 0.0881 (0.0010)
& GG-NB (polynomial): 0.6492 (0.0018) \\

\hline
\end{tabular}%
}
\end{table*}
%
\clearpage
% USG records the post-\clearpage counter as the last page; restore the last physical page.
\addtocounter{page}{-1}
\fi

\end{document}

% \documentclass[AMA,twocolumn]{USG}

% \usepackage{float}
% \usepackage{bm}
% \usepackage{booktabs}
% \usepackage{algorithm}
% \usepackage{algpseudocode}
% \usepackage{rotating}
% \usepackage{url}
% \def\UrlBreaks{\do\/\do-\do.\do:\do_\do\~}

% \graphicspath{{./images/}{./figures/review/}}
% \newif\ifggkmsubmission
% \ifdefined\GGKMSubmissionLayout
%   \ggkmsubmissiontrue
% \else
%   \ggkmsubmissionfalse
% \fi
% \ifggkmsubmission
%   \newcounter{submissionfigure}
% \fi

% \newcommand{\GGKMTable}[1]{%
%   \ifggkmsubmission\else\input{#1}\fi}
% \newcommand{\GGKMWideFigure}[4]{%
%   \ifggkmsubmission
%     \begin{figure*}[tbp]
%       \noindent\textit{[Figure~\ref{#3} near here]}\par
%       \caption{#2}\label{#3}
%     \end{figure*}%
%   \else
%     \begin{figure*}[tbp]
%       \centerline{\includegraphics[width=\linewidth]{#1}}%
%       \caption{#2}\label{#3}
%     \end{figure*}%
%   \fi}

% \articletype{RESEARCH ARTICLE}
% \journal{Statistics in Medicine}
% \volume{}
% \startpage{1}
% \articledoi{}

% \begin{document}
% \title{A generalized gamma kernel-machine promotion-time cure model}
% \titlemark{Kernel-machine promotion-time cure model}
% \shrttitle{Kernel-machine promotion-time cure model}
% \author[1]{Gabriel de Jesus Pereira}[https://orcid.org/0009-0003-3254-4854]
% \author[1]{Vera L. D. Tomazella}[https://orcid.org/0000-0002-6780-2089]
% \author[2]{Pedro R. D. Marinho}[https://orcid.org/0000-0003-1591-8300]
% \author[2]{Marcelo Rodrigo Portela Ferreira}[https://orcid.org/0000-0003-0242-0255]
% \author[2]{Agatha S. Rodrigues}[https://orcid.org/0000-0003-1112-8264]
% \authormark{PEREIRA \textsc{et al.}}
% \address[1]{\orgdiv{Department of Statistics, }\orgname{Universidade Federal de S\~ao Carlos, }\orgaddress{\state{S\~ao Carlos, }\country{Brazil}}}
% \address[2]{\orgdiv{Department of Statistics, }\orgname{Universidade Federal da Para\'iba, }\orgaddress{\state{Jo\~ao Pessoa, }\country{Brazil}}}
% \corres{Vera L. D. Tomazella, Department of Statistics, Universidade Federal de S\~ao Carlos, S\~ao Carlos, Brazil. \textbf{AUTHOR CONFIRMATION REQUIRED:} corresponding-author email.}
% \keywords{cure model | generalized gamma | kernel machine | promotion time | survival analysis}
% \abstract[ABSTRACT]{{\color{red}Cure models allow the long-term event-free fraction to be modelled explicitly,
% but the commonly used log-linear predictor in promotion-time cure models may
% be inadequate when covariate effects on the latent-cause mechanism are nonlinear
% or involve interactions. We propose a generalized gamma kernel-machine
% promotion-time cure model (GG-KM) that combines a flexible nonlinear predictor
% for the latent-cause intensity with a parametric generalized gamma latency
% distribution. The latent-cause intensity is represented in a reproducing-kernel
% Hilbert space, while Poisson and Negative-Binomial specifications allow the
% latent number of causes to accommodate both equidispersion and overdispersion.
% Estimation is performed using penalized expectation-maximization or
% expectation/conditional maximization either, followed by bounded
% observed-likelihood refinement and explicit convergence diagnostics. We evaluate the proposed framework through simulation experiments involving
% linear, quadratic, and periodic covariate effects, as well as a four-covariate
% mechanism design, and through an application to a melanoma cohort. Predictive
% performance is assessed using event-stratified training/test splits,
% training-only tuning, and inverse-probability-of-censoring weighted integrated
% Brier scores, together with calibration and cure-fraction accuracy. The results
% show that no single predictor is uniformly superior: simpler parametric
% predictors perform best when the underlying covariate structure is correctly
% specified, whereas the GG-KM provides competitive performance across different
% nonlinear mechanisms and performs particularly well for interaction and
% radial structures. These findings support the GG-KM as a flexible and robust
% alternative when the functional form of the latent-cause mechanism is unknown.
% The melanoma application illustrates the practical use of the proposed framework
% for out-of-sample prediction.}}

% \maketitle

% \section{Introduction}\label{sec:introduction}

Standard survival models assign vanishing probability to remaining event
free as follow-up grows arbitrarily long. That limit is inappropriate
when a clinically meaningful subgroup will never experience the event of
interest, a situation indicated by a sustained plateau in the survival
curve under mature follow-up. Cure models describe this setting by
separating the long-term event-free probability from the timing of
events among individuals who remain susceptible, which makes the cured
proportion an interpretable quantity in oncology and other long-term-survivor
studies \cite{berkson1952,tsodikov2003}.

{
Accommodating this plateau requires a model that allows a positive
limiting survival probability, and three main constructions are available
in the literature for this purpose, differing in the mechanism by which
they generate the cure mass. The first, the mixture cure model,
introduces a latent binary indicator that separates long-term event-free
individuals from those who remain susceptible, so that the cured
proportion and the survival of the susceptible subpopulation are
specified as separate components, typically a regression model for
incidence together with a proper distribution for the latency
\cite{boag1949,berkson1952,farewell1982,sy2000}. The second, the
defective-distribution model, does not represent the cured fraction
through a separate component; instead, the cured fraction arises
intrinsically from the parameter space of a single parametric survival
distribution. For most parameter values the distribution is proper and
the survival function decays to zero, as in conventional survival models;
over a distinct region of the same parameter space the distribution
becomes improper and the survival function instead converges to a
strictly positive limit, which is the cured fraction. A single family
then governs both the failure times of the susceptible individuals and
the size of the cured group, with no separate incidence model and with
covariates entering through its parameters
\cite{balka2009,rocha2017,scudilio2019}. The third, based on latent
competing risks, generates the cure mass from an unobserved count of
competing causes, so that the cured fraction is The probability that this
count is zero \cite{yakovlev1996,chen1999,tsodikov2003,rodrigues2009unif}.
} %This construction can be re-expressed as a mixture whose cured proportion equals that zero-count probability and whose susceptible latency is induced by the count law, so it is a structured special case of the mixture model in which the incidence and latency components are linked through shared parameters rather than specified separately \cite{rodrigues2009}.
{
This mechanism admits a two-stage description. An initiation stage sets
the unobserved number $N$ of latent causes carried by an individual, and
a maturation stage assigns each cause an activation time; the earliest
activation produces the event, so $N=0$ describes an individual who never
fails and identifies the cure mass \cite{rodrigues2025}. Encoding the
count law through its probability-generating function links the latency
distribution to both the population survival and the cure fraction, and
it is this representation on which the model developed here is built.
}

{
In the promotion-time construction the mean number of latent causes is
almost always taken to be log-linear in the covariates
\cite{chen1999,rodrigues2009}. This specification is convenient and
keeps the cure fraction in closed form, but it forces every covariate
to act monotonically on the latent-cause intensity, so threshold effects,
interactions, and periodic or otherwise nonlinear associations cannot be
represented. Because the intensity governs the cure fraction, a
misspecified incidence model biases both the estimated covariate effects
and the cured proportion \cite{amico2018}. The activation-time component
carries another issue: fixing the latency to a restrictive family can be
equally misspecified, and in a cure model this error propagates to the
covariate effects and to the cure fraction. A flexible parametric family
mitigates this while preserving interpretability and a tractable
likelihood, and the generalized gamma is a natural choice, since it
nests the exponential, Weibull, gamma, and log-normal distributions and
spans a wide range of hazard shapes
\cite{stacy1962,prentice1974,cox2007}. Together, these concerns motivate
treating both components flexibly, through a flexible latent-intensity
predictor for the incidence and a sufficiently general parametric latency
for the activation times.
}

{
Flexible predictors have been proposed in cure modelling: smoothers,
support-vector machines, neural networks, and deep architectures have all
been used to relax the usual linear specification. In the promotion-time
setting, Pal and Aselisewine\cite{pal2023} introduced a support-vector
formulation for the incidence component, and Xie and Yu\cite{xie2021}
used a neural network for an unstructured promotion-time component,
while deep promotion-time models and structured nonlinear cure models
further broaden the available high-dimensional architectures\cite{medina2024,feng2026}.
These methods make the covariate effect on the latent-cause intensity
nonlinear, but they differ substantially in how they treat that
intensity, the latency distribution, and the count structure. To our
knowledge, no existing promotion-time approach combines all three within
a single likelihood-based estimation and prediction framework: a
flexible nonlinear predictor for the latent-cause intensity, which shapes
both the cured fraction and the timing of events; a transparent and
sufficiently general parametric latency; and an explicit count-family
treatment of overdispersion.
}

To address this gap, we develop a generalized gamma kernel-machine
promotion-time cure model (GG-KM) that places the log-intensity of
the latent-cause process in a reproducing-kernel Hilbert space (RKHS)
\cite{kimeldorf1971}, while retaining a flexible parametric
generalized gamma latency distribution \cite{stacy1962} and allowing
either Poisson or negative-binomial latent counts. This formulation
separates two sources of modelling flexibility: the kernel predictor
captures potentially complex nonlinear and interaction effects of
covariates on the latent-cause mechanism, whereas the generalized
gamma distribution provides a flexible parametric representation of
the activation-time process.


The remainder of this article is organized as follows. Section~\ref{sec:methods}
presents the proposed generalized gamma kernel-machine promotion-time cure
model, covering the general promotion-time construction, the generalized gamma
latency distribution, the kernel-based latent-cause predictor, and the penalized
EM and ECME estimation procedures with their observed-likelihood refinement and
convergence assessment. Section~\ref{sec:simulation} describes the simulation
study and evaluates the predictive performance of the proposed model against
log-linear, additive-spline, and quadratic alternatives under different
covariate structures and latent-count distributions. Section~\ref{sec:melanoma}
applies the model to the melanoma data and reports its predictive behaviour.
Section~\ref{sec:discussion} discusses the main findings, limitations, and
directions for future research.

% \section{Methods}\label{sec:methods}

This section first recalls, as background, the general promotion-time cure construction that the model builds on, and then develops the proposed GG-KM model: a generalized gamma latency, a regularized reproducing-kernel Hilbert space predictor for the latent-cause intensity, and Poisson or negative-binomial latent counts. The remaining subsections give the penalized estimation algorithms and the prediction and evaluation protocol. Observed likelihoods and conditional objective functions are stated in the main text, while lengthy derivations and analytical gradients are deferred to the Supporting Information.

%{ Figure~\ref{fig:ggkm-workflow} provides an overview of the GG-KM model and its estimation and evaluation framework. Panel A illustrates the kernel-based latent-cause intensity $\theta(\bm{x})=\exp\{f(\bm{x})\}$ and its link to the Poisson and negative-binomial latent-count distributions. Panel B depicts the promotion-time mechanism, in which individuals with $N_i=0$ are cured, whereas for $N_i\geq1$ the event time is the minimum of the latent activation times. Panel C shows the resulting population survival curves and the cure plateau $p_c(\bm{x})$. Panel D summarizes the estimation and prediction pipeline, including the event-stratified data split, training-only preprocessing, penalized EM or ECME estimation, out-of-sample prediction, and IPCW integrated Brier evaluation.}

 %Figure~\ref{fig:ggkm-workflow} lays out the GG-KM model and its estimation and evaluation path. Panel A shows the latent-cause intensity $\theta(\bm{x})=\exp\{f(\bm{x})\}$ built by a radial-basis kernel expansion over $(x_1,x_2)\in[-1,1]^2$, with the six kernel centres $\bm{x}_j$ marked and the count law named beneath the surface. Panel B plots four parallel time axes for two cured individuals ($N_i=0$) and two susceptible ones ($N_i\ge1$). Its key maps every glyph to a quantity: an open circle is an activation time $Z_{ij}$ from the generalized gamma latency, a filled circle is the event time $T_i=\min_jZ_{ij}$, and a vertical bar is the censoring time $C_i$; the resulting pair $(Y_i,\delta_i)$ is printed at the right of each axis. Panel C draws the matched Poisson and negative-binomial population survival curves, the cure plateau $p_c(\bm{x})$, and the shaded interval at or below the scoring horizon $\tau$. Panel D traces the leakage-controlled pipeline: event-stratified 70/30 split, Gram matrix construction, penalized EM or ECME with bounded refinement, kernel-coefficient prediction, and IPCW Brier evaluation. The dashed arc in Panel D is the only feedback path and marks the single iterative update that revises the kernel coefficients, the latency parameters, and the negative-binomial dispersion when applicable.

%\GGKMWideFigure{ggkm-method-diagram.pdf}{Structure of the GG-KM model and of its estimation and evaluation. Panel A shows the generative mechanism: a kernel predictor $f(\bm{x})=\beta_0+\sum_j\alpha_jK(\bm{x},\bm{x}_j)$ over the marked kernel centres $\bm{x}_j$ determines the positive latent-cause intensity $\theta(\bm{x})=\exp\{f(\bm{x})\}$, which indexes the Poisson or negative-binomial count $N_i$ named beneath the surface. Panel B shows the observation mechanism. Given $N_i\ge1$ the promotion times $Z_{ij}$ are independent generalized gamma variables and $T_i=\min_jZ_{ij}$, whereas $N_i=0$ gives $T_i=\infty$, so only $Y_i=\min(T_i,C_i)$, $\delta_i=I(T_i\leq C_i)$ and $\bm{x}_i$ are observed; the key maps the open circle, the filled circle and the vertical bar to $Z_{ij}$, $T_i$ and $C_i$, and the realized pair $(Y_i,\delta_i)$ is printed at the right of each time axis. Panel C shows the population survival curves, where P denotes Poisson and NB negative-binomial counts. They use the simulation latency $(a,d,p)=(1.6,1.3,1.1)$, a matched cure fraction of $0.45$ and dispersion $\phi=1$; they are illustrative rather than fitted, and line style with markers distinguishes them in addition to colour. Shading marks times at or below the frozen horizon $\tau$ within which predictions are scored, so the approach of $S_{\mathrm{pop}}(t\mid\bm{x})$ to $p_c(\bm{x})$ is extrapolated rather than evaluated. Panel D shows the estimation and evaluation path: one event-stratified 70/30 split with training-only preprocessing, the validated Gram matrix, penalized EM or ECME with an ascent safeguard, out-of-sample prediction, and IPCW integrated Brier scoring with paired differences. The dashed arc in Panel D is the only feedback path and marks the iterative update of $(\beta_0,\bm{\alpha})$, $(a,d,p)$ and $\phi$.}{fig:ggkm-workflow}{Four-panel schematic. Panel A: a kernel predictor with marked centres gives the latent-cause intensity that indexes a Poisson or negative-binomial latent count. Panel B: four time axes for two susceptible and two cured individuals, with a key mapping open circles to activation times, filled circles to event times and vertical bars to censoring times, and the observed time and event indicator printed at the right of each axis. Panel C: population survival curves for both count laws approach a common cure plateau beyond the shaded scored horizon. Panel D: event-stratified split with training-only preprocessing, validated Gram matrix, penalized EM or ECME with ascent safeguard, out-of-sample prediction, and IPCW integrated Brier evaluation, with one dashed feedback path.}

\subsection{General promotion-time cure construction}\label{sec:promotion-time}

{
	%Let $\mathcal{D}_n=\{(Y_i,\delta_i,\bm{x}_i):i=1,\dots,n\}$ denote an independent sample of size $n$, where $T_i$ is the event time, $C_i$ the right-censoring time, $Y_i=\min(T_i,C_i)$, $\delta_i=I(T_i\leq C_i)$, and $\bm{x}_i$ a covariate vector. Censoring is assumed conditionally independent, $T_i\perp C_i\mid\bm{x}_i$. The latent variable $N_i$ counts the potential causes of the event, and $N_i=0$ identifies a cured individual.

Let $T_i$ denote the event time and $C_i$ the right-censoring time. The observed time is the random variable $Y_i=\min(T_i,C_i)$, whose realized value is written $t_i$, and $\delta_i=I(T_i\leq C_i)$ is the event indicator, with $\bm{x}_i$ a covariate vector, for $i=1,\ldots,n$. The observed data are therefore $\mathcal{D}_{\text{obs}}=\{(t_i,\delta_i,\bm{x}_i)\}^{n}_{i=1}$. Censoring is assumed conditionally independent, $T_i\perp C_i\mid\bm{x}_i$. The latent variable $N_i$ counts the potential causes of the event, and $N_i=0$ identifies a cured individual.

Conditional on the covariates, $N_i\mid\bm{x}_i$ has probability mass function $\pi_k(\bm{x}_i)=P(N_i=k\mid\bm{x}_i)$, $k=0,1,2,\dots$, and probability-generating function $A_i(u)=\mathbb{E}\left(u^{N_i}\mid\bm{x}_i\right)=\sum_{k\ge0}\pi_k(\bm{x}_i)u^k$. Given $N_i=k\ge1$, the activation times $Z_{i1},\dots,Z_{ik}$ are independent and identically distributed with a proper distribution function $F_0$, survival function $S_0=1-F_0$, and density $f_0$; conditional on $\bm{x}_i$ they are independent of $N_i$, and their distribution does not depend on $N_i$. In the formulation considered here the covariates act on the latent-count distribution but not on the latency $F_0$. The event time is the earliest activation, $T_i=\min_{1\le j\le N_i}Z_{ij}$ when $N_i\ge1$, whereas $T_i=\infty$ when $N_i=0$.

Conditional on $N_i=k$, the event exceeds $t$ only if all $k$ activation times do, so $P(T_i>t\mid N_i=k,\bm{x}_i)=S_0(t)^k$. Averaging over the latent-count distribution gives the population survival function and, as $t\to\infty$ with $S_0(t)\to0$, the cure fraction,
\begin{equation}
\begin{aligned}
S_{\mathrm{pop}}(t\mid\bm{x}_i)
&=P(T_i>t\mid\bm{x}_i)
=\sum_{k\ge0}\pi_k(\bm{x}_i)S_0(t)^k
=A_i\{S_0(t)\},\\
p_c(\bm{x}_i)
&=\lim_{t\to\infty}S_{\mathrm{pop}}(t\mid\bm{x}_i)
=A_i(0)=P(N_i=0\mid\bm{x}_i).
\end{aligned}
\label{eq:promotion-population}
\end{equation}
The population survival is the generating function evaluated at the latency survival, and the cure fraction is its value at zero, equal to the probability of no latent cause. Differentiating with respect to time, and using $S_0'=-f_0$, gives the population density and hazard,
\begin{equation}
\begin{aligned}
f_{\mathrm{pop}}(t\mid\bm{x}_i)
&=f_0(t)\,A_i'\{S_0(t)\},\\
h_{\mathrm{pop}}(t\mid\bm{x}_i)
&=\frac{f_{\mathrm{pop}}(t\mid\bm{x}_i)}
        {S_{\mathrm{pop}}(t\mid\bm{x}_i)},
\end{aligned}
\label{eq:promotion-density-hazard}
\end{equation}
where $A_i'$ denotes differentiation with respect to its argument. When $p_c(\bm{x}_i)>0$, $f_{\mathrm{pop}}$ is an improper subdensity that integrates to $1-p_c(\bm{x}_i)$ rather than to one. The construction satisfies the required boundary behaviour: $S_{\mathrm{pop}}(0\mid\bm{x})=1$, $S_{\mathrm{pop}}(t\mid\bm{x})$ is non-increasing in $t$, and $\lim_{t\to\infty}S_{\mathrm{pop}}(t\mid\bm{x})=p_c(\bm{x})$. % A degenerate cure fraction $p_c(\bm{x})=1$ corresponds to the latent-count mass concentrating entirely on $N_i=0$, that is $\pi_0(\bm{x})=1$, which is a limiting boundary case rather than a value attained at interior count parameters.

Equations~\eqref{eq:promotion-population} and \eqref{eq:promotion-density-hazard} give the standard promotion-time construction for a generic activation-time distribution. The proposed GG-KM model, developed in the following subsection, specializes the latency $F_0$ to the Stacy generalized gamma distribution and models the latent-count mean $\theta_i=\theta(\bm{x}_i)$ through a reproducing-kernel Hilbert space predictor, with the count taken to be Poisson or negative binomial.}

\subsection{The proposed GG-KM model}\label{sec:model-specification}

%{The proposed GG-KM model specializes the general construction of Section~\ref{sec:promotion-time} in three linked components: the activation-time distribution, the latent-cause intensity, and the latent-count law. The activation times follow a Stacy generalized gamma latency with parameters $(a,d,p)$ shared across individuals; the latent-cause intensity is $\theta(\bm{x})=\exp\{f(\bm{x})\}$, where $f$ is a regularized reproducing-kernel Hilbert space predictor; and the latent count $N_i\mid\bm{x}_i$ is Poisson or negative binomial with mean $\theta_i=\theta(\bm{x}_i)$. Covariates act on the model only through the intensity $\theta(\bm{x})$, and hence only through $f$. The parameters satisfy $a,d,p>0$ and $\theta_i>0$, and the negative-binomial count adds a dispersion parameter $\phi>0$. The three components are specified in turn below.}

The activation time of each latent cause follows the Stacy generalized gamma distribution with scale $a>0$ and shape parameters $d>0$ and $p>0$.\cite{stacy1962} Its density, distribution function, and survival function are
\begin{equation}
	\begin{aligned}
		f_{GG}\left(t;a,d,p\right)
		&=\frac{p t^{d-1}\exp\{-(t/a)^p\}}
		{a^d\Gamma(d/p)},\\
		F_{GG}\left(t;a,d,p\right)
		&=P\{d/p,(t/a)^p\},\\
		S_{GG}\left(t;a,d,p\right)
		&=Q\{d/p,(t/a)^p\},
	\end{aligned}
	\label{eq:gg-distribution}
\end{equation}
where $P$ and $Q$ are the regularized lower and upper incomplete gamma
functions. This parameterization contains Weibull models when $d=p$,
gamma models when $p=1$, and the exponential model when $d=p=1$. The
lognormal distribution is not a special case or limit of the Stacy
parameterization at fixed $(a,d,p)$; obtaining it requires a separate
reparameterized limit and is therefore not claimed here. The
implementation verifies numerically that the density integrates to one
and that the derivative of the distribution function equals the density,
so $F_{GG}$, $f_{GG}$, and $S_{GG}$ remain mutually consistent. Figure~\ref{fig:gg-latency-flexibility} illustrates how the verified special cases and a general parameter setting produce distinct density and hazard shapes.

\GGKMWideFigure{gg-latency-flexibility.pdf}{Generalized gamma latency flexibility under the Stacy parameterization. Panel A shows the density $f_{GG}\left(t; a,d,p\right)$ and Panel B the corresponding hazard $h_{GG}\left(t;a,d,p\right)=f_{GG}\left(t;a,d,p\right)/S_{GG}\left(t;a,d,p\right)$ for the exponential case $(a,d,p)=(1.6,1,1)$, Weibull case $(1.6,0.7,0.7)$, gamma case $(1.6,2,1)$, and the general generalized gamma setting $(1.6,1.3,1.1)$ used in the simulation study. Line type and markers distinguish the four cases in addition to colour.}{fig:gg-latency-flexibility}{Two panels compare four generalized gamma latency distributions. The left panel shows their densities over time and the right panel shows their hazards. Exponential, Weibull, gamma, and general generalized gamma curves use distinct line types and markers.}


{

	The GG-KM model relaxes the usual log-linear specification of the latent-cause
	intensity by representing
	$
	    \eta(\bm{x})
	    =
	    \log\{\theta(\bm{x})\}
	$
	through a reproducing kernel Hilbert space (RKHS). Here,
	$\theta(\bm{x})>0$ denotes the conditional mean of the latent number of causes,
	\begin{equation}
	    \theta(\bm{x})
	    =
	    \mathbb{E}(N\mid\bm{x}),
	\end{equation}
	under both the Poisson and negative-binomial parameterizations considered in
	this work. Modelling $\eta(\bm{x})$ in an RKHS allows the association between
	the covariates and the latent-cause mechanism to be nonlinear without requiring
	a prespecified finite-dimensional functional form.

	% The covariate effects enter the model through the mean parameter of the
	% latent number of competing causes. Let
	% $\mathcal{X}\subseteq\mathbb{R}^{r}$ denote the covariate space and $\bm{x},\bm{x}'\in\mathcal{X}$ two covariate vectors, and let $K:\mathcal{X}\times\mathcal{X}\to\mathbb{R}$ be a kernel measuring their similarity. The
	% nonlinear component of the predictor is modelled through a reproducing
	% kernel selected from the following family:
	% \begin{equation}
	% 	\begin{aligned}
	% 		K_{\mathrm{lin}}(\bm{x},\bm{x}')
	% 		&=\bm{x}^{\mathsf T}\bm{x}',\\
	% 		K_{\mathrm{poly}}(\bm{x},\bm{x}')
	% 		&=(\gamma\,\bm{x}^{\mathsf T}\bm{x}'+c_0)^{q},\\
	% 		K_{\mathrm{rbf}}(\bm{x},\bm{x}')
	% 		&=\exp\{-\gamma\|\bm{x}-\bm{x}'\|_2^2\},\\
	% 		K_{\mathrm{exp}}(\bm{x},\bm{x}')
	% 		&=\exp\{-\gamma\|\bm{x}-\bm{x}'\|_2\},\\
	% 		K_{\mathrm{lap}}(\bm{x},\bm{x}')
	% 		&=\exp\{-\gamma\|\bm{x}-\bm{x}'\|_1\},\\
	% 		K_{\mathrm{cau}}(\bm{x},\bm{x}')
	% 		&=\{1+\gamma\|\bm{x}-\bm{x}'\|_2^2\}^{-1},\\
	% 		K_{\mathrm{wav}}(\bm{x},\bm{x}')
	% 		&=
	% 		\prod_{j=1}^{r}
	% 		\left[
	% 		\cos\!\left(
	% 		1.75\,\frac{x_j-x'_j}{A}
	% 		\right)
	% 		\exp\!\left\{
	% 		-\frac{1}{2}
	% 		\left(
	% 		\frac{x_j-x'_j}{A}
	% 		\right)^2
	% 		\right\}
	% 		\right].
	% 	\end{aligned}
	% 	\label{eq:kernels}
	% \end{equation}
	% Here $\gamma>0$ is a scale parameter, $q\geq1$ is an integer polynomial
	% degree, $c_0\geq0$ is an offset, and $A>0$ is the scale parameter of the
	% wavelet kernel. The RBF kernel corresponds to the Gaussian kernel. The
	% exponential and Laplacian kernels differ in the distance used in their
	% construction: the former uses the Euclidean ($L_2$) norm, whereas the
	% latter uses the $L_1$ norm.

	% Each kernel in Equation~\eqref{eq:kernels} is positive semidefinite and
	% therefore defines an associated reproducing-kernel Hilbert space
	% $\mathcal H_K$. In particular, the wavelet kernel is constructed as a
	% product across coordinates of cosine-modulated Gaussian kernels;
	% positive semidefiniteness is preserved under products of
	% positive-semidefinite kernels.

	% The semiparametric version of the generalized representer theorem
	% \cite{scholkopf2001generalized} allows the RKHS component to be
	% augmented by a finite-dimensional, unpenalized parametric component.
	% Let
	% $
	% \{\psi_\ell\}_{\ell=1}^{M}
	% $
	% be a collection of real-valued functions on the covariate space
	% $\mathcal X$, and suppose that the $n\times M$ matrix
	% $
	% \bm\Psi
	% =
	% \bigl(\psi_\ell(\bm{x}_i)\bigr)_
	% {\substack{i\in\left[n\right]\\\ell\in\left[M\right]}}
	% $
	% has rank $M$. Consider a predictor of the form
	% \begin{equation*}
	% 	\widetilde f=f+h,
	% 	\qquad
	% 	f\in\mathcal H_K,
	% 	\qquad
	% 	h\in\operatorname{span}\{\psi_1,\ldots,\psi_M\},
	% \end{equation*}
	% obtained by minimizing a regularized criterion of the form
	% \begin{equation*}
	% 	\mathcal C
	% 	\left(
	% 	(\bm{x}_1,y_1,\widetilde f(\bm{x}_1)),
	% 	\ldots,
	% 	(\bm{x}_n,y_n,\widetilde f(\bm{x}_n))
	% 	\right)
	% 	+
	% 	G\!\left(\|f\|_{\mathcal H_K}\right),
	% \end{equation*}
	% where $G$ is strictly increasing. Under these conditions, the
	% semiparametric representer theorem implies that any minimizer admits
	% the representation
	% \begin{equation}
	% 	\widetilde f^{\,\star}(\bm{x})
	% 	=
	% 	\sum_{j=1}^{n}
	% 	\alpha_j K(\bm{x},\bm{x}_j)
	% 	+
	% 	\sum_{\ell=1}^{M}
	% 	\beta_\ell\psi_\ell(\bm{x}).
	% 	\label{eq:semiparametric-representer}
	% \end{equation}
	% Thus, only the component belonging to $\mathcal H_K$ is regularized,
	% whereas the finite-dimensional component spanned by the functions
	% $\psi_\ell$ remains unpenalized.

	A semiparametric extension of the generalized representer theorem
	\cite{scholkopf2001generalized} allows an RKHS component to be combined with
	an unpenalized finite-dimensional parametric component. Let
	$
	    \mathcal H_0
	    =
	    \operatorname{span}\{\psi_1,\ldots,\psi_M\},
	$
	where $\psi_1,\ldots,\psi_M:\mathcal X\to\mathbb R$, and define
	\begin{equation*}
	    \Psi_{i\ell}
	    =
	    \psi_\ell(\bm{x}_i),
	    \qquad
	    i\in \left[n\right],
	    \quad
	    \ell\in\left[M\right].
	\end{equation*}
	Assume that the resulting $n\times M$ matrix $\bm\Psi$ has full column rank
	$M$, so that the coefficients of the finite-dimensional component are
	identifiable on the observed sample.

	Consider predictors of the form
	\begin{equation*}
	    \widetilde f
	    =
	    f+h,
	    \qquad
	    f\in\mathcal H_K,
	    \qquad
	    h\in\mathcal H_0,
	\end{equation*}
	and suppose that estimation is based on minimizing a criterion of the form
	\begin{equation*}
	    \mathcal C
	    \left(
	        \widetilde f(\bm{x}_1),
	        \ldots,
	        \widetilde f(\bm{x}_n)
	    \right)
	    +
	    G\!\left(\|f\|_{\mathcal H_K}\right),
	\end{equation*}
	where the data are regarded as fixed arguments of $\mathcal C$ and
	$G:[0,\infty)\to\mathbb R$ is strictly increasing. Under these conditions,
	the generalized representer theorem implies that the RKHS component of any
	minimizer can be represented using only kernel sections evaluated at the
	training observations. Consequently, a minimizer admits the representation
	\begin{equation}
	    \widetilde f^{\,\star}(\bm{x})
	    =
	    \sum_{j=1}^{n}
	        \alpha_j K(\bm{x},\bm{x}_j)
	    +
	    \sum_{\ell=1}^{M}
	        \beta_\ell\psi_\ell(\bm{x}).
	    \label{eq:semiparametric-representer}
	\end{equation}
	Thus, regularization is applied only to the RKHS component, whereas the
	finite-dimensional component remains unpenalized.

	For the GG-KM model, estimation is formulated as the minimization of a
	penalized negative log-likelihood over the latent-cause predictor
	$\eta$. Let
	$\ell(\eta;\mathcal D_{\text{obs}})$ denote the corresponding observed-data
	log-likelihood. The corresponding penalized estimation problem is
	\begin{equation}
	    \min_{\eta}
	    \left\{
	        -\ell(\eta;\mathcal D_{\text{obs}})
	        +
	        \frac{\lambda}{2}
	        \|f\|_{\mathcal H_K}^{2}
	    \right\},
	    \qquad
	    \lambda>0,
	    \label{eq:ggkm-functional-objective}
	\end{equation}
	where the penalization is applied only to the RKHS component of
	$\eta$.

	To place this optimization problem within the semiparametric
	representer-theorem framework, the unpenalized finite-dimensional
	space is chosen to contain only an intercept. Specifically, setting
	$
	    M=1,
		\text{ }
	    \psi_1(\bm{x})=1,
	$
	so that
	$
	    \mathcal H_0
	    =
	    \operatorname{span}\{1\}.
	$
	Consequently, the class of predictors over which
	Equation~\eqref{eq:ggkm-functional-objective} is minimized is
	\begin{equation}
	    \eta(\bm{x})
	    =
	    \beta_0+f(\bm{x}),
	    \qquad
	    \beta_0\in\mathbb R,
	    \qquad
	    f\in\mathcal H_K.
	    \label{eq:ggkm-semiparametric-predictor}
	\end{equation}
	Thus, the intercept constitutes the unpenalized finite-dimensional
	component, whereas $f$ contains the possibly nonlinear covariate
	effect represented in the RKHS. In this case, the matrix
	$\bm\Psi$ reduces to the $n$-vector of ones and has rank one for
	every nonempty sample.

	Substituting
	Equation~\eqref{eq:ggkm-semiparametric-predictor} into the likelihood,
	the estimation problem can be written explicitly as
	\begin{equation}
	    (\widehat{\beta}_0,\widehat f)
	    =
	    \arg\min_{
	        \beta_0\in\mathbb R,\,
	        f\in\mathcal H_K
	    }
	    \left\{
	        -\ell(\beta_0,f;\mathcal D_{\text{obs}})
	        +
	        \frac{\lambda}{2}
	        \|f\|_{\mathcal H_K}^{2}
	    \right\}.
	    \label{eq:ggkm-penalized-objective}
	\end{equation}
	Here, $\ell(\beta_0,f;\mathcal D_{\text{obs}})$ denotes the same observed-data
	log-likelihood evaluated through the predictor
	$\eta(\bm{x})=\beta_0+f(\bm{x})$.

	The regularization term in
	Equation~\eqref{eq:ggkm-penalized-objective} is of the form required by
	the representer theorem. Indeed, for
	\begin{equation*}
	    G(\|f\|_{\mathcal H_K})
	    =
	    \frac{\lambda}{2}\|f\|_{\mathcal H_K}^{2},
	    \qquad
	    \|f\|_{\mathcal H_K}\geq 0,
	\end{equation*}
	$G$ is strictly increasing when $\lambda>0$. Therefore, the
	semiparametric representer theorem applies to the minimization problem
	in Equation~\eqref{eq:ggkm-penalized-objective}. It follows that the
	RKHS component of any minimizer can be expressed as
	$
	    \widehat f(\bm{x})
	    =
	    \sum_{j=1}^{n}
	        \widehat{\alpha}_j
	        K(\bm{x},\bm{x}_j).
	$
	Hence, although
	Equation~\eqref{eq:ggkm-penalized-objective} is initially posed as an
	optimization problem over the infinite-dimensional space
	$\mathcal H_K$, its solution can be sought within the finite-dimensional
	family
	\begin{equation}
	    \widehat{\eta}(\bm{x})
	    =
	    \widehat{\beta}_0
	    +
	    \sum_{j=1}^{n}
	        \widehat{\alpha}_j
	        K(\bm{x},\bm{x}_j).
	    \label{eq:rkhs-predictor}
	\end{equation}

	Let $\bm K$ denote the $n\times n$ training Gram matrix, with entries
	$
	    K_{ij}
	    =
	    K(\bm{x}_i,\bm{x}_j).
	$
	For
	$
	    f(\cdot)
	    =
	    \sum_{j=1}^{n}
	        \alpha_j K(\cdot,\bm{x}_j),
	$
	the squared RKHS norm satisfies
	$
	    \|f\|_{\mathcal H_K}^{2}
	    =
	    \bm{\alpha}^{\mathsf T}
	    \bm K
	    \bm{\alpha}.
	$
	Accordingly, the penalized estimation problem reduces to the
	finite-dimensional optimization
	\begin{equation}
	    \min_{\beta_0\in\mathbb R,\,
	         \bm{\alpha}\in\mathbb R^n}
	    \left\{
	        -\ell(\beta_0,\bm{\alpha};\mathcal D_{\text{obs}})
	        +
	        \frac{\lambda}{2}
	        \bm{\alpha}^{\mathsf T}
	        \bm K
	        \bm{\alpha}
	    \right\}.
	    \label{eq:ggkm-finite-objective}
	\end{equation}
	Thus, increasing $\lambda$ imposes stronger regularization on the
	nonparametric RKHS component of the predictor, while the intercept
	$\beta_0$ remains unpenalized.

	Finally, the GG-KM predictor is linked to the conditional mean of the
	latent number of causes through
	\begin{equation}
	    \theta(\bm{x})
	    =
	    \exp\{\eta(\bm{x})\}.
	    \label{eq:theta-link}
	\end{equation}
	Therefore, $\eta(\bm{x})$ is unrestricted on $\mathbb R$, whereas the
	exponential link guarantees $\theta(\bm{x})>0$ without requiring an
	explicit positivity constraint.


	% The estimation of $\eta\left(\bm{x}\right) = \log\left(\theta\left(\bm{x}\right)\right)$
	% is formulated through the minimization of a penalized negative log-likelihood function:

	% \begin{equation}
	% 	\min \left\{- \ell\left(\bm{x}; \beta_0, f\right) + \frac{\lambda}{2} \|f\|_{\mathcal H_K}^{2}\right\},
	% \end{equation}

	% For the GG-KM model, the finite-dimensional component consists only of
	% an intercept. Setting
	% $M=1,
	% \text{ }
	% \psi_1(\bm{x})=1,
	% $
	% gives
	% $
	% h(\bm{x})=\beta_0.
	% $
	% The corresponding matrix $\bm\Psi$ is the $n$-vector of ones and has
	% rank one for any nonempty sample, so the rank condition of the
	% semiparametric representer theorem is satisfied.

	% Identifying the complete predictor $\widetilde f$ with the GG-KM
	% predictor $\eta$, therefore
	% \begin{equation*}
	% 	\eta(\bm{x})
	% 	=
	% 	\widetilde f(\bm{x})
	% 	=
	% 	f(\bm{x})+\beta_0,
	% 	\qquad
	% 	f\in\mathcal H_K.
	% \end{equation*}
	% The corresponding penalized objective is
	% \begin{equation*}
	% 	\mathcal J(\beta_0,f)
	% 	=
	% 	\mathcal L\!\left(
	% 	\beta_0+f(\bm{x}_1),\ldots,
	% 	\beta_0+f(\bm{x}_n)
	% 	\right)
	% 	+
	% 	\frac{\lambda}{2}
	% 	\|f\|_{\mathcal H_K}^{2},
	% 	\qquad
	% 	\lambda>0.
	% \end{equation*}
	% Since the penalty depends only on the RKHS component $f$, the intercept
	% $\beta_0$ is not regularized. By
	% Equation~\eqref{eq:semiparametric-representer}, the RKHS component of
	% the minimizer has the finite-dimensional representation
	% \begin{equation*}
	% 	f^\star(\bm{x})
	% 	=
	% 	\sum_{j=1}^{n}
	% 	\alpha_jK(\bm{x},\bm{x}_j),
	% \end{equation*}
	% and consequently the GG-KM predictor takes the form
	% \begin{equation}
	% 	\eta(\bm{x})
	% 	=
	% 	\beta_0
	% 	+
	% 	\sum_{j=1}^{n}
	% 	\alpha_jK(\bm{x},\bm{x}_j).
	% 	\label{eq:rkhs-predictor}
	% \end{equation}

	% If $\bm K$ denotes the training Gram matrix, with
	% $
	% K_{ij}=K(\bm{x}_i,\bm{x}_j),
	% $
	% then
	% \begin{equation*}
	% 	\|f\|_{\mathcal H_K}^{2}
	% 	=
	% 	\bm{\alpha}^{\mathsf T}
	% 	\bm K
	% 	\bm{\alpha},
	% \end{equation*}
	% and the regularization term becomes
	% $
	% \frac{\lambda}{2}
	% \bm{\alpha}^{\mathsf T}
	% \bm K
	% \bm{\alpha}.
	% $
	% Hence, increasing $\lambda$ imposes stronger regularization on the
	% RKHS component of the predictor while leaving the intercept
	% $\beta_0$ unpenalized.

	% For $\lambda>0$, the quadratic penalty is a strictly increasing
	% function of $\|f\|_{\mathcal H_K}$ and therefore satisfies the
	% condition required by the representer theorem.

	% Finally, the mean parameter of the latent number of competing causes
	% is related to the predictor through
	% \begin{equation*}
	% 	\theta(\bm{x})
	% 	=
	% 	\exp\{\eta(\bm{x})\},
	% \end{equation*}
	% or equivalently, $\eta(\bm{x})=\log\theta(\bm{x})$.
	% Thus, $\eta(\bm{x})$ is unrestricted on the real line, while the
	% exponential link guarantees
	% $\theta(\bm{x})>0$ without requiring an additional positivity
	% constraint.
}

{The count law converts the intensity into a population survival function and a cure fraction through its generating function. Writing $\theta_i=\theta(\bm{x}_i)$, Poisson counts have generating function
\begin{equation*}
A_i(u)=\exp\{\theta_i(u-1)\},
\end{equation*}
so that substituting $u=S_{GG}(t)$ in Equation~\eqref{eq:promotion-population} for the population survival and $u=0$ for the cure fraction gives
\begin{equation}
\begin{aligned}
S_{\mathrm{P}}(t\mid\bm{x})
&=\exp\{-\theta(\bm{x})F_{GG}(t)\},\\
p_{c,\mathrm{P}}(\bm{x})
&=\exp\{-\theta(\bm{x})\},
\end{aligned}
\label{eq:poisson-survival-cure}
\end{equation}
with population density and hazard
\begin{equation*}
\begin{aligned}
f_{\mathrm{pop}}(t\mid\bm{x})
&=\theta(\bm{x})f_{GG}(t)\exp\{-\theta(\bm{x})F_{GG}(t)\},\\
h_{\mathrm{pop}}(t\mid\bm{x})
&=\theta(\bm{x})f_{GG}(t),
\end{aligned}
\end{equation*}
so the Poisson population hazard is the latency density scaled by the intensity, and $\mathrm{Var}(N_i\mid\bm{x}_i)=\theta_i$. Negative-binomial counts with dispersion $\phi>0$ have generating function
\begin{equation*}
A_i(u)=\{1+\phi\theta_i(1-u)\}^{-1/\phi},
\end{equation*}
with mean $\theta_i$ and variance $\theta_i(1+\phi\theta_i)$, so $\phi$ controls overdispersion relative to the Poisson case; the same substitutions give
\begin{equation}
\begin{aligned}
S_{\mathrm{NB}}(t\mid\bm{x})
&=\{1+\phi\theta(\bm{x})F_{GG}(t)\}^{-1/\phi},\\
p_{c,\mathrm{NB}}(\bm{x})
&=\{1+\phi\theta(\bm{x})\}^{-1/\phi},
\end{aligned}
\label{eq:nb-survival-cure}
\end{equation}
with population density and hazard
\begin{equation*}
\begin{aligned}
f_{\mathrm{pop}}(t\mid\bm{x})
&=\theta(\bm{x})f_{GG}(t)\{1+\phi\theta(\bm{x})F_{GG}(t)\}^{-1/\phi-1},\\
h_{\mathrm{pop}}(t\mid\bm{x})
&=\frac{\theta(\bm{x})f_{GG}(t)}{1+\phi\theta(\bm{x})F_{GG}(t)}.
\end{aligned}
\end{equation*}
The negative-binomial law reduces to the Poisson law as $\phi\to0$, so that $S_{\mathrm{NB}}\to S_{\mathrm{P}}$ and $p_{c,\mathrm{NB}}\to p_{c,\mathrm{P}}$, giving a direct check on the two count models. The cure fraction is strictly positive at every finite $\theta(\bm{x})$ and approaches one only as $\theta(\bm{x})\to0$.}

\subsection{Estimation Procedures}\label{sec:penalized-estimation}

Estimation of the GG-KM model is based on the EM algorithm, with
regularization incorporated through a penalized EM procedure. The
Poisson specification is estimated using penalized EM, whereas the
negative-binomial specification is estimated using the
Expectation/Conditional Maximization Either (ECME) algorithm
\cite{mclachlan2008algorithm}. The latter is particularly useful for
the negative-binomial model because the additional dispersion parameter
$\phi$ leads to a more complicated complete-data maximization step.

Let
$
    \mathcal D_{\mathrm{comp}}
    =
    \{(t_i,\delta_i,N_i,\bm{x}_i)\}_{i=1}^{n}
$
denote the corresponding complete data obtained by augmenting the
observations with the latent number of causes $N_i$. The parameter
vectors for the Poisson and negative-binomial models are respectively
$
    \bm\Theta_{\mathrm P}
    =
    (\beta_0,\bm{\alpha},a,d,p)
$
and
$
    \bm\Theta_{\mathrm{NB}}
    =
    (\beta_0,\bm{\alpha},a,d,p,\phi).
$
For both specifications,
$
    \eta_i=
    \eta(\bm{x}_i)
    =
    \beta_0+\bm K_i^{\mathsf T}\bm{\alpha},
    \text{ }
    \theta_i=\exp(\eta_i),
$
where
$
    \bm K_i
    =
    \bigl(
        K(\bm{x}_i,\bm{x}_1),
        \ldots,
        K(\bm{x}_i,\bm{x}_n)
    \bigr)^{\mathsf T}.
$

The target criterion is the observed penalized log-likelihood, obtained
by combining the count distribution, the generalized gamma latency
distribution in Equation~\eqref{eq:gg-distribution}, the kernel
predictor in Equation~\eqref{eq:rkhs-predictor}, and the RKHS penalty.
For the Poisson model,
\begin{equation}
    \ell^{\mathrm{pen}}_{\mathrm P}
    (\bm\Theta_{\mathrm P})
    =
    \sum_{i=1}^{n}
    \left[
        \delta_i
        \{\eta_i+\log f_{GG}(t_i)\}
        -
        \theta_iF_{GG}(t_i)
    \right]
    -
    \frac{\lambda}{2}
    \bm{\alpha}^{\mathsf T}
    \bm K
    \bm{\alpha}.
    \label{eq:poisson-observed-likelihood}
\end{equation}
For the negative-binomial model,
\begin{equation}
    \begin{aligned}
        \ell^{\mathrm{pen}}_{\mathrm{NB}}
        (\bm\Theta_{\mathrm{NB}})
        ={}&
        \sum_{i=1}^{n}
        \left[
            \delta_i
            \{\eta_i+\log f_{GG}(t_i)\}
        \right.\\
        &\left.
            \qquad-
            \left(
                \frac{1}{\phi}+\delta_i
            \right)
            \log
            \{1+\phi\theta_iF_{GG}(t_i)\}
        \right]
        -
        \frac{\lambda}{2}
        \bm{\alpha}^{\mathsf T}
        \bm K
        \bm{\alpha},
    \end{aligned}
    \label{eq:nb-observed-likelihood}
\end{equation}
with $\phi>0$. Here and throughout this subsection,
$f_{GG}(t_i)$, $F_{GG}(t_i)$, and $S_{GG}(t_i)$ are evaluated at the
current generalized gamma parameters $(a,d,p)$.

Direct optimization of these observed-data criteria is complicated by
the latent counts $N_i$. The EM construction instead uses the
conditional expectation of the complete-data log-likelihood. Given the
current parameter estimate $\bm\Theta^{(r)}$, define
\begin{equation}
    Q_{\lambda}
    \left(
        \bm\Theta
        \mid
        \bm\Theta^{(r)}
    \right)
    =
    \mathbb E_{\bm\Theta^{(r)}}
    \left[
        \ell_c
        \left(
            \bm\Theta;
            \mathcal D_{\mathrm{comp}}
        \right)
        \,\middle|\,
        \mathcal D_{\mathrm{obs}}
    \right]
    -
    \frac{\lambda}{2}
    \bm{\alpha}^{\mathsf T}
    \bm K
    \bm{\alpha}.
    \label{eq:penalized-Q-function}
\end{equation}
The E-step therefore replaces functions of the unobserved $N_i$ by
their conditional expectations under the current parameter values.

For the Poisson specification, the complete-data log-likelihood,
up to terms that do not depend on the model parameters, is
\begin{equation}
    \begin{aligned}
        \ell_{c,\mathrm P}
        ={}&
        \sum_{i=1}^{n}
        \left\{
            N_i\eta_i-\theta_i
        \right\}\\
        &+
        \sum_{i=1}^{n}
        \left\{
            \delta_i\log f_{GG}(t_i)
            +
            (N_i-\delta_i)
            \log S_{GG}(t_i)
        \right\}.
    \end{aligned}
    \label{eq:poisson-complete-data}
\end{equation}
The omitted terms, such as those involving the factorial of $N_i$,
are constant with respect to $\bm\Theta_{\mathrm P}$ when the complete
data are held fixed.

At iteration $r$, the conditional expectation of the latent count for
the Poisson model is
\begin{equation}
    \widehat N_{i,\mathrm P}^{(r)}
    =
    \mathbb E
    \left[
        N_i
        \mid
        t_i,\delta_i,\bm{x}_i;
        \bm\Theta_{\mathrm P}^{(r)}
    \right]
    =
    \delta_i
    +
    \theta_i^{(r)}
    S_{GG}^{(r)}(t_i).
    \label{eq:poisson-estep}
\end{equation}
For the negative-binomial specification, the corresponding conditional
expectation is
\begin{equation}
    \widehat N_{i,\mathrm{NB}}^{(r)}
    =
    \delta_i
    +
    \frac{
        \bigl(1+\phi^{(r)}\delta_i\bigr)
        \theta_i^{(r)}
        S_{GG}^{(r)}(t_i)
    }{
        1+
        \phi^{(r)}
        \theta_i^{(r)}
        F_{GG}^{(r)}(t_i)
    }.
    \label{eq:nb-estep}
\end{equation}

Once the conditional counts have been computed, the expected
complete-data objective separates naturally into a predictor block and
a generalized gamma latency block. For notational simplicity, the
iteration index on $\widehat N_i$ is suppressed in the following
expressions.

For the Poisson model, the terms depending on
$(\beta_0,\bm{\alpha})$ are
\begin{equation}
    Q_{\eta,\mathrm P}
    =
    \sum_{i=1}^{n}
    \left\{
        \widehat N_i\eta_i
        -
        \exp(\eta_i)
    \right\}
    -
    \frac{\lambda}{2}
    \bm{\alpha}^{\mathsf T}
    \bm K
    \bm{\alpha}.
    \label{eq:poisson-predictor-block}
\end{equation}
For the negative-binomial model, conditional on the current value of
$\phi$, the corresponding predictor block is
\begin{equation}
    Q_{\eta,\mathrm{NB}}
    =
    \sum_{i=1}^{n}
    \left[
        \widehat N_i\eta_i
        -
        \left(
            \widehat N_i+\frac{1}{\phi}
        \right)
        \log\{1+\phi\exp(\eta_i)\}
    \right]
    -
    \frac{\lambda}{2}
    \bm{\alpha}^{\mathsf T}
    \bm K
    \bm{\alpha},
    \label{eq:nb-predictor-block}
\end{equation}
up to terms that do not depend on $(\beta_0,\bm{\alpha})$.

For fixed conditional counts, both count models share the same
generalized gamma latency block,
\begin{equation}
    Q_{GG}
    =
    \sum_{i=1}^{n}
    \left\{
        \delta_i\log f_{GG}(t_i)
        +
        (\widehat N_i-\delta_i)
        \log S_{GG}(t_i)
    \right\}.
    \label{eq:latency-block}
\end{equation}
Consequently, $(\beta_0,\bm{\alpha})$ appears only in the corresponding
predictor block, whereas $(a,d,p)$ appears only in
Equation~\eqref{eq:latency-block} once $\widehat N_i$ is held fixed.
This separation permits the predictor and latency parameters to be
updated in distinct maximization blocks without altering the underlying
penalized EM objective.

For the Poisson specification, these conditional maximizations
constitute the penalized EM update. The negative-binomial specification
contains the additional dispersion parameter $\phi$. Its complete-data
count contribution contains terms involving $\phi$ and gamma functions
of the latent count, which makes direct maximization of the expected
complete-data criterion with respect to $\phi$ considerably less
convenient. The ECME algorithm avoids this additional expectation by
updating $\phi$ through the observed-data log-likelihood rather than
through the corresponding expected complete-data block.

Specifically, after updating
$(\beta_0,\bm{\alpha})$ and $(a,d,p)$ conditionally on the current
dispersion, the ECME update is
\begin{equation}
    \phi^{(r+1)}
    =
    \underset{\phi>0}{\arg\max}
    \;
    \ell^{\mathrm{pen}}_{\mathrm{NB}}
    \left(
        \beta_0^{(r+1)},
        \bm{\alpha}^{(r+1)},
        a^{(r+1)},
        d^{(r+1)},
        p^{(r+1)},
        \phi
    \right).
    \label{eq:nb-ecme-phi-update}
\end{equation}
Because the RKHS penalty does not depend on $\phi$, this update is
equivalent to maximizing the unpenalized observed-data log-likelihood
with respect to $\phi$ while all remaining parameters are held fixed.

The Supporting Information provides the complete-data derivations,
the conditional expectations in
Equations~\eqref{eq:poisson-estep}--\eqref{eq:nb-estep}, and the
analytical gradients used to optimize the predictor and latency
blocks.

\begin{algorithm*}[tbp]
	\caption{Penalized EM algorithm for the Poisson GG-KM model}
	\label{alg:poisson-em}
	\begin{algorithmic}[1]

		\Require Training data $(\bm X,\bm t,\bm\delta)$, kernel configuration,
		$\lambda>0$, tolerance $\varepsilon>0$, and maximum number of
		iterations $k_{\max}$

		\State Construct $\bm K$ and initialize
		$(\beta_0^{(0)},\bm\alpha^{(0)},\bm\xi^{(0)})$, with
		$\bm\xi=(\log a,\log d,\log p)^\top$; evaluate
		$\ell_{\mathrm P}^{\mathrm{pen},(0)}$.

		\For{$k=0,\ldots,k_{\max}-1$}

			\State Set
			$\eta_i^{(k)}
			=\beta_0^{(k)}+\bm K_i^\top\bm\alpha^{(k)}$ and
			$(a^{(k)},d^{(k)},p^{(k)})=\exp\{\bm\xi^{(k)}\}$.

			\State \textbf{E-step:} for $i=1,\ldots,n$, compute
			\[
			\widehat N_i^{(k+1)}
			=
			\delta_i+
			\exp\{\eta_i^{(k)}\}
			S_{GG}\!\left(
				t_i;a^{(k)},d^{(k)},p^{(k)}
			\right).
			\]

			\State \textbf{M-step:} with $\widehat N_i^{(k+1)}$ fixed, update
			\[
			(\beta_0^{(k+1)},\bm\alpha^{(k+1)})
			=
			\arg\max_{\beta_0,\bm\alpha}
			\left\{
			\sum_{i=1}^{n}
			\left[
				\widehat N_i^{(k+1)}\eta_i-e^{\eta_i}
			\right]
			-
			\frac{\lambda}{2}
			\bm\alpha^\top\bm K\bm\alpha
			\right\}.
			% \qquad
			% \eta_i=\beta_0+\bm K_i^\top\bm\alpha.
			\]

			\State Update the latency parameters on the log scale:
			\[
			\bm\xi^{(k+1)}
			=
			\arg\max_{\bm\xi}
			\sum_{i=1}^{n}
			\left[
				\delta_i
				\log f_{GG}\!\left(t_i;e^{\xi_a},e^{\xi_d},e^{\xi_p}\right)
				+
				(\widehat N_i^{(k+1)}-\delta_i)
				\log S_{GG}\!\left(t_i;e^{\xi_a},e^{\xi_d},e^{\xi_p}\right)
			\right].
			\]

			\State Set
			$(a^{(k+1)},d^{(k+1)},p^{(k+1)})
			=\exp\{\bm\xi^{(k+1)}\}$ and evaluate
			$\ell_{\mathrm P}^{\mathrm{pen},(k+1)}$.

			\If{
			$\left|
			\ell_{\mathrm P}^{\mathrm{pen},(k+1)}
			-
			\ell_{\mathrm P}^{\mathrm{pen},(k)}
			\right|
			\leq\varepsilon$
			}
				\State \textbf{break}
			\EndIf

		\EndFor

		\State \Return
		$(\widehat\beta_0,\widehat{\bm\alpha},
		\widehat a,\widehat d,\widehat p)$.

	\end{algorithmic}
\end{algorithm*}

\begin{algorithm*}[tbp]
	\caption{Penalized ECME algorithm for the negative-binomial GG-KM model}
	\label{alg:nb-ecme}
	\begin{algorithmic}[1]

		\Require Training data $(\bm X,\bm t,\bm\delta)$, kernel configuration,
		$\lambda>0$, tolerance $\varepsilon>0$, and maximum number of
		iterations $k_{\max}$

		\State Construct $\bm K$ and initialize
		$(\beta_0^{(0)},\bm\alpha^{(0)},\bm\xi^{(0)},\zeta^{(0)})$, with
		$\bm\xi=(\log a,\log d,\log p)^\top$ and $\zeta=\log\phi$;
		evaluate $\ell_{\mathrm{NB}}^{\mathrm{pen},(0)}$.

		\For{$k=0,\ldots,k_{\max}-1$}

			\State Set
			$\eta_i^{(k)}
			=\beta_0^{(k)}+\bm K_i^\top\bm\alpha^{(k)}$,
			$(a^{(k)},d^{(k)},p^{(k)})=\exp\{\bm\xi^{(k)}\}$, and
			$\phi^{(k)}=e^{\zeta^{(k)}}$.

			\State \textbf{E-step:} for $i=1,\ldots,n$, compute
			\[
			\widehat N_i^{(k+1)}
			=
			\delta_i+
			\frac{
				(1+\phi^{(k)}\delta_i)
				e^{\eta_i^{(k)}}
				S_{GG}(t_i;a^{(k)},d^{(k)},p^{(k)})
			}{
				1+\phi^{(k)}
				e^{\eta_i^{(k)}}
				F_{GG}(t_i;a^{(k)},d^{(k)},p^{(k)})
			}.
			\]

			\State \textbf{M-step:} with $\widehat N_i^{(k+1)}$ and
			$\phi^{(k)}$ fixed, update
			\[
			(\beta_0^{(k+1)},\bm\alpha^{(k+1)})
			=
			\arg\max_{\beta_0,\bm\alpha}
			\left\{
			\sum_{i=1}^{n}
			\left[
				\widehat N_i^{(k+1)}\eta_i
				-
				\left(
					\widehat N_i^{(k+1)}
					+\frac{1}{\phi^{(k)}}
				\right)
				\log\{1+\phi^{(k)}e^{\eta_i}\}
			\right]
			-
			\frac{\lambda}{2}
			\bm\alpha^\top\bm K\bm\alpha
			\right\}.
			% \qquad
			% \eta_i=\beta_0+\bm K_i^\top\bm\alpha.
			\]

			\State Update the latency parameters on the log scale:
			\[
			\bm\xi^{(k+1)}
			=
			\arg\max_{\bm\xi}
			\sum_{i=1}^{n}
			\left[
				\delta_i
				\log f_{GG}(t_i;e^{\xi_a},e^{\xi_d},e^{\xi_p})
				+
				(\widehat N_i^{(k+1)}-\delta_i)
				\log S_{GG}(t_i;e^{\xi_a},e^{\xi_d},e^{\xi_p})
			\right].
			\]

			\State \textbf{ECME-step:} update the dispersion on the log scale
			by maximizing the observed penalized log-likelihood:
			\[
			\zeta^{(k+1)}
			=
			\arg\max_{\zeta\in\mathbb R}
			\ell_{\mathrm{NB}}^{\mathrm{pen}}
			\left(
			    \beta_0^{(k+1)},
			    \bm\alpha^{(k+1)},
			    e^{\bm\xi^{(k+1)}},
			    e^\zeta
			\right).
			\]

			\State Set
			$(a^{(k+1)},d^{(k+1)},p^{(k+1)})
			=\exp\{\bm\xi^{(k+1)}\}$ and
			$\phi^{(k+1)}=e^{\zeta^{(k+1)}}$;
			evaluate $\ell_{\mathrm{NB}}^{\mathrm{pen},(k+1)}$.

			\If{
			$\left|
			\ell_{\mathrm{NB}}^{\mathrm{pen},(k+1)}
			-
			\ell_{\mathrm{NB}}^{\mathrm{pen},(k)}
			\right|
			\leq\varepsilon$
			}
				\State \textbf{break}
			\EndIf

		\EndFor

		\State \Return
		$(\widehat\beta_0,\widehat{\bm\alpha},
		\widehat a,\widehat d,\widehat p,\widehat\phi)$.

	\end{algorithmic}
\end{algorithm*}

Algorithms~\ref{alg:poisson-em} and \ref{alg:nb-ecme} summarize the
penalized EM and ECME procedures for the Poisson and negative-binomial
GG-KM models, respectively.

The simulations and application developed throughout this work use different
families of kernels to represent distinct forms of similarity between
covariate vectors. Let $\bm{x},\bm{x}'\in\mathcal X\subseteq\mathbb R^r$.
The kernel families considered are
\begin{equation}
	\begin{aligned}
		K_{\mathrm{lin}}(\bm{x},\bm{x}')
		&=\bm{x}^{\mathsf T}\bm{x}',\\
		K_{\mathrm{poly}}(\bm{x},\bm{x}')
		&=(\gamma\,\bm{x}^{\mathsf T}\bm{x}'+c_0)^q,\\
		K_{\mathrm{rbf}}(\bm{x},\bm{x}')
		&=\exp\{-\gamma\|\bm{x}-\bm{x}'\|_2^2\},\\
		K_{\mathrm{exp}}(\bm{x},\bm{x}')
		&=\exp\{-\gamma\|\bm{x}-\bm{x}'\|_2\},\\
		K_{\mathrm{lap}}(\bm{x},\bm{x}')
		&=\exp\{-\gamma\|\bm{x}-\bm{x}'\|_1\},\\
		K_{\mathrm{cau}}(\bm{x},\bm{x}')
		&=\{1+\gamma\|\bm{x}-\bm{x}'\|_2^2\}^{-1},\\
		K_{\mathrm{wav}}(\bm{x},\bm{x}')
		&=
		\prod_{j=1}^{r}
		\left[
		\cos\!\left(
		1.75\,\frac{x_j-x_j'}{A}
		\right)
		\exp\!\left\{
		-\frac{1}{2}
		\left(
		\frac{x_j-x_j'}{A}
		\right)^2
		\right\}
		\right].
	\end{aligned}
	\label{eq:kernels}
\end{equation}
Here, $\gamma>0$ is a kernel scale parameter, $q$ is the polynomial degree,
$c_0$ is the polynomial offset, and $A>0$ controls the scale of the
wavelet kernel. The RBF kernel corresponds to the Gaussian kernel and
depends on the squared Euclidean distance, whereas the exponential and
Laplacian kernels use the Euclidean and $L_1$ distances, respectively.
Kernel-specific hyperparameters are selected within the corresponding
tuning procedure.

The GG-KM model is implemented in the open-source Python package \texttt{scikit-cure},
available at \href{https://github.com/cowvin0/scikit-cure}{GitHub}.

\subsection{Prediction, tuning, and evaluation}
\label{sec:prediction-evaluation}

Prediction for an unseen covariate vector uses the training expansion
$
\eta(\bm{x}_{\mathrm{new}})
=
\widehat\beta_0
+
K(\bm{x}_{\mathrm{new}},\bm{X}_{\mathrm{train}})
\widehat{\bm{\alpha}}.
$
Equations~\eqref{eq:poisson-survival-cure} and
\eqref{eq:nb-survival-cure} then produce the individual cure fraction
and population survival curve. Covariate preprocessing is estimated
from the training data and applied unchanged to validation and test
data.

Predictive evaluation uses common event-stratified outer splits with
$70\%$ of observations for training and $30\%$ for testing. Kernel
parameters, the penalty, model variants, and all preprocessing choices
are selected inside the outer training set by a shared inner tuning
protocol and comparable search budget. The same outer test observations
and evaluation grid are used for every paired method comparison.

Let $\widehat G$ be the Kaplan--Meier estimate of the censoring survival
distribution fitted from the outer training data, with censoring treated
as the event. The weighting convention follows the prediction-error
construction of Graf and colleagues.\cite{graf1999} For a test set of
size $m$, the inverse-probability-of-censoring weighted Brier score is
\begin{equation}
	\begin{aligned}
		\widehat{\mathrm{BS}}(t)
		=
		\frac{1}{m}\sum_{i=1}^{m}
		\Biggl\{
		&\frac{I(t_i\leq t,\delta_i=1)}
		{\widehat G(t_i-)}
		\widehat S_i(t)^2\\
		+{}&
		\frac{I(t_i>t)}
		{\widehat G(t)}
		[1-\widehat S_i(t)]^2
		\Biggr\}.
	\end{aligned}
	\label{eq:ipcw-brier}
\end{equation}
The integration horizon is selected from the outer training data and
restricted to times with reliable censoring support. The integrated
Brier score is the time average obtained by trapezoidal integration of
Equation~\eqref{eq:ipcw-brier} over the frozen grid and division by the
horizon length, without renormalizing by the realized sum of weights.

Discrimination with respect to long-term cure status is additionally
assessed using an imputation-based area under the receiver operating
characteristic curve for cure models, denoted by
$\mathrm{AUC}_{\mathrm{cure}}$ \cite{medina2024}. Let
\begin{equation}
	\widehat\pi_i
	=
	\lim_{t\to\infty}
	\widehat S_i(t)
\end{equation}
denote the estimated cure probability for subject $i$. For the Poisson
GG-KM model,
$
\widehat\pi_i
=
\exp\{-\exp(\widehat\eta_i)\},
$
whereas for the negative-binomial GG-KM model,
$
\widehat\pi_i
=
\left(
1+\widehat\phi\exp(\widehat\eta_i)
\right)^{-1/\widehat\phi}
$.
Because cure status is not directly observed for censored individuals,
$\widehat\pi_i$ and $1-\widehat\pi_i$ are used as imputation weights for
the cured and susceptible states, respectively. For a cut-off
$c\in[0,1]$, the estimated true-positive and false-positive rates are
\begin{equation}
	\widehat{\mathrm{TPR}}(c)
	=
	\frac{
		\displaystyle
		\sum_{i=1}^{m}
		I(\widehat\pi_i\leq c)
		(1-\widehat\pi_i)
	}{
		\displaystyle
		\sum_{i=1}^{m}
		(1-\widehat\pi_i)
	},
	\label{eq:auc-cure-tpr}
\end{equation}
\begin{equation}
	\widehat{\mathrm{FPR}}(c)
	=
	\frac{
		\displaystyle
		\sum_{i=1}^{m}
		I(\widehat\pi_i\leq c)
		\widehat\pi_i
	}{
		\displaystyle
		\sum_{i=1}^{m}
		\widehat\pi_i
	}.
	\label{eq:auc-cure-fpr}
\end{equation}
Here, lower values of $\widehat\pi_i$ correspond to stronger evidence
that an individual is susceptible, so susceptibility is treated as the
positive class. The $\mathrm{AUC}_{\mathrm{cure}}$ is obtained by
trapezoidal integration of the receiver operating characteristic curve
formed by Equations~\eqref{eq:auc-cure-tpr} and \eqref{eq:auc-cure-fpr}, that is
		$\{\widehat{\mathrm{FPR}}(c),\widehat{\mathrm{TPR}}(c)\}$
over $c\in[0,1]$. Larger values indicate better discrimination between
individuals with higher and lower probabilities of belonging to the
susceptible subpopulation.

Evaluation summaries retain paired score differences, cure-fraction
error in simulations, $\mathrm{AUC}_{\mathrm{cure}}$, calibration,
convergence, failures, and runtime. A common generalized gamma latency
is used throughout the reported simulation so that the generating
latency matches the fitted latency structure and predictor-shape
contrasts are not confounded with unmodelled latency heterogeneity.
Covariate-dependent latency was not evaluated in the reported scenarios
and remains a future misspecification stress test.

% \subsection{Prediction, tuning, and evaluation}\label{sec:prediction-evaluation}

% Prediction for an unseen covariate vector uses the training expansion $\eta(\bm{x}_{\mathrm{new}})=\widehat\beta_0+K(\bm{x}_{\mathrm{new}},\bm{X}_{\mathrm{train}})\widehat{\bm{\alpha}}$. Equations~\eqref{eq:poisson-survival-cure} and \eqref{eq:nb-survival-cure} then produce the individual cure fraction and population survival curve. Covariate preprocessing is estimated from the training data and applied unchanged to validation and test data.

% Predictive evaluation uses common event-stratified outer splits with $70\%$ of observations for training and $30\%$ for testing. Kernel parameters, the penalty, model variants, and all preprocessing choices are selected inside the outer training set by a shared inner tuning protocol and comparable search budget. The same outer test observations and evaluation grid are used for every paired method comparison.

% Let $\widehat G$ be the Kaplan--Meier estimate of the censoring survival distribution fitted from the outer training data, with censoring treated as the event. The weighting convention follows the prediction-error construction of Graf and colleagues.\cite{graf1999} For a test set of size $m$, the inverse-probability-of-censoring weighted Brier score is
% \begin{equation}
	% \begin{aligned}
		% \widehat{\mathrm{BS}}(t)=\frac{1}{m}\sum_{i=1}^{m}
		% \Biggl\{
		% &\frac{I(t_i\leq t,\delta_i=1)}{\widehat G(t_i-)}
		% \widehat S_i(t)^2\\
		% +{}&\frac{I(t_i>t)}{\widehat G(t)}
		% [1-\widehat S_i(t)]^2
		% \Biggr\}.
		% \end{aligned}
	% \label{eq:ipcw-brier}
	% \end{equation}
% The integration horizon is selected from the outer training data and restricted to times with reliable censoring support. The integrated Brier score is the time average obtained by trapezoidal integration of Equation~\eqref{eq:ipcw-brier} over the frozen grid and division by the horizon length, without renormalizing by the realized sum of weights.

% Evaluation summaries retain paired score differences, cure-fraction error in simulations, calibration, convergence, failures, and runtime. A common generalized gamma latency is used throughout the reported simulation so that the generating latency matches the fitted latency structure and predictor-shape contrasts are not confounded with unmodelled latency heterogeneity. Covariate-dependent latency was not evaluated in the reported scenarios and remains a future misspecification stress test.

% \section{Simulation study}\label{sec:simulation}

The simulation study was designed to evaluate the predictive performance
of the proposed GG-KM model under different cure mechanisms, latent-count
distributions, and sample sizes. Specifically, we investigated three
questions. First, we examined whether a simpler predictor performs best
when its functional form is correctly specified. Second, we evaluated
whether GG-KM provides robust predictions when the true cure-probability
function is not adequately represented by the simpler models. Third, we
assessed the extent to which predictive performance depends on the
latent-count distribution and its degree of overdispersion. The
generalized gamma latency distribution was correctly specified and held
fixed across all scenarios. Therefore, differences in predictive
performance primarily reflect the ability of each predictor to recover
the covariate-dependent cure mechanism.

Each data-generating scenario was defined by one cure-probability surface,
one latent-count distribution, and one sample size. We considered three
cure-probability surfaces, linear, quadratic, and periodic. We also
considered three latent-count distributions, Poisson, negative binomial
with dispersion $\phi=0.25$, and negative binomial with dispersion
$\phi=1.00$. The two negative-binomial settings represent modest and
stronger overdispersion, respectively, whereas the Poisson distribution
represents the equidispersed limiting case. The sample sizes were
$n=300$ and $n=600$, resulting in $18$ data-generating scenarios.

In every scenario, two covariates were generated independently as
$x_1,x_2\sim\operatorname{Uniform}(-1,1)$. The bounded uniform domain
deliberately differs from the historical source because it gives all
three geometries a common domain for level-set and recovery comparisons.
The source used standard-normal covariates and a covariate-dependent
Weibull latency; retaining those choices would confound cure-predictor
shape with unmodelled latency heterogeneity.

The true cure probability was generated from one of three fixed surfaces.
Specifically, with $\operatorname{expit}(v)=\{1+\exp(-v)\}^{-1}$, the
designs were
\begin{equation}
	\begin{aligned}
		p_c^{\mathrm{lin}}(\bm{x})
		&=\operatorname{expit}(0.15-1.00x_1+0.75x_2),\\
		p_c^{\mathrm{quad}}(\bm{x})
		&=\operatorname{expit}(0.60-1.45x_1^2-1.05x_2^2+0.55x_1x_2),\\
		p_c^{\mathrm{per}}(\bm{x})
		&=\operatorname{expit}\{0.10+1.15\sin(\pi x_1)+0.80\cos(\pi x_2)\}.
	\end{aligned}
	\label{eq:simulation-cure-surfaces}
\end{equation}
For Poisson generation, the latent-count mean was
$\theta(\bm{x})=-\log p_c(\bm{x})$; for negative-binomial generation, it
was $\theta(\bm{x})=\{p_c(\bm{x})^{-\phi}-1\}/\phi$. This mapping
preserves the same cure-probability surface across count distributions,
so comparisons do not confound dispersion with a different cure target.

Figure~\ref{fig:cure-surfaces} shows the three surfaces on the common
covariate square. The linear surface increases along one tilted gradient
and has parallel straight level sets, so the log-linear predictor is
correctly specified for it. The quadratic surface has closed elliptical
level sets generated by the two squared terms and the interaction, and no
additive predictor reproduces that geometry. The periodic surface reverses
its gradient twice inside the domain and produces alternating regions of
high and low cure probability, a shape an additive spline can follow one
coordinate at a time but a degree-two map cannot. The three designs
therefore separate correct linear specification, smooth nonadditivity and
oscillation while holding the latency distribution fixed. They are not
matched on prevalence: the mean of $p_c(\bm{x})$ over the square is
0.533 for the linear surface, 0.448 for the quadratic surface and 0.520
for the periodic surface, with attainable ranges $[0.168,0.870]$,
$[0.079,0.646]$ and $[0.136,0.886]$. Scores are consequently comparable
between methods within a scenario and not between scenarios.

\GGKMWideFigure{cure-surfaces.pdf}{Frozen cure-probability surfaces used
in the simulation design. Panels show the linear, quadratic, and
periodic functions in Equation~\eqref{eq:simulation-cure-surfaces} over
$(x_1,x_2)\in[-1,1]^2$. All panels use the same $p_c(\bm{x})$ colour
scale; superimposed contour lines provide a colour-independent encoding
of the surface geometry. Contours are labelled at $p_c=0.25$, $0.50$ and
$0.60$, the same levels used in Figure~\ref{fig:cure-surface-recovery},
and each panel title reports the mean of $p_c(\bm{x})$ over the square.
}{fig:cure-surfaces}{Three contour panels show cure probability over two
covariates. The linear panel has diagonal straight contours, the quadratic
panel has curved contours with an interaction, and the periodic panel has
repeated peaks and troughs. Each panel title states the mean cure
probability for that surface.
}

All latent-cause activation times followed the generalized gamma
distribution with $(a,d,p)=(1.6,1.3,1.1)$. Independent censoring times
were exponential with mean 3.5, and observed time was the minimum of the
event and censoring times. The inverse population distribution was used
to generate finite event times directly for susceptible individuals, while
individuals in the limiting cure mass were assigned no event. Every
replicate retained the true individual cure probability and population
survival curve for evaluation.

Algorithm~\ref{alg:simulation-generation} states the implemented
inverse-transform generator. Sampling from the improper population
distribution avoids simulating a potentially large latent count while
preserving its cure mass and finite-event distribution exactly.
\begin{algorithm*}[tbp]
\caption{Right-censored simulation from a promotion-time population distribution}\label{alg:simulation-generation}
\begin{algorithmic}[1]
\Require Scenario, sample size $n$, generalized gamma parameters $(a,d,p)$, censoring scale, and deterministic replication seed
\For{$i=1,\ldots,n$}
  \State Draw $x_{i1},x_{i2}\stackrel{\mathrm{ind}}{\sim}\operatorname{Uniform}(-1,1)$ and evaluate the declared $p_c(\bm{x}_i)$
  \State Set $\theta_i=-\log p_c(\bm{x}_i)$ for Poisson counts, or $\theta_i=\{p_c(\bm{x}_i)^{-\phi}-1\}/\phi$ for negative-binomial counts
  \State Draw $U_i\sim\operatorname{Uniform}(0,1)$ from the population event-time distribution
  \If{$U_i\geq1-p_c(\bm{x}_i)$}
    \State Set $T_i=\infty$ (the draw falls in the cure mass)
  \Else
    \State Set $q_i=-\log(1-U_i)/\theta_i$ for Poisson counts
    \State Set $q_i=\{(1-U_i)^{-\phi}-1\}/(\phi\theta_i)$ for negative-binomial counts
    \State Set $T_i=F_{GG}^{-1}(q_i;a,d,p)$
  \EndIf
  \State Draw $C_i$ from the declared independent censoring distribution
  \State Record the observed time $t_i=\min(T_i,C_i)$, indicator $\delta_i=I(T_i\leq C_i)$, $p_c(\bm{x}_i)$, and the true population survival grid
\EndFor
\State Return the complete observed sample, generating truth, scenario identifier, and replication seed
\end{algorithmic}
\end{algorithm*}

The fitted comparison set shared the same generalized gamma latency and comprised
GG-KM predictors with linear, radial-basis, Laplacian, exponential, Cauchy,
polynomial, and wavelet kernels, each with Poisson and negative-binomial latent
counts, together with additive cubic-spline Poisson and negative-binomial models
and the classical mixture-cure and promotion-time comparators. Linear-kernel
candidates used $\lambda\in\{0.3,3\}$. Radial-basis, Laplacian, exponential, and
Cauchy candidates used $\gamma\in\{0.25,1\}$ with
$\lambda\in\{0.3,3\}$. Polynomial candidates additionally used
degree $q\in\{2,3\}$ with $c_0=1$, while wavelet candidates used
$A\in\{0.5,1,2\}$, in both cases with $\lambda\in\{0.3,3\}$.
Spline candidates used three or five training-fold quantile knots with
$\lambda\in\{0.3,3\}$. All tuned methods used the same two-fold inner tuning
criterion and evaluation-grid size. The promotion-time support-vector method of
Pal and Aselisewine was not included because the reproducibility audit found no
public executable implementation at the article or arXiv records, and a new
reimplementation would not provide a controlled software comparison.
\cite{pal2023}

Each generated sample was split once into a common event-stratified
$70\%$ training set and $30\%$ test set for all fitted methods. Tuning,
feature-map fitting, and preprocessing were confined to the training set,
and the IPCW integration horizon was set from that set at the smaller of
its 90th observed-time percentile and the limit at which estimated
censoring survival remained at least 0.05. Primary test-set estimands were
integrated Brier score, cure AUC, and cure-probability mean squared error.

A target-size diagnostic pilot used five planned replicates per pilot
scenario and was not pooled with the final experiment; at least four
replicates were successful in every pilot cell. The largest observed
between-replicate standard deviations were 0.0418 for integrated Brier
score and 0.0332 for cure-probability mean squared error. Requiring
projected Monte Carlo standard errors no larger than 0.003 and 0.005,
respectively, implied at least 195 replicates for integrated Brier score
and 45 for cure-probability mean squared error. The more restrictive
criterion therefore required 195 replicates, so the final count was fixed
at 250 replicates per scenario before final estimation. At
$R=250$, the corresponding projected Monte Carlo standard errors were
0.00265 for integrated Brier score and 0.00210 for cure-probability mean
squared error. For cure AUC, the largest pilot standard deviation was
0.1211, corresponding to a projected Monte Carlo standard error of
0.00766 at $R=250$. Final summaries report the realized Monte Carlo
standard error for every primary metric and paired differences computed
within each shared replication.

\subsection{Simulation results}\label{sec:simulation-results}

The primary final experiment contained 81,000 method rows from
250 replications in each of 18 scenarios. The largest realized
Monte Carlo standard errors were 0.00175 for the IPCW integrated
Brier score, 0.00434 for cure-probability RMSE, and
0.00694 for cure AUC. Of the 81,000 fits, 80,842 completed
successfully and 158 failed; among the successful fits, 80,832
satisfied the explicit convergence criterion and 10 did not.
Method-specific successful-fit totals ranged from 4,467 to 4,498
of 4,500. Mean per-fit runtime ranged from 1.13 to 40.32 seconds
across methods, with a maximum of 520.48 seconds.

Table~\ref{tab:simulation-best} reports the best observed mean IPCW
integrated Brier score, cure root mean squared error, and cure AUC in each
scenario. For the linear surface, a linear-kernel GG-KM attained the lowest
mean IPCW integrated Brier score and cure root mean squared error in all six
scenarios. For the quadratic surface, a polynomial-kernel GG-KM attained the
lowest cure root mean squared error in all six scenarios and the lowest
integrated Brier score in five of six scenarios, with an RBF-kernel GG-KM
attaining the lowest integrated Brier score in the remaining case. For the
periodic surface, an additive spline model attained the lowest integrated
Brier score and cure root mean squared error in all six scenarios. For cure
AUC, a polynomial-kernel GG-KM attained the highest mean in all six linear
scenarios and all six periodic scenarios. On the quadratic surface,
polynomial-kernel GG-KM models attained the highest mean cure AUC in four
of six scenarios, whereas wavelet-kernel GG-KM models attained the highest
mean in the remaining two.

\GGKMTable{sections/generated/simulation-best-table}

Figure~\ref{fig:simulation-paired-differences} shows paired differences
between the RBF GG-KM and the count-matched linear-kernel GG-KM,
additive-spline, and polynomial-kernel GG-KM comparators across the six
generating-law and sample-size cells. Each contrast is defined as the RBF
GG-KM value minus the comparator value; hence, negative differences favor
the RBF model for IPCW integrated Brier score and cure root mean squared
error, whereas positive differences favor it for cure AUC.

For the linear cure surface, the linear-kernel GG-KM consistently
outperformed the RBF model in integrated Brier score and cure root mean
squared error, while the spline and polynomial comparisons were generally
closer. Cure AUC also favored the linear and polynomial kernels, whereas
the RBF--spline differences remained comparatively small. Under the
quadratic surface, the RBF GG-KM substantially improved on the linear
kernel across all three metrics, but its differences from the spline and
polynomial models were much smaller; the polynomial kernel generally had
lower cure root mean squared error. For the periodic surface, the RBF
model improved markedly on the linear and polynomial kernels in integrated
Brier score and cure root mean squared error, particularly at $n=600$,
but remained inferior to the spline. Cure AUC showed a different pattern,
with the RBF model clearly outperforming the linear kernel and, at
$n=600$, the Poisson polynomial-kernel GG-KM, while its differences from
the spline remained comparatively small. Overall, the paired comparisons
show that relative performance depends strongly on both the geometry of
the cure surface and the evaluation metric.

\GGKMWideFigure
{simulation_rbf_paired_differences_with_auc.pdf}
{Count-matched paired predictive differences of the RBF GG-KM against the linear-kernel GG-KM, additive-spline, and polynomial-kernel GG-KM comparators. Columns correspond to the linear, quadratic, and periodic cure surfaces, while rows show IPCW integrated Brier score, cure root mean squared error, and cure AUC differences. Each contrast is defined as the RBF GG-KM value minus the corresponding comparator value, with Poisson and negative-binomial fits compared separately. Error bars span $\pm1.96$ paired Monte Carlo standard errors, and a horizontal line marks zero. Negative differences favor the RBF GG-KM for integrated Brier score and cure root mean squared error, whereas positive differences favor it for cure AUC.}
{fig:simulation-paired-differences}
{
A three-row by three-column grid displays count-matched paired
differences between the RBF GG-KM and linear-kernel GG-KM,
additive-spline, and polynomial-kernel GG-KM comparators. Columns
represent linear, quadratic, and periodic cure surfaces; rows represent
integrated Brier score, cure root mean squared error, and cure AUC.
Poisson and negative-binomial contrasts are distinguished by marker and
line type, with error bars spanning $\pm1.96$ paired Monte Carlo standard
errors around a horizontal zero line. Negative values favor the RBF model
for integrated Brier score and cure root mean squared error, whereas
positive values favor it for cure AUC.
}

Figure~\ref{fig:cure-surface-recovery} maps where in the
covariate square each predictor loses accuracy. Every panel
fills the $41\times41$ evaluation grid with the pointwise root
mean squared error of the predicted cure probability against
the generating surface, averaged over the valid Poisson
replications at $n=600$, on a scale shared by all twelve panels.
Solid lines repeat the true level sets at $p_c=0.25$, $0.50$
and $0.60$, and the dashed line is the mean estimated
$p_c=0.50$ level, which separates an error of magnitude from a
displacement of shape. Failed fits reduce the pointwise
$n_{\mathrm{valid}}$ rather than being removed from the total
number of replications. The panel means summarize the gridwise
recovery of the cure surface. On the linear surface the
linear-kernel GG-KM attains 0.0490, compared with 0.0671 for
the additive spline, 0.0694 for the polynomial-kernel GG-KM,
and 0.0730 for the RBF GG-KM. On the quadratic surface, the
RBF GG-KM attains the lowest mean error, 0.0802, followed
closely by the additive spline at 0.0816 and the
polynomial-kernel GG-KM at 0.0827, whereas the linear-kernel
GG-KM increases to 0.1282. On the periodic surface, the
additive spline performs best at 0.0824, followed by the RBF
GG-KM at 0.0959, while the linear- and polynomial-kernel GG-KM
predictors increase to 0.1504 and 0.1510, respectively. The
spatial pattern adds what these scalar summaries omit. A
misaligned predictor does not spread its error evenly: the
linear-kernel GG-KM shows concentrated error on the quadratic
and periodic surfaces, and its estimated $p_c=0.50$ level is
absent on the quadratic surface because the fitted cure surface
does not reach that level. The polynomial-kernel GG-KM
similarly deteriorates on the periodic surface. In contrast,
the RBF GG-KM maintains comparatively stable recovery across
the three generating mechanisms relative to the linear and
polynomial kernels, with its estimated $p_c=0.50$ contour
closely following the true level wherever it is present.
Figure~\ref{fig:cure-surface-recovery} therefore complements
the regret summary in Table~\ref{tab:simulation-regret} by
showing where prediction errors occur spatially: no single
predictor dominates across all three surfaces, with the
linear kernel performing best on the linear surface, the RBF
kernel on the quadratic surface, and the additive spline on
the periodic surface.

\GGKMWideFigure{simulation_cure_surface_recovery.pdf}{Pointwise accuracy of the three primary cure surfaces at $n=600$ under Poisson generation. Rows show the linear, quadratic and periodic generating surfaces; columns show the linear-kernel GG-KM, additive-spline, polynomial-kernel GG-KM and radial-basis-function (RBF) GG-KM predictors. The filled field is the root mean squared error of the predicted cure probability at each of the $41\times41$ grid points, taken across all valid replications, on one scale shared by every panel and truncated above 0.30. Solid labelled lines are the true cure levels $0.25$, $0.50$ and $0.60$; the dashed line is the mean estimated $0.50$ level and is annotated as absent when the fitted range does not contain it. The value printed in each panel is the mean of the displayed field. Failed fits remain in $n_{\mathrm{total}}$ and reduce the pointwise $n_{\mathrm{valid}}$ rather than disappearing.}{fig:cure-surface-recovery}{A three-row by four-column grid of shaded panels shows the pointwise root mean squared error of the predicted cure probability for the linear, quadratic and periodic generating surfaces. Columns correspond to the linear-kernel GG-KM, additive-spline, polynomial-kernel GG-KM and RBF GG-KM predictors. Darker shading marks larger error, and the printed value in each panel is the mean of the displayed field. Solid labelled lines mark the true cure levels, while a dashed line marks the mean estimated 0.50 level and is explicitly reported as absent wherever that level is outside the fitted range.}

The complete mean IPCW integrated Brier scores for all eight methods,
with their Monte Carlo standard errors, appear in the Supporting
Information. They confirm the geometry-specific ordering in
Table~\ref{tab:simulation-best} and retain both fitted count laws for
every predictor.

To complement the scenario-specific comparisons, we
summarized performance using regret relative to the
best-performing method within the retained comparison set.
Table~\ref{tab:simulation-regret} restricts this analysis to
the Poisson and negative-binomial spline models and RBF GG-KM
models across the 18 scenarios. For IPCW integrated Brier score,
the Poisson spline had the smallest maximum paired regret,
0.00065 with a 95\% Monte Carlo half-width of 0.00056, whereas
the negative-binomial spline had the smallest equally weighted
mean regret, 0.00012 (half-width 0.00010). For cure root mean
squared error, the Poisson spline had the smallest maximum regret,
0.00464 (half-width 0.00210), whereas the negative-binomial spline
had the smallest mean regret, 0.00073 (half-width 0.00025). In
contrast, the negative-binomial RBF GG-KM was most competitive
for cure AUC, attaining the smallest maximum regret, 0.00687
(half-width 0.00464), the smallest mean regret, 0.00144
(half-width 0.00060), and the best descriptive mean rank of
2.06. Thus, the regret summaries favor the spline models for
prediction error but show greater competitiveness of the RBF
GG-KM for cure discrimination.

\GGKMTable{sections/generated/simulation-regret-table}

\subsection{Higher-dimensional interaction surfaces}\label{sec:highdim-simulation}

The four-covariate experiment is a secondary mechanism study developed
after the primary analysis. It probes interaction structures that
an additive spline cannot represent, but it is neither
independent confirmation nor evidence of a general kernel
advantage. Covariates satisfy
$x_1,\dots,x_4\sim\operatorname{Uniform}(-1,1)$, and the first
target is the \emph{interaction} surface
\begin{equation}
p_c(\bm{x})=\operatorname{expit}\bigl(C_{\mathrm{int}}+1.60\,x_1x_2-1.30\,x_3x_4+0.70\,x_1-0.50\,x_4\bigr),
\end{equation}
and the \emph{radial} surface is
\begin{equation}
p_c(\bm{x})=\operatorname{expit}\bigl(C_{\mathrm{rad}}+3.00\exp\{-0.80(x_1^2+x_2^2+x_3^2+x_4^2)\}\bigr).
\end{equation}
Intercepts $C_{\mathrm{int}}$ and $C_{\mathrm{rad}}$ were
calibrated by bisection on a fixed sample of $200\,000$
observations (seed 20260901) so that each surface yields a mean
cure fraction of $0.45$. The latency remained the common
generalized gamma with $(a,d,p)=(1.6,1.3,1.1)$ and no
covariate-dependent scale; censoring was exponential with mean
3.5. The same fitted methods as in the primary experiment were
compared under Poisson, negative-binomial ($\phi=0.25$), and
negative-binomial ($\phi=1.00$) generating laws at $n=300$ and
$n=600$, giving 12 scenarios. The final experiment used 250
replicates per scenario, and every declared cell remains in the
analysis regardless of its ordering.

Table~\ref{tab:highdim-best} and
Figure~\ref{fig:highdim-paired-differences} report the directed
comparisons with paired Monte Carlo uncertainty. Their
interpretation is limited to the evaluated interaction and
radial mechanisms, and it retains every declared cell,
including comparisons that do not favor GG-KM.

    \GGKMTable{sections/generated/highdim-best-table}
\GGKMWideFigure
{highdim_rbf_vs_linear_spline_polynomial_paired_differences_with_auc.pdf}
{Paired predictive differences of the radial-basis GG-KM predictor against the count-matched GG-KM linear-kernel, spline, and GG-KM polynomial-kernel comparators on the interaction and radial cure surfaces. The top row shows paired mean IPCW integrated Brier score differences and the bottom row paired mean cure root mean squared error differences. Each plotted contrast is defined as the RBF GG-KM value minus the corresponding comparator value, so negative differences favor the RBF GG-KM and positive differences favor the comparator. Solid lines with filled markers denote Poisson fits, whereas dashed lines with open markers denote negative-binomial fits. Each panel displays the six generating-law and sample-size cells, with error bars spanning $\pm1.96$ paired Monte Carlo standard errors.}
{fig:highdim-paired-differences}
{Two-row by two-column grid of paired-difference plots for the interaction and radial cure surfaces, comparing RBF GG-KM with GG-KM linear-kernel, spline, and GG-KM polynomial-kernel models under Poisson and negative-binomial latent-count specifications.}

The complete per-scenario means with Monte Carlo standard errors, together with the numerical diagnostics, are
reported in the Supporting Information.

% \section{Melanoma application}\label{sec:melanoma}

The empirical application used the malignant melanoma cohort distributed
as the \texttt{melanoma} data set in version 1.3-32 of the \texttt{boot}
R package and documented in the original counting-process source.
\cite{boot2025,andersen1993} The cohort contains 205 patients operated
on in Odense, Denmark between 1962 and 1977, with follow-up recorded through
1977. There were 57 deaths from melanoma, 14 deaths from other causes,
and 134 patients alive at the end of follow-up. The analysis treated
death from melanoma as the event of interest and censored both administrative
survivors and deaths from other causes, so the fitted model is cause
specific in its event definition and is not a competing-risks analysis;
$p_c(\bm{x})$ is the probability of no melanoma death from the modeled
process, not an all-cause survival probability.
The predictor set comprised sex, age at operation, calendar year of
operation, the natural logarithm of tumour thickness in millimetres, and
ulceration status. Follow-up time was converted from days to years, and
no observations or covariate values were missing. Continuous and binary
predictors passed through the same training-fold standardization step so
that no test-set information entered preprocessing. The versioned
analysis copy was confirmed row for row against \texttt{boot::melanoma}
and audited against its recorded SHA-256 checksum before estimation.

Predictive performance was assessed in 50 event-stratified outer holdouts,
each assigning 143 patients to training and 62 to testing. The same split
indices were used for every fitted method, and all tuning choices were made
within three inner training folds using the IPCW integrated Brier score
and the horizon rule defined in Section~\ref{sec:prediction-evaluation}.
The comparison comprised GG-KM predictors with linear, radial-basis,
Laplacian, exponential, Cauchy, polynomial, and wavelet kernels, each
with Poisson and negative-binomial latent-count laws, together with
additive cubic-spline Poisson and negative-binomial models and the
classical mixture-cure and promotion-time comparators. Kernel and spline
hyperparameters were selected from prespecified grids.

\GGKMTable{sections/generated/melanoma-performance-table}

Table~\ref{tab:melanoma-performance} reports split-level means and
variation. The Poisson and negative-binomial GG-KM models with the RBF
kernel attained the lowest mean IPCW integrated Brier score, both at
0.1418 with split standard deviation 0.0206, followed closely by the
Cauchy kernel under both count laws, with mean IBS 0.1419. The
mixture-cure model had the highest mean cure AUC overall, 0.8492, while
among the GG-KM models the negative-binomial polynomial kernel had the
highest mean cure AUC, 0.8117. Across the 900 planned split-model fits,
897 completed successfully and 893 satisfied the explicit convergence
criteria.

Figure~\ref{fig:melanoma-performance} places the paired score
distributions beside descriptive survival calibration. The comparison of
split-averaged predicted survival with the Kaplan--Meier curve is a
calibration check on the population survival function; it does not
identify individual cure status. Repeated holdouts are dependent because
the same patients recur across splits, so empirical quantiles describe
variation across training sets rather than independent-sample confidence
intervals. Similar mean scores do not establish an approximately linear
latent-cause structure, and a smaller split standard deviation alone does
not prove greater predictive stability.

\GGKMWideFigure{melanoma_predictive_assessment_ibs_auc.pdf}{Predictive assessment in 50 paired repeated holdouts of the melanoma cohort. Panel A compares the full-cohort Kaplan--Meier curve with split-averaged out-of-sample survival predictions over the common supported horizon. Panel B shows split-level IPCW integrated Brier scores, with each point representing one outer split and lower values indicating better prediction. Line type, marker or box encoding, and labels distinguish all eight methods without relying on color alone. Panel A shows mean curves without an uncertainty band; the dependence caveat for repeated holdouts is stated in the text.}{fig:melanoma-performance}{Two-panel melanoma prediction figure. The left panel shows a black Kaplan-Meier step curve and eight mean predicted survival curves from log-linear, spline, quadratic, and GG-KM models with Poisson or negative-binomial count laws. The right panel shows split-level IPCW integrated Brier scores as box-and-strip distributions for every method, with each point representing one outer split and lower values indicating better prediction.}
\GGKMWideFigure{melanoma_ggkm_kernel_comparison_ibs_auc.pdf}{Kernel-specific predictive performance of the GG-KM models in 50 paired repeated holdouts of the melanoma cohort. Panel A shows the split-level IPCW integrated Brier score and Panel B the cure-discrimination area under the curve ($\mathrm{AUC}_{\mathrm{cure}}$) across the seven candidate kernels. Within each kernel, Poisson and negative-binomial GG-KM results are shown side by side, with color and hatching distinguishing the two latent-count laws. Each point represents one outer split, and the boxplots summarize the corresponding split-level distributions.}{fig:melanoma-performance-ggkm-kernel}{Two-panel comparison of GG-KM kernels in the melanoma application. Panel A displays split-level IPCW integrated Brier scores and Panel B cure AUC values for the linear, RBF, Laplacian, exponential, Cauchy, polynomial, and wavelet kernels. For each kernel, Poisson and negative-binomial GG-KM boxplots are shown side by side and distinguished by color and hatching. Boxplots summarize the repeated-holdout distributions, with individual outer-split values overlaid as points.}

Mean predicted cure fractions ranged from 0.572 for the Poisson
GG-KM with the wavelet kernel to 0.639 for the Poisson GG-KM with the
polynomial kernel. These values are model-based extrapolations because
individual cure status is unobserved and follow-up is finite, and the
Kaplan--Meier comparison cannot identify a cure fraction. Across-split
prediction summaries were retained for every individual to expose
variation caused by repeated training samples rather than to create
surrogate cure labels. The application consequently evaluates
reproducible out-of-sample prediction and numerical behavior while
leaving biological interpretation of individual cure probabilities
outside its evidential scope.

Figure~\ref{fig:melanoma-performance-ggkm-kernel} first indicates a marked
similarity between the Poisson and negative-binomial GG-KM specifications.
For most kernels, the corresponding boxplots have strongly overlapping
distributions for both the IPCW integrated Brier score and
$\mathrm{AUC}_{\mathrm{cure}}$, and their mean performances are generally
close.
% This pattern provides empirical evidence, in this application, that
% the predictive behavior of the GG-KM model is relatively insensitive to the
% specific latent-count distribution used. It should not, however, be
% interpreted as establishing distributional independence of the model, since
% the comparison is limited to the Poisson and negative-binomial
% specifications and to this data set. Rather, the results suggest a degree of
% robustness to the latent-count specification, with the kernel choice having
% a more visible influence on predictive performance.

For the IPCW integrated Brier score (Panel A), the distributions overlap
substantially across kernels, and the differences between the Poisson and
negative-binomial versions of the same kernel are generally small. Mean IBS
values were relatively close across specifications, with both the Poisson
and negative-binomial RBF models attaining the lowest mean value, 0.1418.
The Poisson polynomial model had the largest mean IBS, 0.1515, and also
showed comparatively large between-split variability. Although several
kernels exhibited isolated splits with substantially higher Brier scores,
the central parts of the distributions remained similar. Overall, the IBS
results do not indicate a clearly dominant kernel.

Panel B shows greater sensitivity to the kernel choice for
$\mathrm{AUC}_{\mathrm{cure}}$, while still retaining substantial agreement
between the two latent-count laws within most kernels. The linear kernel
provided consistently strong cure discrimination under both count laws,
with mean cure AUCs of approximately 0.809. The negative-binomial
polynomial model achieved the highest mean cure AUC among the GG-KM
specifications, 0.8117, and showed a relatively concentrated distribution,
whereas the corresponding Poisson polynomial model was substantially more
variable across splits. The RBF, Laplacian, exponential, and Cauchy kernels
generally produced lower cure AUCs, concentrated around 0.68--0.70, while
the wavelet kernel showed the weakest discrimination, with mean cure AUCs
of approximately 0.62 under both count laws. Thus, the results suggest that
kernel specification has a stronger practical influence than the choice
between Poisson and negative-binomial latent-count laws in this application.
No single kernel dominated simultaneously in prediction error and cure
discrimination: the RBF kernel was most favorable in terms of mean IBS
under both count laws, whereas the polynomial and linear kernels provided
the strongest cure discrimination.

% \input{sections/05_discussion}
% \input{sections/06_declarations}

% \bibliography{references}

% \ifggkmsubmission
% \clearpage
% \begin{table*}[tbp]
\caption{Predictive performance in 50 paired repeated holdouts of the melanoma cohort.
Values for the IPCW integrated Brier score (IBS), cure-discrimination area under the curve
($\mathrm{AUC}_{\mathrm{cure}}$), and mean predicted cure fraction are means across successful
splits, with split standard deviations in parentheses. Convergence is reported as the number
of converged fits among successful fits. Runtime is the mean with the maximum in parentheses.
}
\label{tab:melanoma-performance}
\centering
\begin{tabular}{llccccc}
\hline
Predictor & Count law & Converged & IPCW IBS &
$\mathrm{AUC}_{\mathrm{cure}}$ & Mean cure prediction & Runtime (s) \\
\hline

Mixture cure
& Mixture cure
& 49/50
& 0.1430 (0.0185)
& 0.8492 (0.0491)
& 0.607 (0.079)
& 2.28 (12.39) \\

Promotion time
& Poisson
& 49/49
& 0.1456 (0.0239)
& 0.8240 (0.0283)
& 0.633 (0.047)
& 2.21 (7.29) \\

\hline

Spline
& Poisson
& 50/50
& 0.1432 (0.0217)
& 0.7912 (0.0290)
& 0.613 (0.049)
& 42.56 (99.65) \\

Spline
& Negative binomial
& 50/50
& 0.1434 (0.0216)
& 0.7911 (0.0291)
& 0.614 (0.049)
& 46.37 (112.34) \\

\hline

GG-KM (linear)
& Poisson
& 50/50
& 0.1426 (0.0224)
& 0.8091 (0.0274)
& 0.627 (0.046)
& 15.48 (37.21) \\

GG-KM (linear)
& Negative binomial
& 50/50
& 0.1426 (0.0224)
& 0.8091 (0.0274)
& 0.627 (0.046)
& 16.29 (38.53) \\

GG-KM (RBF)
& Poisson
& 50/50
& 0.1418 (0.0206)
& 0.6937 (0.0340)
& 0.586 (0.047)
& 45.85 (97.33) \\

GG-KM (RBF)
& Negative binomial
& 50/50
& 0.1418 (0.0206)
& 0.6937 (0.0341)
& 0.586 (0.047)
& 54.00 (184.25) \\

GG-KM (Laplacian)
& Poisson
& 48/49
& 0.1426 (0.0253)
& 0.6987 (0.0490)
& 0.578 (0.050)
& 47.26 (97.01) \\

GG-KM (Laplacian)
& Negative binomial
& 50/50
& 0.1438 (0.0257)
& 0.7021 (0.0526)
& 0.575 (0.062)
& 55.12 (140.77) \\

GG-KM (exponential)
& Poisson
& 49/50
& 0.1450 (0.0242)
& 0.6984 (0.0800)
& 0.580 (0.048)
& 43.01 (89.43) \\

GG-KM (exponential)
& Negative binomial
& 50/50
& 0.1448 (0.0241)
& 0.6949 (0.0768)
& 0.580 (0.049)
& 49.45 (110.53) \\

GG-KM (Cauchy)
& Poisson
& 50/50
& 0.1419 (0.0246)
& 0.6817 (0.0390)
& 0.585 (0.047)
& 46.38 (92.49) \\

GG-KM (Cauchy)
& Negative binomial
& 50/50
& 0.1419 (0.0246)
& 0.6823 (0.0390)
& 0.586 (0.046)
& 51.95 (173.17) \\

GG-KM (polynomial)
& Poisson
& 50/50
& 0.1515 (0.0274)
& 0.7457 (0.1338)
& 0.639 (0.063)
& 113.78 (232.80) \\

GG-KM (polynomial)
& Negative binomial
& 50/50
& 0.1477 (0.0280)
& 0.8117 (0.0261)
& 0.613 (0.054)
& 368.69 (806.30) \\

GG-KM (wavelet)
& Poisson
& 48/49
& 0.1488 (0.0137)
& 0.6176 (0.0337)
& 0.572 (0.061)
& 64.57 (146.75) \\

GG-KM (wavelet)
& Negative binomial
& 50/50
& 0.1482 (0.0131)
& 0.6193 (0.0320)
& 0.573 (0.060)
& 77.38 (219.31) \\

\hline
\end{tabular}
\end{table*}
%
% \begin{table*}[tbp]
\caption{
Best observed mean predictive performance in each simulation scenario
among Poisson and negative-binomial models.
For IPCW IBS and cure RMSE, smaller values are better, whereas for cure AUC,
larger values are better. Parenthesized values are Monte Carlo standard errors.
The displayed optimum is a descriptive comparison and is not a claim of
superiority when paired uncertainty includes zero.
GG, Spl, P, and NB denote generalized-gamma kernel machine, spline,
Poisson, and negative binomial, respectively.
For GG-KM fits, the selected kernel is shown in parentheses.
}
\label{tab:simulation-best}
\centering
\small
\setlength{\tabcolsep}{4pt}
\begin{tabular}{lllccc}
\hline
True surface
& Generating law
& $n$
& Lowest mean IPCW IBS
& Lowest mean cure RMSE
& Highest mean cure AUC \\
\hline

Linear
& P
& 300
& GG-NB (linear): 0.1929 (0.0013)
& GG-P (linear): 0.0630 (0.0019)
& GG-NB (poly): 0.6845 (0.0025) \\

Linear
& NB(0.25)
& 300
& GG-NB (linear): 0.1965 (0.0015)
& GG-NB (linear): 0.0602 (0.0017)
& GG-NB (poly): 0.6930 (0.0025) \\

Linear
& NB(1.00)
& 300
& GG-NB (linear): 0.1975 (0.0014)
& GG-P (linear): 0.0606 (0.0017)
& GG-NB (poly): 0.7032 (0.0025) \\

Linear
& P
& 600
& GG-NB (linear): 0.1937 (0.0010)
& GG-P (linear): 0.0458 (0.0011)
& GG-NB (poly): 0.6875 (0.0020) \\

Linear
& NB(0.25)
& 600
& GG-NB (linear): 0.1944 (0.0010)
& GG-P (linear): 0.0448 (0.0011)
& GG-NB (poly): 0.6910 (0.0019) \\

Linear
& NB(1.00)
& 600
& GG-NB (linear): 0.1968 (0.0011)
& GG-NB (linear): 0.0453 (0.0012)
& GG-NB (poly): 0.7045 (0.0018) \\

\hline

Quadratic
& P
& 300
& GG-NB (poly): 0.2097 (0.0014)
& GG-NB (poly): 0.0990 (0.0018)
& GG-NB (poly): 0.6411 (0.0029) \\

Quadratic
& NB(0.25)
& 300
& GG-P (rbf): 0.2128 (0.0014)
& GG-NB (poly): 0.0962 (0.0016)
& GG-P (poly): 0.6443 (0.0031) \\

Quadratic
& NB(1.00)
& 300
& GG-P (poly): 0.2145 (0.0016)
& GG-NB (poly): 0.0933 (0.0018)
& GG-P (poly): 0.6586 (0.0031) \\

Quadratic
& P
& 600
& GG-NB (poly): 0.2079 (0.0010)
& GG-NB (poly): 0.0720 (0.0013)
& GG-P (wavelet): 0.6422 (0.0018) \\

Quadratic
& NB(0.25)
& 600
& GG-NB (poly): 0.2096 (0.0009)
& GG-NB (poly): 0.0687 (0.0013)
& GG-P (wavelet): 0.6452 (0.0018) \\

Quadratic
& NB(1.00)
& 600
& GG-NB (poly): 0.2128 (0.0010)
& GG-NB (poly): 0.0646 (0.0012)
& GG-NB (poly): 0.6555 (0.0019) \\

\hline

Periodic
& P
& 300
& Spl-NB: 0.1913 (0.0017)
& Spl-NB: 0.1091 (0.0017)
& GG-NB (poly): 0.7237 (0.0034) \\

Periodic
& NB(0.25)
& 300
& Spl-NB: 0.1906 (0.0015)
& Spl-NB: 0.1050 (0.0017)
& GG-NB (poly): 0.7216 (0.0033) \\

Periodic
& NB(1.00)
& 300
& Spl-NB: 0.1905 (0.0016)
& Spl-NB: 0.1005 (0.0016)
& GG-NB (poly): 0.7345 (0.0033) \\

Periodic
& P
& 600
& Spl-NB: 0.1851 (0.0010)
& Spl-P: 0.0794 (0.0012)
& GG-NB (poly): 0.7301 (0.0023) \\

Periodic
& NB(0.25)
& 600
& Spl-NB: 0.1858 (0.0011)
& Spl-NB: 0.0778 (0.0012)
& GG-NB (poly): 0.7361 (0.0022) \\

Periodic
& NB(1.00)
& 600
& Spl-P: 0.1878 (0.0012)
& Spl-P: 0.0728 (0.0013)
& GG-NB (poly): 0.7501 (0.0020) \\

\hline
\end{tabular}
\end{table*}
%
% % Generated from the Poisson/NB regret results,
% restricted to spline models and GG-KM with the RBF kernel.
\begin{table*}[tbp]
\caption{
Paired regret and descriptive rank summaries over the 18 primary simulation
scenarios for the Poisson and negative-binomial models.
The comparison is restricted to spline models and GG-KM models using the RBF
kernel. Within each scenario and metric, the reference is the retained method
with the best finite mean. For IPCW integrated Brier score and cure RMSE,
regret is the method value minus the reference value; for cure AUC, regret is
the reference value minus the method value, so that smaller regret is better
for all metrics. Regret uses only replication pairs for which both methods are
finite. The worst regret identifies its scenario and reports a 95\% Monte Carlo
half-width; mean regret weights scenarios equally and propagates their Monte
Carlo standard errors. Mean rank is based on scenario means and is descriptive.
}
\label{tab:simulation-regret}
\centering
\resizebox{\textwidth}{!}{%
\begin{tabular}{llccc}
\hline
Metric
& Method
& Worst paired regret ($\pm$95\%; scenario)
& Mean paired regret ($\pm$95\%)
& Mean rank \\
\hline

IPCW integrated Brier score
& Spline Negative binomial
& 0.00067 $\pm$ 0.00056 (Quadratic, P, $n=600$)
& 0.00012 $\pm$ 0.00010
& 2.06 \\

IPCW integrated Brier score
& Spline Poisson
& 0.00065 $\pm$ 0.00056 (Quadratic, P, $n=600$)
& 0.00013 $\pm$ 0.00010
& 2.06 \\

IPCW integrated Brier score
& GG-KM Poisson (RBF)
& 0.00348 $\pm$ 0.00119 (Periodic, NB(1.00), $n=300$)
& 0.00108 $\pm$ 0.00017
& 3.22 \\

IPCW integrated Brier score
& GG-KM Negative binomial (RBF)
& 0.00359 $\pm$ 0.00120 (Periodic, NB(1.00), $n=300$)
& 0.00108 $\pm$ 0.00017
& 2.67 \\

\hline

Cure RMSE
& Spline Negative binomial
& 0.00464 $\pm$ 0.00210 (Quadratic, NB(1.00), $n=600$)
& 0.00073 $\pm$ 0.00025
& 1.94 \\

Cure RMSE
& Spline Poisson
& 0.00464 $\pm$ 0.00210 (Quadratic, NB(1.00), $n=600$)
& 0.00076 $\pm$ 0.00025
& 1.94 \\

Cure RMSE
& GG-KM Poisson (RBF)
& 0.01651 $\pm$ 0.00268 (Periodic, NB(1.00), $n=300$)
& 0.00631 $\pm$ 0.00054
& 3.17 \\

Cure RMSE
& GG-KM Negative binomial (RBF)
& 0.01627 $\pm$ 0.00272 (Periodic, NB(1.00), $n=300$)
& 0.00632 $\pm$ 0.00053
& 2.94 \\

\hline

Cure AUC
& Spline Negative binomial
& 0.00768 $\pm$ 0.00529 (Quadratic, P, $n=300$)
& 0.00208 $\pm$ 0.00069
& 2.61 \\

Cure AUC
& Spline Poisson
& 0.00693 $\pm$ 0.00538 (Quadratic, P, $n=300$)
& 0.00224 $\pm$ 0.00070
& 2.94 \\

Cure AUC
& GG-KM Poisson (RBF)
& 0.00795 $\pm$ 0.00464 (Periodic, P, $n=300$)
& 0.00157 $\pm$ 0.00061
& 2.39 \\

Cure AUC
& GG-KM Negative binomial (RBF)
& 0.00687 $\pm$ 0.00464 (Periodic, P, $n=300$)
& 0.00144 $\pm$ 0.00060
& 2.06 \\

\hline
\end{tabular}%
}
\end{table*}
%
% % Generated from highdim_best_poisson_nb_parametric_updated.csv
\begin{table*}[tbp]
\caption{
Best observed mean predictive performance in each high-dimensional simulation
scenario. For IPCW IBS and cure RMSE, smaller values are better, whereas for
cure AUC, larger values are better. Parenthesized values are Monte Carlo
standard errors. The comparison includes Poisson and negative-binomial GG-KM
models, Poisson and negative-binomial spline models, the mixture-cure model,
and the promotion-time model. Int and Rad denote Interaction and Radial
surfaces, respectively. GG, Spl, P, and NB denote generalized-gamma kernel
machine, spline, Poisson, and negative binomial, respectively. For GG-KM fits,
the selected kernel is shown in parentheses.
}
\label{tab:highdim-best}
\centering
\resizebox{\textwidth}{!}{%
\begin{tabular}{lllccc}
\hline
True surface
& Generating law
& $n$
& Lowest mean IPCW IBS
& Lowest mean cure RMSE
& Highest mean cure AUC \\
\hline

Interaction
& P
& 300
& GG-NB (polynomial): 0.2045 (0.0017)
& GG-NB (polynomial): 0.1250 (0.0017)
& GG-NB (polynomial): 0.7052 (0.0024) \\

Interaction
& NB(0.25)
& 300
& GG-NB (polynomial): 0.2075 (0.0015)
& GG-NB (polynomial): 0.1254 (0.0017)
& GG-NB (polynomial): 0.7179 (0.0024) \\

Interaction
& NB(1.00)
& 300
& GG-NB (polynomial): 0.2093 (0.0017)
& GG-NB (polynomial): 0.1179 (0.0018)
& GG-NB (polynomial): 0.7289 (0.0023) \\

Interaction
& P
& 600
& GG-NB (polynomial): 0.2018 (0.0011)
& GG-NB (polynomial): 0.0961 (0.0014)
& GG-NB (polynomial): 0.7043 (0.0017) \\

Interaction
& NB(0.25)
& 600
& GG-NB (polynomial): 0.2027 (0.0011)
& GG-NB (polynomial): 0.0912 (0.0013)
& GG-NB (polynomial): 0.7091 (0.0017) \\

Interaction
& NB(1.00)
& 600
& GG-NB (polynomial): 0.2027 (0.0011)
& GG-NB (polynomial): 0.0869 (0.0012)
& GG-P (polynomial): 0.7250 (0.0017) \\

\hline

Radial
& P
& 300
& Spl-P: 0.2138 (0.0014)
& Spl-P: 0.1168 (0.0017)
& GG-NB (polynomial): 0.6521 (0.0023) \\

Radial
& NB(0.25)
& 300
& Spl-P: 0.2167 (0.0014)
& Spl-P: 0.1150 (0.0016)
& GG-NB (polynomial): 0.6602 (0.0023) \\

Radial
& NB(1.00)
& 300
& Spl-NB: 0.2226 (0.0015)
& Spl-NB: 0.1117 (0.0015)
& GG-NB (polynomial): 0.6633 (0.0024) \\

Radial
& P
& 600
& Spl-P: 0.2138 (0.0009)
& Spl-NB: 0.0926 (0.0011)
& GG-NB (polynomial): 0.6426 (0.0019) \\

Radial
& NB(0.25)
& 600
& Spl-P: 0.2136 (0.0009)
& Spl-NB: 0.0897 (0.0010)
& GG-NB (polynomial): 0.6480 (0.0018) \\

Radial
& NB(1.00)
& 600
& Spl-P: 0.2174 (0.0010)
& Spl-P: 0.0881 (0.0010)
& GG-NB (polynomial): 0.6492 (0.0018) \\

\hline
\end{tabular}%
}
\end{table*}
%
% \clearpage
% % USG records the post-\clearpage counter as the last page; restore the last physical page.
% \addtocounter{page}{-1}
% \fi

% \end{document}
